%==========================================================
%  Contemporary Economic Issues: A Canadian Perspective
%  ECON 204 / INTS 211 — University of Northern British Columbia
%  Muhebullah Karimzada — Winter 2027
%==========================================================

\documentclass[12pt,twoside,openright]{book}

%----------------------------------------------------------
% ENCODING & FONTS
%----------------------------------------------------------
\usepackage[T1]{fontenc}
\usepackage[utf8]{inputenc}
\usepackage{lmodern}
\usepackage{microtype}
\usepackage{palatino}

%----------------------------------------------------------
% PAGE GEOMETRY
%----------------------------------------------------------
\usepackage[
  top=1in,
  bottom=1.1in,
  inner=1.25in,
  outer=0.875in,
  headheight=15pt
]{geometry}

%----------------------------------------------------------
% MATHEMATICS
%----------------------------------------------------------
\usepackage{amsmath,amssymb,amsthm}

%----------------------------------------------------------
% COLORS  (UNBC green: RGB 0,122,51)
%----------------------------------------------------------
\usepackage[dvipsnames]{xcolor}
\definecolor{unbcgreen}{RGB}{0,122,51}
\definecolor{unbcgreenlight}{RGB}{228,244,234}
\definecolor{unbcgreenmid}{RGB}{178,222,194}
\definecolor{policyblue}{RGB}{0,60,120}
\definecolor{policyblulight}{RGB}{224,234,248}
\definecolor{exampleteal}{RGB}{0,118,115}
\definecolor{exampleteallight}{RGB}{222,244,242}
\definecolor{thinkpurple}{RGB}{88,20,138}
\definecolor{thinkpurplelight}{RGB}{238,226,252}
\definecolor{dataorange}{RGB}{172,84,0}
\definecolor{dataorangelight}{RGB}{255,237,216}
\definecolor{chaptergray}{RGB}{75,75,75}
\definecolor{lightgray}{RGB}{245,245,245}
\definecolor{casered}{RGB}{120,20,40}
\definecolor{caseredlight}{RGB}{248,224,228}
\definecolor{everygold}{RGB}{148,96,0}
\definecolor{everygoldlight}{RGB}{255,248,220}

%----------------------------------------------------------
% GRAPHICS & TIKZ
%----------------------------------------------------------
\usepackage{graphicx}
\usepackage{tikz}
\usepackage{pgfplots}
\pgfplotsset{compat=1.15}
\usetikzlibrary{
  arrows.meta,
  positioning,
  shapes.geometric,
  shapes.misc,
  calc,
  patterns,
  decorations.pathreplacing,
  decorations.markings,
  fit,
  backgrounds,
  shadows
}

%----------------------------------------------------------
% TABLES
%----------------------------------------------------------
\usepackage{booktabs}
\usepackage{tabularx}
\usepackage{multirow}
\usepackage{array}
\usepackage{longtable}
\renewcommand{\arraystretch}{1.35}

%----------------------------------------------------------
% LISTS
%----------------------------------------------------------
\usepackage[shortlabels]{enumitem}
\setlist[itemize]{leftmargin=*,topsep=4pt,itemsep=2pt}
\setlist[enumerate]{leftmargin=*,topsep=4pt,itemsep=2pt}

%----------------------------------------------------------
% HEADERS & FOOTERS
%----------------------------------------------------------
\usepackage{fancyhdr}
\pagestyle{fancy}
\fancyhf{}
\fancyhead[LE]{\small\thepage\quad\textcolor{unbcgreen}{\textit{\leftmark}}}
\fancyhead[RO]{\textcolor{unbcgreen}{\textit{\rightmark}}\quad\small\thepage}
\fancyfoot{}
\renewcommand{\headrulewidth}{0.5pt}
\renewcommand{\headrule}{\hbox to\headwidth{\color{unbcgreen}\leaders\hrule height \headrulewidth\hfill}}

%----------------------------------------------------------
% CHAPTER & SECTION FORMATTING
%----------------------------------------------------------
\usepackage{titlesec}

\titleformat{\chapter}[display]
  {\normalfont\bfseries\color{unbcgreen}}
  {\large\MakeUppercase{\chaptertitlename}\ \thechapter}
  {8pt}
  {\Huge}
  [\vspace{6pt}\color{unbcgreen}\titlerule]

\titlespacing*{\chapter}{0pt}{-20pt}{32pt}

\titleformat{\section}
  {\normalfont\Large\bfseries\color{unbcgreen!85!black}}
  {\thesection}{1em}{}

\titleformat{\subsection}
  {\normalfont\large\bfseries\color{chaptergray}}
  {\thesubsection}{1em}{}

\titleformat{\subsubsection}
  {\normalfont\normalsize\bfseries\itshape\color{chaptergray}}
  {\thesubsubsection}{1em}{}

%----------------------------------------------------------
% CAPTIONS
%----------------------------------------------------------
\usepackage{caption}
\captionsetup{
  font=small,
  labelfont={bf,color=unbcgreen},
  labelsep=period,
  margin=12pt,
  justification=justified
}

%----------------------------------------------------------
% FLOATS
%----------------------------------------------------------
\usepackage{float}
\usepackage[section]{placeins}

%----------------------------------------------------------
% FOOTNOTES
%----------------------------------------------------------
\usepackage{footmisc}

%----------------------------------------------------------
% SPACING
%----------------------------------------------------------
\usepackage{setspace}
\onehalfspacing

%----------------------------------------------------------
% HYPERLINKS
%----------------------------------------------------------
\usepackage[
  colorlinks=true,
  linkcolor=unbcgreen,
  citecolor=policyblue,
  urlcolor=policyblue,
  pdftitle={Contemporary Economic Issues: A Canadian Perspective},
  pdfauthor={Muhebullah Karimzada},
  pdfsubject={ECON 204 / INTS 211, UNBC, Winter 2027},
  bookmarks=true,
  bookmarksopen=true
]{hyperref}

%----------------------------------------------------------
% TCOLORBOX ENVIRONMENTS
%----------------------------------------------------------
\usepackage[most]{tcolorbox}

%--- Learning Objectives (green, chapter opener) ----------
\newtcolorbox{learningobjectives}{
  enhanced,
  colback=unbcgreenlight,
  colframe=unbcgreen,
  colbacktitle=unbcgreen,
  coltitle=white,
  fonttitle=\bfseries\large,
  title={Learning Objectives},
  arc=4pt,
  boxrule=1pt,
  left=10pt, right=10pt, top=8pt, bottom=8pt,
  attach boxed title to top left={yshift=-2mm,xshift=8mm},
  boxed title style={arc=3pt},
  before upper={\begin{itemize}[leftmargin=1.2em,topsep=2pt,itemsep=3pt,parsep=0pt]},
  after upper={\end{itemize}}
}

%--- Definition Box (gray) --------------------------------
\newtcolorbox{defbox}[1]{
  enhanced,
  colback=lightgray,
  colframe=gray!55,
  colbacktitle=gray!30,
  coltitle=black,
  fonttitle=\bfseries,
  title={Definition: #1},
  arc=3pt,
  boxrule=0.7pt,
  left=8pt, right=8pt, top=6pt, bottom=6pt
}

%--- Canadian Example Box (teal) -------------------------
\newtcolorbox{canexample}[1]{
  enhanced,
  colback=exampleteallight,
  colframe=exampleteal,
  colbacktitle=exampleteal,
  coltitle=white,
  fonttitle=\bfseries,
  title={Canadian Example: #1},
  arc=3pt,
  boxrule=1pt,
  left=8pt, right=8pt, top=6pt, bottom=6pt
}

%--- Think Like an Economist (purple) --------------------
\newtcolorbox{thinkbox}{
  enhanced,
  colback=thinkpurplelight,
  colframe=thinkpurple,
  colbacktitle=thinkpurple,
  coltitle=white,
  fonttitle=\bfseries,
  title={Think Like an Economist},
  arc=3pt,
  boxrule=1pt,
  left=8pt, right=8pt, top=6pt, bottom=6pt
}

%--- Policy Debate (blue) --------------------------------
\newtcolorbox{policydebate}[1]{
  enhanced,
  breakable,
  colback=policyblulight,
  colframe=policyblue,
  colbacktitle=policyblue,
  coltitle=white,
  fonttitle=\bfseries\large,
  title={Policy Debate: #1},
  arc=4pt,
  boxrule=1pt,
  left=10pt, right=10pt, top=8pt, bottom=8pt
}

%--- Case Study (dark red) --------------------------------
\newtcolorbox{casestudy}[1]{
  enhanced,
  breakable,
  colback=caseredlight,
  colframe=casered,
  colbacktitle=casered,
  coltitle=white,
  fonttitle=\bfseries\large,
  title={Case Study: #1},
  arc=4pt,
  boxrule=1pt,
  left=10pt, right=10pt, top=8pt, bottom=8pt
}

%--- Economics in Everyday Life (gold) -------------------
\newtcolorbox{everydaylife}[1]{
  enhanced,
  breakable,
  colback=everygoldlight,
  colframe=everygold,
  colbacktitle=everygold,
  coltitle=white,
  fonttitle=\bfseries\large,
  title={Economics in Everyday Life: #1},
  arc=4pt,
  boxrule=1pt,
  left=10pt, right=10pt, top=8pt, bottom=8pt
}

%--- Key Takeaway (green) --------------------------------
\newtcolorbox{keytakeaway}{
  enhanced,
  colback=unbcgreenlight,
  colframe=unbcgreen,
  colbacktitle=unbcgreen,
  coltitle=white,
  fonttitle=\bfseries,
  title={Key Takeaway},
  arc=3pt,
  boxrule=0.8pt,
  left=8pt, right=8pt, top=6pt, bottom=6pt
}

%--- Data Spotlight (orange) -----------------------------
\newtcolorbox{databox}[1]{
  enhanced,
  colback=dataorangelight,
  colframe=dataorange,
  colbacktitle=dataorange,
  coltitle=white,
  fonttitle=\bfseries,
  title={Data Spotlight: #1},
  arc=3pt,
  boxrule=0.8pt,
  left=8pt, right=8pt, top=6pt, bottom=6pt
}

%----------------------------------------------------------
% THEOREM-LIKE ENVIRONMENTS
%----------------------------------------------------------
\theoremstyle{definition}
\newtheorem{proposition}{Proposition}[chapter]
\newtheorem{example}{Example}[chapter]

%----------------------------------------------------------
% CUSTOM COMMANDS
%----------------------------------------------------------
\newcommand{\term}[1]{\textbf{\textit{#1}}}
\newcommand{\source}[1]{\par\vspace{4pt}\noindent\footnotesize\textit{Source: #1}}
\newcommand{\figexplain}[1]{\par\vspace{4pt}\noindent\footnotesize #1}
\newcommand{\dollar}{\$}
\newcommand{\percent}{\%}

%----------------------------------------------------------
% BIBLIOGRAPHY
%----------------------------------------------------------
\usepackage[authoryear,round,sort]{natbib}
\bibliographystyle{apalike}

%==========================================================
\begin{document}
%==========================================================

%----------------------------------------------------------
% FRONT MATTER
%----------------------------------------------------------
\frontmatter
\pagestyle{empty}

%--- TITLE PAGE -------------------------------------------
\begin{titlepage}
\begin{tikzpicture}[remember picture,overlay]
  % Top green bar
  \fill[unbcgreen] (current page.north west) rectangle
    ([yshift=-3.5cm]current page.north east);
  % Bottom green bar
  \fill[unbcgreen!30] (current page.south west) rectangle
    ([yshift=1.5cm]current page.south east);
\end{tikzpicture}

\vspace*{2.2cm}

\begin{center}
  {\color{white}\fontsize{28}{34}\selectfont\bfseries
    Contemporary Economic Issues:\\[6pt]
    A Canadian Perspective}

  \vspace{1cm}
  \color{unbcgreen}
  \rule{0.85\textwidth}{1.5pt}

  \vspace{1cm}
  {\Large\bfseries\color{black} Muhebullah Karimzada}\\[6pt]
  {\large\color{chaptergray} School of Economics}\\[4pt]
  {\large\color{chaptergray} University of Northern British Columbia}

  \vspace{1.5cm}
  {\normalsize\color{chaptergray}
    Prepared for \textbf{ECON 204 / INTS 211}\\
    \textit{Contemporary Economic Issues}\\[4pt]
    Winter 2027}

  \vfill
  {\small\color{chaptergray}
    University of Northern British Columbia\\
    Prince George, British Columbia, Canada}
\end{center}
\end{titlepage}

%--- COPYRIGHT PAGE ---------------------------------------
\newpage
\thispagestyle{empty}
\vspace*{\fill}
\begin{small}
\noindent\textcopyright\ 2027 Muhebullah Karimzada. University of Northern British Columbia.

\vspace{8pt}
\noindent This textbook is prepared for use in ECON 204 / INTS 211 at UNBC. All rights reserved.
No part of this publication may be reproduced or distributed without the written permission
of the author.

\vspace{8pt}
\noindent\textbf{Suggested Citation:}\\
Karimzada, M. (2027). \textit{Contemporary Economic Issues: A Canadian Perspective}.
University of Northern British Columbia.

\vspace{8pt}
\noindent Data sources cited throughout this textbook include Statistics Canada, the Bank of Canada,
the Organisation for Economic Co-operation and Development (OECD), the International Monetary
Fund (IMF), and the World Bank. All data are for illustrative and educational purposes.

\vspace{8pt}
\noindent Printed and distributed at Prince George, British Columbia, Canada.
\end{small}
\vspace*{\fill}

%--- PREFACE ----------------------------------------------
\chapter*{Preface}
\addcontentsline{toc}{chapter}{Preface}
\markboth{Preface}{Preface}

\noindent This textbook was written to fill a gap that became apparent in teaching ECON 204 at the University
of Northern British Columbia: most available economics textbooks treat contemporary issues through
an American or generic international lens, leaving Canadian students without a primary reference
that speaks directly to their economy, their institutions, and their policy debates.

The Canadian economy has its own character. It is deeply integrated with the United States yet
markedly different in its institutional architecture---universal healthcare, a federal carbon price,
a publicly owned central bank with an explicit inflation target, and a tradition of provincially
diverse labour markets. Housing affordability in Vancouver and Toronto is not the same problem
as in Houston. Canada's experience with post-COVID inflation, while similar in some dimensions
to the American experience, has its own drivers, timeline, and policy response. Understanding
these differences requires a textbook that takes Canada seriously as its own economic case study.

\vspace{10pt}
\noindent\textbf{Who This Book Is For}

This textbook is written for second-year undergraduate students in economics or related social
science disciplines who have completed introductory microeconomics and macroeconomics. It
assumes you know what supply and demand curves are, understand the basics of GDP, and have
been introduced to the concepts of inflation and unemployment. It does not assume you have
studied advanced mathematics or econometrics.

The goal is to help you move from understanding basic economic concepts to using them as
analytical tools for reasoning about real, complex, messy issues---the kind you read about in
the news every day.

\vspace{10pt}
\noindent\textbf{What Makes This Textbook Different}

Three features distinguish this textbook:

First, it is \textbf{Canadian}. More than 80\% of the examples, case studies, and data refer to
Canada. When international comparisons are useful, they are made explicitly and for a reason.

Second, it is \textbf{contemporary}. Each chapter addresses an issue that is actively shaping
Canada's economic future: the inflation surge of 2022--2023, the housing crisis, artificial
intelligence and the future of work, and climate policy. Economic theory is the tool;
contemporary Canada is the subject.

Third, it is \textbf{honest about uncertainty}. Economics does not always deliver clean answers.
Every chapter includes a ``Policy Debate'' section that presents multiple perspectives, weighs
the evidence, and acknowledges what we do not yet know. This is closer to how economics is
actually practised than many textbooks suggest.

\vspace{10pt}
\noindent\textbf{Acknowledgements}

I am grateful to students in ECON 204 at UNBC whose questions and discussions have shaped
the structure and tone of every chapter. Northern British Columbia, with its resource economy,
Indigenous communities, and distinctive labour markets, provides a uniquely rich context
for studying contemporary economic issues, and I have tried to honour that richness throughout.

\vspace{10pt}
\hfill\textit{Muhebullah Karimzada}\\
\hfill Prince George, British Columbia\\
\hfill January 2027

%--- HOW TO USE THIS TEXTBOOK ----------------------------
\chapter*{How to Use This Textbook}
\addcontentsline{toc}{chapter}{How to Use This Textbook}
\markboth{How to Use This Textbook}{How to Use This Textbook}

\noindent Each chapter in this textbook follows a consistent structure designed to help you build
understanding progressively. Here is a guide to the features you will encounter.

\vspace{10pt}
\noindent\textbf{Chapter Features}

\begin{description}[leftmargin=2cm,style=nextline]

\item[\textcolor{unbcgreen}{\bfseries Learning Objectives}]
Every chapter opens with a list of what you should be able to do after studying the chapter.
Return to this list when reviewing for exams.

\item[\textcolor{unbcgreen}{\bfseries Definition Boxes}]
Key terms are defined in gray boxes the first time they appear. These terms form the
foundation of economic reasoning and will appear in exam questions.

\item[\textcolor{exampleteal}{\bfseries Canadian Example Boxes}]
Teal boxes highlight real Canadian applications of the theory being discussed. These help
connect abstract concepts to recognizable events and institutions.

\item[\textcolor{thinkpurple}{\bfseries Think Like an Economist}]
Purple boxes invite you to apply economic reasoning to a question or scenario. These are
excellent preparation for analytical exam questions.

\item[\textcolor{policyblue}{\bfseries Policy Debate Boxes}]
Blue boxes present competing perspectives on a policy question. They show that economics
involves genuine disagreement and trade-offs, not just formulas.

\item[\textcolor{unbcgreen}{\bfseries Key Takeaway}]
Green summary boxes at strategic points in the chapter capture the most important insight
from the preceding section.

\item[\textcolor{dataorange}{\bfseries Data Spotlight}]
Orange boxes present important Canadian economic data, explain how to read it, and
discuss what it tells us.

\end{description}

\vspace{10pt}
\noindent\textbf{End-of-Chapter Material}

Each chapter ends with:
\begin{itemize}
  \item \textbf{Chapter Summary} --- a concise review of the main ideas.
  \item \textbf{Review Questions} --- concept and comprehension questions aligned with the
        learning objectives.
  \item \textbf{Discussion Questions} --- open-ended questions for seminar discussion or
        short written responses.
  \item \textbf{Exercises} --- quantitative, graphical, and applied problems similar to
        those on midterm and final exams.
\end{itemize}

\vspace{10pt}
\noindent\textbf{Figures}

Every figure includes a descriptive title, labeled axes, and an explanatory caption. Before
looking at the caption, try to interpret the figure on your own. Reading graphs is a skill
that develops with practice and is tested in exams.

\vspace{10pt}
\noindent\textbf{A Note on Data}

Economic data change over time. The figures and statistics in this textbook reflect data
available at the time of writing (late 2026). For the most current data, consult
Statistics Canada (\url{www.statcan.gc.ca}) and the Bank of Canada
(\url{www.bankofcanada.ca}).

%--- TABLE OF CONTENTS ------------------------------------
\tableofcontents
\listoffigures
\listoftables

%----------------------------------------------------------
% MAIN MATTER
%----------------------------------------------------------
\mainmatter
\pagestyle{fancy}

%==========================================================
%  CHAPTER 1: The Canadian Economy — Structure, Institutions,
%              and Analytical Tools
%  ECON 204 / Contemporary Economic Issues: A Canadian Perspective
%  Muhebullah Karimzada — University of Northern British Columbia
%==========================================================

\chapter{The Canadian Economy: Structure, Institutions, and Analytical Tools}
\label{ch:introduction}

%----------------------------------------------------------
%  OPENING VIGNETTE
%----------------------------------------------------------

\noindent In the spring of 2024, four people in Prince George, British Columbia
faced economic decisions shaped by forces none of them fully understood.

Diane, 52, received her mortgage renewal notice in March. Her monthly payment
was rising from \$1,840 to \$2,710 --- an increase of \$870, or more than
\$10,000 per year. Nothing had changed in her household except one number:
the Bank of Canada's overnight interest rate, which had climbed from 0.25\%
in early 2022 to 5.00\% by mid-2023. She had not overspent. She had not lost
her job. But a decision made in an Ottawa boardroom had fundamentally altered
her monthly budget.

Her son Marcus, 24, had just graduated from UNBC with a degree in business
administration. He was earning \$22 per hour at a local firm --- a wage that
would have seemed generous five years ago, but which, after two years of
inflation averaging over 5\%, felt significantly less secure. He was
trying to save for a down payment on an apartment, but rental costs in Prince
George had risen nearly 30\% since 2020. His parents' house had doubled in
value, but that wealth would not be his for decades.

At the same sawmill where Marcus's uncle worked, the plant manager was
navigating a \$3 million investment decision: automate the log-grading line
with computer vision technology, or keep 18 workers. The system would pay for
itself in four years. The workers earned above the minimum wage, had families
in the community, and spent their incomes at local businesses. The economics
pointed one direction; the community consequences pointed another.

And at the Northern Health Authority, administrator Priya was reviewing a
budget that had grown 14\% in two years --- driven by an aging patient
population, nurse recruitment costs, and pharmaceutical price increases ---
while the provincial government's transfer payment had grown by only 6\%.
The gap had to come from somewhere.

These four people live in the same economy. They are affected by the same
interest rates, the same inflation, the same labour market pressures, and the
same fiscal constraints. None of their situations can be understood without
understanding how the Canadian economy works, how its institutions make
decisions, and how those decisions ripple outward into ordinary lives.

That is what this textbook is about.

%----------------------------------------------------------
%  LEARNING OBJECTIVES
%----------------------------------------------------------

\begin{learningobjectives}
  \item Describe the scale and structure of the Canadian economy using key
        aggregate statistics, including GDP, GDP per capita, and sectoral
        composition.
  \item Apply the expenditure approach to GDP and calculate nominal GDP growth,
        real GDP growth, and GDP per capita from raw data.
  \item Identify Canada's major economic institutions and explain the mandate,
        tools, and independence of each.
  \item Distinguish between monetary policy and fiscal policy, explain the
        transmission channels of each, and identify which institution is
        responsible for each in Canada.
  \item Apply the core tools of economic reasoning --- models, opportunity cost,
        marginal analysis, and the positive--normative distinction --- to
        real-world Canadian policy questions.
  \item Situate Canada within an international context by comparing key
        macroeconomic indicators with OECD peers.
  \item Preview the ten major contemporary economic issues addressed in this
        textbook and explain how they are interconnected.
\end{learningobjectives}

%==========================================================
\section{Why Study Contemporary Economic Issues?}
%==========================================================

Economics is often introduced as the study of scarcity: because resources are
limited and human wants are not, societies must make choices about how to
allocate what they have. That definition is correct, but it understates the
subject's reach. A more precise framing for this course is that economics is
the study of \term{incentives}, \term{trade-offs}, and \term{unintended
consequences} --- the three forces that shape virtually every economic outcome
a society produces.

\term{Contemporary economic issues} are economic problems that are actively
unfolding, contested among reasonable people, and consequential for the
well-being of Canadians. They differ from purely theoretical questions
because they involve real institutions making real decisions under genuine
uncertainty. They differ from purely political questions because the
empirical evidence matters --- how markets actually behave, how people
actually respond to incentives, and what policies have actually worked
in comparable contexts.

Consider the four people from our opening vignette. Diane's mortgage crisis
cannot be understood without knowing why the Bank of Canada raised interest
rates, what caused inflation to spike in 2022, and how monetary policy
transmits through the financial system to household balance sheets. Marcus's
wage frustration cannot be understood without knowing how labour markets
determine compensation, why real wages have lagged productivity growth since
the 1980s, and how automation changes the demand for different kinds of work.
The sawmill manager's dilemma cannot be understood without the economics of
technological adoption, capital investment, and labour market adjustment. And
Priya's budget cannot be balanced without understanding fiscal federalism,
demographic aging, and the political economy of healthcare spending.

Economics provides the analytical tools to work through each of these problems
systematically. It does not always deliver clean, comfortable answers. But it
imposes a discipline on reasoning that is absent from purely ideological or
intuitive approaches: it forces you to trace consequences, identify trade-offs,
acknowledge trade-offs on the other side, and ask whether the evidence
supports your conclusions.

\begin{thinkbox}
When you encounter a policy proposal in this course --- a carbon tax, a rent
control ordinance, an immigration target, a minimum wage increase --- train
yourself to ask three questions before forming an opinion:
\begin{enumerate}[(1)]
  \item \textbf{What incentives does this policy create?} Who responds, how,
        and in what direction?
  \item \textbf{What are the trade-offs?} What must be given up to achieve
        this goal?
  \item \textbf{What might happen that the policymaker did not intend?}
        Economics is full of well-intentioned policies that produced
        surprising consequences.
\end{enumerate}
These questions do not determine whether a policy is good or bad. They are
the beginning of a serious analysis, not its conclusion.
\end{thinkbox}

%==========================================================
\section{The Canadian Economy at a Glance}
%==========================================================

Before analysing specific contemporary issues, we need a clear statistical
portrait of the economy we are studying. This section establishes the
foundational facts of Canadian economic structure that every subsequent
chapter will draw upon.

%----------------------------------------------------------
\subsection{Scale and Aggregate Output}
%----------------------------------------------------------

Canada is the world's ninth-largest economy measured by nominal gross domestic
product. In 2024, Canada's GDP reached approximately \$3.0 trillion in
Canadian dollars --- a figure that places Canada ahead of Australia, South
Korea, and Spain, but significantly behind the United States (\$29~trillion
USD) and far smaller than China and the European Union collectively. With a
population of approximately 41.5 million people, Canada's GDP per capita ---
our first approximation of average living standards --- was roughly
\$72,000~CAD, or approximately \$53,000~USD at prevailing exchange rates.

\begin{defbox}{Gross Domestic Product (GDP)}
\textbf{GDP} is the total market value of all final goods and services
produced within a country's borders during a given period, typically one
calendar year or one fiscal quarter.

\textit{Economic intuition:} GDP measures the size of the economy's
productive output. It captures what the economy \textit{produces} --- not
what it owns, not what it earns abroad, and not what flows from one
person to another without new production (such as transfer payments or
the sale of existing assets).

\textit{What GDP excludes:} Household production (cooking, childcare at
home), the underground economy, environmental degradation, the distribution
of income, and non-material well-being.

\textit{Canadian example:} In 2020, Canada's real GDP fell 5.2\% ---
the steepest single-year contraction since the Great Depression --- as
COVID-19 forced businesses to close and household spending collapsed.
By 2021, GDP rebounded by 5.0\% as restrictions lifted. The swing of
more than 10 percentage points in two years illustrates both the fragility
and the resilience of a modern market economy.
\end{defbox}

%----------------------------------------------------------
\subsection{The Expenditure Approach to GDP}
\label{sec:expenditure}
%----------------------------------------------------------

There are three equivalent ways to measure GDP: the expenditure approach,
the income approach, and the production approach. In practice, national
statistical agencies use all three and reconcile any discrepancies.
For macroeconomic analysis, the \term{expenditure approach} is the most
widely used because it decomposes GDP by who is spending.

The fundamental expenditure identity is:

\begin{equation}
  Y \;=\; C \;+\; I \;+\; G \;+\; (X - M)
  \label{eq:gdp_expenditure}
\end{equation}

\noindent where:
\begin{itemize}
  \item $Y$ is nominal GDP (total output, measured in current dollars)
  \item $C$ is \term{household consumption expenditure} --- spending by
        households on goods (food, clothing, vehicles) and services
        (healthcare, education, financial services). This is typically
        the largest component, accounting for roughly 55--60\% of GDP
        in Canada.
  \item $I$ is \term{gross investment} (or gross capital formation) ---
        business spending on plant, equipment, and inventory, plus
        residential construction. This is the most volatile component
        of GDP.
  \item $G$ is \term{government expenditure on goods and services} ---
        note this includes only direct government purchases, \textit{not}
        transfer payments such as Employment Insurance or the Canada
        Child Benefit, which are redistribution rather than new production.
  \item $X$ is exports --- goods and services produced in Canada and
        sold abroad.
  \item $M$ is imports --- goods and services produced abroad and
        purchased in Canada.
  \item $(X - M) = NX$ is \term{net exports} (the trade balance).
        A positive NX means Canada sells more abroad than it buys ---
        a trade surplus. A negative NX means a trade deficit.
\end{itemize}

\begin{databox}{Canada's GDP by Expenditure Component, 2024 (Approximate)}
\begin{center}
\renewcommand{\arraystretch}{1.4}
\begin{tabular}{lrr}
\toprule
\textbf{Component} & \textbf{Value (\$B~CAD)} & \textbf{Share of GDP (\%)} \\
\midrule
Household Consumption ($C$)     & 1{,}652 & 55.1 \\
Gross Investment ($I$)          &   652  & 21.7 \\
Government Expenditure ($G$)    &   718  & 23.9 \\
Net Exports ($NX = X - M$)      &  $-$21 & $-$0.7 \\
\midrule
\textbf{GDP} ($Y = C+I+G+NX$)   & \textbf{3{,}001} & \textbf{100.0} \\
\bottomrule
\end{tabular}
\end{center}
\vspace{2pt}
\textit{Note:} Values are approximate and illustrative, consistent with
Statistics Canada national accounts data. Figures may not sum to exactly
3{,}001 due to rounding.
\source{Statistics Canada, Table 36-10-0104-01 (National Accounts, expenditure
approach), 2024 estimates.}
\end{databox}

\noindent\textbf{Worked Example 1.1 --- Calculating GDP from Components}

Suppose you are given the following data for a simplified economy:

\begin{center}
\begin{tabular}{lr}
\toprule
Item & Value (\$B) \\
\midrule
Household spending on food, housing, and services & 1{,}200 \\
Business investment in machinery and factories   &   400  \\
Residential housing construction                  &   150  \\
Federal, provincial, and municipal government purchases & 500 \\
Exports of goods and services                    &   620  \\
Imports of goods and services                    &   670  \\
\bottomrule
\end{tabular}
\end{center}

\noindent\textit{Step 1:} Identify the components of Equation~\eqref{eq:gdp_expenditure}:
\begin{align*}
C &= 1{,}200 \\
I &= 400 + 150 = 550 \quad \text{(business investment + residential construction)} \\
G &= 500 \\
NX &= 620 - 670 = -50 \quad \text{(a trade deficit of \$50 billion)}
\end{align*}

\noindent\textit{Step 2:} Apply the identity:
\[
Y = 1{,}200 + 550 + 500 + (-50) = \mathbf{2{,}200 \text{ billion}}
\]

\noindent\textit{Interpretation:} This economy produces \$2.2~trillion of
output. The trade deficit (\$50B) means the economy is spending more than
it produces --- the excess is financed by borrowing from abroad or by
foreigners purchasing domestic assets.

\begin{keytakeaway}
GDP measures production, not welfare. A society can have high GDP and
high inequality, high GDP and poor health outcomes, or high GDP and
severe environmental degradation. GDP is an indispensable starting point
for macroeconomic analysis, but it is not the endpoint for evaluating
economic well-being.
\end{keytakeaway}

%----------------------------------------------------------
\subsection{Real GDP, Nominal GDP, and the GDP Deflator}
\label{sec:realgdp}
%----------------------------------------------------------

When we compare GDP across years, we face a measurement challenge: prices
change over time. If Canadian GDP rose from \$2.8~trillion in 2022 to
\$3.0~trillion in 2024, some of that increase reflects more goods and
services being produced, but some reflects higher prices for the same
goods and services. Only the first part represents genuine economic growth.

\term{Nominal GDP} measures output at the \textit{current year's prices}.
\term{Real GDP} measures output at the prices of a fixed \term{base year},
removing the effect of inflation.

The conversion between them uses the \term{GDP deflator}:

\begin{equation}
  \text{GDP Deflator}_t = \frac{\text{Nominal GDP}_t}{\text{Real GDP}_t} \times 100
  \label{eq:gdp_deflator}
\end{equation}

Rearranging, real GDP can be obtained from nominal GDP and the deflator:

\begin{equation}
  \text{Real GDP}_t = \frac{\text{Nominal GDP}_t}{\text{GDP Deflator}_t / 100}
  \label{eq:real_gdp}
\end{equation}

The \term{real GDP growth rate} --- the standard measure of economic
expansion or contraction --- is:

\begin{equation}
  g_t = \frac{\text{Real GDP}_t - \text{Real GDP}_{t-1}}{\text{Real GDP}_{t-1}} \times 100
  \label{eq:gdp_growth}
\end{equation}

\noindent\textbf{Worked Example 1.2 --- Real vs. Nominal GDP Growth, Canada 2022--2023}

\begin{center}
\begin{tabular}{lrrrr}
\toprule
Year & Nominal GDP (\$B) & GDP Deflator & Real GDP (\$B, 2017\$) & Real Growth (\%) \\
\midrule
2022 & 2{,}721 & 110.3 & 2{,}467 & 3.4 \\
2023 & 2{,}879 & 115.2 & 2{,}499 & 1.3 \\
\bottomrule
\end{tabular}
\end{center}

\noindent\textit{Calculation for 2023 real GDP:}
\[
\text{Real GDP}_{2023} = \frac{2{,}879}{115.2/100} = \frac{2{,}879}{1.152} \approx 2{,}499 \text{ billion (2017 dollars)}
\]

\noindent\textit{Calculation for real growth rate:}
\[
g_{2023} = \frac{2{,}499 - 2{,}467}{2{,}467} \times 100 = \frac{32}{2{,}467} \times 100 \approx 1.3\%
\]

\noindent\textit{Interpretation:} Nominal GDP rose 5.8\% between 2022
and 2023, but more than four percentage points of that increase reflected
higher prices rather than more output. Real growth was only 1.3\% ---
an economy moving at a moderate pace, consistent with Canada's estimated
potential growth rate of approximately 1.5--2.0\% per year in the
mid-2020s.

%----------------------------------------------------------
\subsection{GDP Per Capita: A Measure of Living Standards}
%----------------------------------------------------------

Total GDP tells us the size of the economy; it tells us very little about
the typical person's standard of living. An economy of 1.4 billion people
will have a larger GDP than one of 40 million people, even if each of its
citizens is far poorer. For cross-country and cross-time comparisons
of living standards, economists use \term{GDP per capita} --- total output
divided by population:

\begin{equation}
  \text{GDP per capita}_t = \frac{Y_t}{N_t}
  \label{eq:gdp_per_capita}
\end{equation}

\noindent where $N_t$ is the total population.

\noindent\textbf{Worked Example 1.3 --- Canada vs. United States GDP per Capita, 2024}

\begin{center}
\renewcommand{\arraystretch}{1.3}
\begin{tabular}{lrrr}
\toprule
Country & Nominal GDP (USD B) & Population (M) & GDP per Capita (USD) \\
\midrule
United States & 28{,}781 & 337.0 & 85{,}406 \\
Canada        & 2{,}221  &  41.5 & 53{,}518 \\
\midrule
\textbf{Canada--US Gap} & & & \textbf{--\$31{,}888 (--37.3\%)} \\
\bottomrule
\end{tabular}
\end{center}

\noindent\textit{Note:} Canada's nominal GDP is converted to USD at
approximate 2024 exchange rate of 0.74 USD/CAD. GDP in USD thousands:
\(\$2{,}221\text{B} \times 0.74 = \$1{,}643\text{B}\)\ USD; \(\$1{,}643\text{B} / 41.5\text{M} = \$39{,}590\) USD. For PPP-adjusted comparison (which removes exchange rate distortions), the gap narrows but Canada still trails the US by approximately 20--25\%.

This income gap with the United States is one of the most important and
persistent facts about the Canadian economy. It is not primarily explained
by natural resources, education, or geography. A major contributor is
a productivity gap --- Canadians produce less output per hour worked than
Americans --- and that productivity puzzle is the subject of
Chapter~\ref{ch:technology}.

\begin{figure}[H]
\centering
\begin{tikzpicture}
\begin{axis}[
  ybar,
  bar width=0.55cm,
  width=13cm, height=7.5cm,
  symbolic x coords={Italy, Japan, {OECD Avg}, France, UK, Canada,
                     Germany, Australia, Switzerland, USA},
  xtick=data,
  x tick label style={rotate=35, anchor=east, font=\small},
  ylabel={GDP per Capita, PPP-Adjusted (USD thousands), 2024},
  ylabel style={font=\small},
  ymin=0, ymax=100,
  ytick={0,20,40,60,80,100},
  yticklabel style={font=\small},
  ymajorgrids=true,
  grid style={dashed, gray!25},
  enlarge x limits=0.08,
  nodes near coords,
  nodes near coords style={font=\footnotesize, rotate=90, anchor=west},
  every axis plot/.append style={fill=unbcgreen!70, draw=unbcgreen!90}
]
\addplot coordinates {
  (Italy,        46.0)
  (Japan,        52.0)
  ({OECD Avg},   52.5)
  (France,       56.5)
  (UK,           55.0)
  (Canada,       57.0)
  (Germany,      62.0)
  (Australia,    64.5)
  (Switzerland,  78.0)
  (USA,          85.4)
};
\end{axis}
\end{tikzpicture}
\caption{GDP per Capita (PPP-Adjusted, USD thousands), Selected OECD Countries, 2024.
Canada's living standard exceeds the OECD average and surpasses France, the UK,
Japan, and Italy, but trails Germany, Australia, Switzerland, and the United
States. The large Canada--US gap of approximately 37\%, despite similar geography,
institutions, and education levels, reflects a structural productivity difference
explored in Chapter~\ref{ch:technology}.}
\label{fig:gdp_per_capita}
\source{OECD National Accounts at a Glance, 2024; World Bank World Development Indicators.}
\end{figure}

%----------------------------------------------------------
\subsection{Sectoral Structure: The Service Economy}
%----------------------------------------------------------

Canada's economy is predominantly a \textbf{service economy}. Services ---
finance, real estate, retail, healthcare, professional services, education,
and public administration --- account for approximately 77\% of GDP. This
pattern is common among high-income countries: as economies develop, labour
and capital shift from agriculture to manufacturing to services.

\begin{figure}[H]
\centering
\begin{tikzpicture}
\begin{axis}[
  ybar,
  bar width=0.50cm,
  width=14cm, height=7cm,
  symbolic x coords={{Retail}, {Public Admin}, {Healthcare}, {Construction},
                     {Transport}, {Prof.\ Svcs}, {Mining/Oil}, {Manufacturing},
                     {Finance}, {Real Estate}},
  xtick=data,
  x tick label style={rotate=38, anchor=east, font=\small},
  ylabel={Share of GDP (\%)},
  ylabel style={font=\small},
  ymin=0, ymax=16,
  ytick={0,2,4,6,8,10,12,14,16},
  yticklabel style={font=\small},
  ymajorgrids=true,
  grid style={dashed, gray!25},
  enlarge x limits=0.06,
  every axis plot/.append style={fill=unbcgreen!65, draw=unbcgreen!85}
]
\addplot coordinates {
  ({Retail},       5.5)
  ({Public Admin}, 6.0)
  ({Healthcare},   6.2)
  ({Construction}, 7.3)
  ({Transport},    5.4)
  ({Prof.\ Svcs},  7.0)
  ({Mining/Oil},   9.1)
  ({Manufacturing},10.2)
  ({Finance},      7.8)
  ({Real Estate},  13.5)
};
\end{axis}
\end{tikzpicture}
\caption{Canada's Major Industries as a Share of GDP, 2024.
Real estate dominates --- reflecting high property values and the
construction and rental sectors --- while manufacturing and mining/oil
and gas follow. Mining and oil and gas (9.1\%) punches above its weight
politically relative to its actual GDP share, a tension explored in
Chapter~\ref{ch:climate}.
The service-heavy structure means Canada's business cycle is largely
driven by consumer spending, housing, and financial conditions --- all
channels through which monetary policy operates powerfully.}
\label{fig:gdp_by_industry}
\source{Statistics Canada, Table 36-10-0401-01, GDP at basic prices by industry, 2024.}
\end{figure}

Critically, within the goods-producing sector, \textbf{natural resources}
(oil and gas, mining, forestry, and agriculture) play a disproportionately
large role relative to most comparable advanced economies. Canada's oil sands
in Alberta produce over 3 million barrels of crude oil per day. British
Columbia's forest products sector exports billions of dollars of softwood lumber
annually. Saskatchewan supplies approximately 40\% of the world's potash.
These resources generate export revenues, government royalties, and employment ---
but they also make Canada's economy sensitive to global commodity price cycles
and they are central to the climate transition challenge explored in
Chapter~\ref{ch:climate}.

%----------------------------------------------------------
\subsection{Regional Economic Geography}
%----------------------------------------------------------

Canada is a vast country, and its economy is not evenly distributed. The
regional concentration of economic activity creates sharp differences in
employment conditions, wages, housing markets, and fiscal capacity across
provinces --- differences that make national averages sometimes misleading
for understanding local economic conditions.

\begin{figure}[H]
\centering
\begin{tikzpicture}
\begin{axis}[
  ybar,
  bar width=0.45cm,
  width=14cm, height=7cm,
  symbolic x coords={{Terr.}, {PEI}, {N\!L}, {NB}, {NS},
                     {MB}, {SK}, {BC}, {AB}, {QC}, {ON}},
  xtick=data,
  x tick label style={font=\small},
  ylabel={Share of National GDP (\%)},
  ylabel style={font=\small},
  ymin=0, ymax=44,
  ytick={0,5,10,15,20,25,30,35,40},
  yticklabel style={font=\small},
  ymajorgrids=true,
  grid style={dashed, gray!25},
  enlarge x limits=0.06,
  nodes near coords,
  nodes near coords style={font=\footnotesize, rotate=90, anchor=west},
  every axis plot/.append style={fill=policyblue!60, draw=policyblue!80}
]
\addplot coordinates {
  ({Terr.}, 0.5)
  ({PEI},   0.4)
  ({N\!L},  1.2)
  ({NB},    1.4)
  ({NS},    1.6)
  ({MB},    3.0)
  ({SK},    3.1)
  ({BC},   13.6)
  ({AB},   15.9)
  ({QC},   21.1)
  ({ON},   38.9)
};
\end{axis}
\end{tikzpicture}
\caption{Share of National GDP by Province and Territory, 2024.
Ontario and Quebec together account for 60\% of Canadian economic output.
Alberta, despite having one-quarter of Ontario's population, generates 15.9\%
of GDP --- reflecting the outsized economic weight of the oil sands. The three
territories combined produce 0.5\% of GDP despite covering 40\% of Canada's
land mass. This extreme geographic concentration of economic activity motivates
the equalization payment system examined in Chapter~\ref{ch:fiscal}.}
\label{fig:gdp_by_province}
\source{Statistics Canada, Provincial and Territorial Economic Accounts, 2024.}
\end{figure}

The regional diversity of Canada's economy creates significant policy
challenges. When oil prices fall, Alberta's provincial revenue collapses
--- but Ontario's economy is largely unaffected. When the Bank of Canada
raises interest rates to combat national inflation, the impact on heavily
indebted households in Toronto and Vancouver is far more severe than on
renters in Moncton or resource workers in Fort McMurray. National policy
instruments applied to a regionally diverse economy inevitably produce
unequal regional consequences.

\begin{keytakeaway}
Canada's economy is large in aggregate but regionally unequal: Ontario
and Quebec dominate output; Alberta dominates resources; and the Atlantic
provinces and territories depend heavily on federal transfers. National
economic statistics always need to be disaggregated by region for a
complete picture.
\end{keytakeaway}

%==========================================================
\section{Canada's Trade Dependence}
\label{sec:trade_dependence}
%==========================================================

Canada is one of the most trade-dependent large economies in the world.
\term{Trade openness} --- defined as total trade (exports plus imports) as
a share of GDP --- measures how integrated an economy is with the global
trading system:

\begin{equation}
  \text{Trade Openness} = \frac{X + M}{Y} \times 100
  \label{eq:trade_openness}
\end{equation}

For Canada in 2024:
\[
\text{Trade Openness} = \frac{897 + 918}{3{,}001} \times 100 = \frac{1{,}815}{3{,}001} \times 100 \approx 60.5\%
\]

This means that the sum of Canada's exports and imports is equivalent
to more than 60\% of total GDP --- an extraordinary degree of integration
for an economy of Canada's size. By comparison, the United States, with
a much larger domestic market, has a trade openness ratio of approximately
27\%.

\begin{figure}[H]
\centering
\begin{tikzpicture}
\begin{axis}[
  ybar,
  bar width=0.6cm,
  width=12cm, height=7cm,
  symbolic x coords={USA, Japan, Italy, France, UK, Canada, Germany},
  xtick=data,
  x tick label style={font=\small},
  ylabel={Trade Openness: $(X+M)/Y \times 100$ (\%)},
  ylabel style={font=\small},
  ymin=0, ymax=105,
  ytick={0,20,40,60,80,100},
  yticklabel style={font=\small},
  ymajorgrids=true,
  grid style={dashed, gray!25},
  enlarge x limits=0.10,
  nodes near coords,
  every node near coord/.append style={font=\footnotesize},
  every axis plot/.append style={fill=exampleteal!65, draw=exampleteal!85}
]
\addplot coordinates {
  (USA,     27.0)
  (Japan,   38.0)
  (Italy,   56.0)
  (France,  58.0)
  (UK,      62.0)
  (Canada,  60.5)
  (Germany, 88.0)
};
\end{axis}
\end{tikzpicture}
\caption{Trade Openness --- $(X+M)/Y \times 100$ --- G7 Countries, 2024.
Canada's trade openness (60.5\%) is among the highest in the G7, surpassed
only by Germany. The US is far less trade-dependent because its enormous
domestic market satisfies most of its own demand. Canada's high trade
dependence means it is especially vulnerable to US tariff policy, exchange
rate movements, and global commodity price shocks --- themes developed in
Chapter~\ref{ch:trade}.}
\label{fig:trade_openness}
\source{World Bank World Development Indicators; Statistics Canada, Table 12-10-0121-01.}
\end{figure}

The most critical structural fact about Canada's trade is its \textbf{extreme
concentration} in a single market. The United States absorbs approximately
75\% of Canada's merchandise exports, making Canada more dependent on one
trading partner than almost any other advanced economy in the world. This
concentration is a source of enormous prosperity --- geographic proximity,
a shared language and legal tradition, and preferential market access under
CUSMA lower transaction costs dramatically --- but it is also a source of
strategic vulnerability.

\begin{table}[H]
\centering
\caption{Canada's Top 10 Trading Partners, 2024}
\label{tab:trading_partners}
\renewcommand{\arraystretch}{1.35}
\begin{tabular}{lrrr}
\toprule
\textbf{Country} & \textbf{Export Share (\%)} & \textbf{Import Share (\%)}
                 & \textbf{Trade Balance (\$B~CAD)} \\
\midrule
United States    & 75.4 & 51.2 & +130.2 \\
China            &  4.2 & 12.8 & $-$37.9 \\
United Kingdom   &  2.8 &  2.1 & $+$3.1 \\
Japan            &  2.1 &  2.8 & $-$4.0 \\
Germany          &  0.8 &  2.8 & $-$7.1 \\
Mexico           &  0.8 &  4.4 & $-$13.8 \\
South Korea      &  0.8 &  1.5 & $-$2.3 \\
Netherlands      &  1.2 &  0.6 & $+$2.2 \\
India            &  0.7 &  1.0 & $-$1.1 \\
France           &  0.6 &  1.2 & $-$2.4 \\
\midrule
\textbf{All others} & 10.6 & 19.6 & $-$30.5 \\
\bottomrule
\end{tabular}
\source{Statistics Canada, Table 12-10-0121-01; Global Affairs Canada,
Trade Data Online, 2024.}
\end{table}

\noindent Table~\ref{tab:trading_partners} reveals two important asymmetries.
First, Canada runs a substantial trade surplus with the United States (\$130B),
driven almost entirely by energy exports and auto-sector trade. Second, Canada
runs trade deficits with most other major partners, including China and Germany.
The concentration of the surplus in one country and one set of sectors creates
the vulnerability analysed in Chapter~\ref{ch:trade}: a shift in US trade
policy --- such as the imposition of tariffs on Canadian steel, lumber, or
automobiles --- can rapidly transform Canada's trade balance from surplus
to deficit.

%==========================================================
\section{Canada's Major Economic Institutions}
\label{sec:institutions}
%==========================================================

Economic outcomes do not emerge spontaneously from markets alone. They are
shaped by institutions --- the formal rules, organizations, and governance
frameworks within which economic decisions are made. Understanding Canada's
key economic institutions is a prerequisite for understanding how economic
policy is made and why its effects are what they are.

%----------------------------------------------------------
\subsection{The Bank of Canada}
%----------------------------------------------------------

The \term{Bank of Canada} is Canada's central bank. Established under the
\textit{Bank of Canada Act} in 1934, it is structured as a Crown corporation
that operates at formal arm's length from the elected federal government.
In practice, this independence is one of the most consequential institutional
features of the Canadian economy.

The Bank of Canada's primary mandate is to \textbf{promote the economic and
financial welfare of Canada}. Since 1991, this mandate has been operationalized
through a formal \term{inflation-targeting framework}, renewed most recently
in 2021: the Bank commits to keeping CPI inflation within a \textbf{1--3\%
control range}, with a \textbf{2\% midpoint target}.

The Bank conducts monetary policy primarily through its \term{policy interest
rate} --- the target for the \term{overnight rate}, the interest rate at which
major financial institutions lend to each other on an overnight basis. When
the Bank raises the overnight rate, it becomes more expensive to borrow
throughout the financial system. Mortgage rates rise. Business loans cost more.
The exchange rate tends to appreciate. All of these effects reduce spending
and investment, which eventually reduces inflation. When the Bank lowers the
overnight rate, the opposite occurs.

\begin{canexample}{The 2022--2024 Rate-Hike Cycle: The Fastest in Modern History}
Between March 2022 and July 2023, the Bank of Canada raised its policy
interest rate from 0.25\% to 5.00\% --- an increase of 475 basis points
(4.75 percentage points) in just 16 months. This was the most aggressive
monetary tightening cycle in the Bank of Canada's modern history, undertaken
to combat inflation that had reached 8.1\% in June 2022, the highest reading
since September 1983.

The rate hike cycle illustrates the transmission mechanism of monetary policy
in practice. By late 2023, inflation had fallen to approximately 3\%. By mid-2024,
it had approached the 2\% target. But the rate hikes also raised mortgage
payments for millions of Canadian households --- an estimated \$700 on average
for variable-rate mortgage holders --- and contributed to a sharp slowdown in
housing market activity. The Bank began cutting rates in June 2024.

This episode is the central case study of Chapter~\ref{ch:inflation}. It also
directly shapes the housing analysis in Chapter~\ref{ch:housing} and the
distributional analysis in Chapter~\ref{ch:inequality}.
\end{canexample}

%----------------------------------------------------------
\subsection{Department of Finance Canada}
%----------------------------------------------------------

The \term{Department of Finance Canada} is the federal government ministry
responsible for \term{fiscal policy} --- the use of government spending and
taxation to influence the economy. The Minister of Finance presents an
annual federal budget, typically in March, which outlines revenue plans,
spending priorities, and projections for the federal deficit or surplus.

Unlike the Bank of Canada, the Department of Finance operates under direct
democratic accountability: its decisions are made by elected politicians and
must pass through Parliament. This creates a fundamentally different set of
incentives than those governing the Bank of Canada, and it is one reason why
monetary and fiscal policy are deliberately separated in Canada's institutional
design.

Key agencies operating alongside Finance Canada include:

\begin{itemize}
  \item \textbf{Parliamentary Budget Officer (PBO):} An independent officer
        of Parliament established in 2006, mandated to provide non-partisan
        economic and fiscal analysis. The PBO has become an important check
        on government budget projections, regularly producing independent
        assessments that sometimes differ substantially from government
        forecasts.
  \item \textbf{Canada Revenue Agency (CRA):} Administers tax collection
        and enforces tax law. Revenue collection is the CRA's core function;
        it does not make tax policy, which remains with Finance Canada.
  \item \textbf{Employment and Social Development Canada (ESDC):} Administers
        transfer programs including Employment Insurance (EI), the Canada
        Pension Plan (CPP), the Canada Child Benefit (CCB), and Old Age
        Security (OAS/GIS). These programs collectively transfer hundreds
        of billions of dollars per year from taxpayers to households.
\end{itemize}

%----------------------------------------------------------
\subsection{Statistics Canada}
%----------------------------------------------------------

\term{Statistics Canada} (\textit{Statistique Canada}) is the national
statistical agency, established under the \textit{Statistics Act}. It
is the primary data source for the analysis conducted throughout this
textbook. Its publications include:

\begin{itemize}
  \item \textbf{Consumer Price Index (CPI):} Published monthly; the primary
        measure of inflation (Table 18-10-0004-01). Analysed in depth in
        Chapter~\ref{ch:inflation}.
  \item \textbf{Labour Force Survey (LFS):} Published monthly; measures the
        employment rate, unemployment rate, and labour force participation
        rate (Table 14-10-0287-01). Central to Chapter~\ref{ch:labour}.
  \item \textbf{National Income and Expenditure Accounts:} Quarterly GDP by
        expenditure, income, and industry (Table 36-10-0104-01). The
        foundation for Chapters~\ref{ch:introduction} through \ref{ch:fiscal}.
  \item \textbf{Survey of Consumer Finances / Canadian Income Survey:}
        Annual data on household income distribution, poverty, and inequality
        (Table 11-10-0190-01). Central to Chapter~\ref{ch:inequality}.
  \item \textbf{Census of Population:} Every five years; provides
        comprehensive demographic, labour market, and income data. The 2021
        Census is heavily used in Chapters~\ref{ch:inequality}
        and \ref{ch:immigration}.
\end{itemize}

Statistics Canada operates with formal independence from political direction
in its data collection and methodological decisions. The integrity of its
data is a fundamental public good: without reliable statistics, evidence-based
policymaking is impossible.

%----------------------------------------------------------
\subsection{Canada Mortgage and Housing Corporation (CMHC)}
%----------------------------------------------------------

The \term{Canada Mortgage and Housing Corporation} (CMHC) is a federal
Crown corporation with two distinct functions. First, it operates a
\term{mortgage insurance} program: Canadians who purchase a home with
a down payment of less than 20\% must purchase CMHC mortgage insurance,
which protects the lender against default. This insurance program
allows lenders to extend mortgages to buyers with limited savings, making
homeownership accessible to a broader population. Second, CMHC is the
primary source of Canadian housing market data --- including housing starts,
rental vacancy rates, affordability indices, and housing market outlooks ---
and it is the implementing agency for the federal government's housing programs.
CMHC is central to the analysis in Chapter~\ref{ch:housing}.

%----------------------------------------------------------
\subsection{Provincial Governments}
%----------------------------------------------------------

In Canada's federal system, provincial governments bear primary responsibility
for several of the largest areas of economic policy:

\begin{itemize}
  \item \textbf{Healthcare:} Provinces operate and fund the publicly administered
        insurance plans mandated by the \textit{Canada Health Act}. Healthcare
        is typically the single largest provincial spending item --- consuming
        35--45\% of provincial budgets.
  \item \textbf{Education:} Provinces regulate and largely fund education from
        kindergarten through post-secondary.
  \item \textbf{Labour market regulation:} Minimum wages, occupational safety,
        labour relations, and most employment standards are set provincially.
        This is why minimum wages vary from \$13.85 in some Atlantic provinces
        to \$17.40 in British Columbia.
  \item \textbf{Natural resources:} Under the Constitution Act, provinces own
        the natural resources within their borders. Provincial royalty regimes
        for oil, gas, and minerals are major determinants of resource sector
        economics.
  \item \textbf{Housing regulation:} Zoning, land use planning, and building
        codes are primarily municipal functions delegated by provinces. This
        makes provinces central players in the housing affordability crisis
        analysed in Chapter~\ref{ch:housing}.
\end{itemize}

\begin{table}[H]
\centering
\caption{Canada's Major Economic Institutions: Mandate, Tools, and Independence}
\label{tab:institutions}
\renewcommand{\arraystretch}{1.35}
\begin{tabular}{p{3.1cm}p{3.6cm}p{2.8cm}p{2.8cm}}
\toprule
\textbf{Institution} & \textbf{Primary Mandate} & \textbf{Main Tools}
                     & \textbf{Independence Level} \\
\midrule
Bank of Canada
  & Low and stable inflation (2\% target); financial stability
  & Overnight rate; quantitative easing; financial regulation
  & High --- Crown corporation, arm's length from government \\
\addlinespace
Department of Finance Canada
  & Fiscal policy; budget; taxation
  & Government spending; tax rates; transfer programs
  & Low --- elected minister, accountable to Parliament \\
\addlinespace
Statistics Canada
  & Collect, produce, and disseminate national statistics
  & Surveys; censuses; administrative data
  & High --- statistical independence protected by law \\
\addlinespace
CMHC
  & Housing finance; mortgage insurance; housing policy
  & Mortgage insurance; housing programs; market data
  & Moderate --- Crown corporation with federal mandate \\
\addlinespace
Parliamentary Budget Officer
  & Independent fiscal and economic analysis
  & Economic modelling; budget assessments; reports
  & High --- independent officer of Parliament \\
\addlinespace
Provincial Governments
  & Healthcare; education; labour; natural resources
  & Spending; regulation; taxation; royalties
  & Full --- constitutionally sovereign within their jurisdiction \\
\bottomrule
\end{tabular}
\source{Bank of Canada Act; Statistics Act; Canada Mortgage and Housing
Corporation Act; Constitution Act, 1867.}
\end{table}

%==========================================================
\section{The Policy Toolkit: Monetary and Fiscal Policy}
\label{sec:policy_toolkit}
%==========================================================

Two broad categories of macroeconomic policy tools are the primary instruments
through which governments and central banks attempt to influence the
performance of the economy. Understanding the distinction between them ---
and the interaction between them --- is foundational for every topic in
this course.

%----------------------------------------------------------
\subsection{Monetary Policy}
%----------------------------------------------------------

\begin{defbox}{Monetary Policy}
\textbf{Monetary policy} refers to the actions taken by a central bank
to influence the money supply, interest rates, and credit conditions,
with the ultimate goals of maintaining price stability and supporting
sustainable economic growth.

\textit{Economic intuition:} When the central bank raises its policy
interest rate, borrowing becomes more expensive throughout the economy.
Households take out smaller mortgages. Firms postpone capital investment.
The exchange rate tends to appreciate, making imports cheaper and exports
less competitive. These effects reduce aggregate demand and, over time,
inflation. The reverse occurs when the Bank cuts rates.

\textit{Key constraint --- the zero lower bound:} The policy rate cannot
be reduced significantly below zero (depositors would withdraw cash and
hold it physically rather than accept negative returns on deposits). This
creates an asymmetry: the Bank can always raise rates to fight inflation,
but its ability to cut rates to fight recession is constrained at zero.
After 2008 and again in 2020, the Bank cut its rate to 0.25\% --- its
effective lower bound --- and was forced to rely on other tools (forward
guidance and, to a limited extent, quantitative easing) to provide
additional stimulus.

\textit{Canadian example:} The Bank cut the overnight rate from 1.75\%
to 0.25\% in March 2020 in response to COVID-19, then raised it to 5.00\%
between March 2022 and July 2023 to fight inflation --- the full range of
monetary policy in action within a four-year period.
\end{defbox}

The primary transmission channels through which a Bank of Canada interest
rate change affects the broader economy are:

\begin{enumerate}
  \item \textbf{Credit cost channel:} Higher overnight rates are passed
        through to mortgage rates, car loan rates, and business borrowing
        costs, reducing household and business spending.
  \item \textbf{Asset price channel:} Higher interest rates reduce the
        present value of future asset cash flows, lowering stock prices
        and real estate values, which reduces household wealth and
        consumption (the ``wealth effect'').
  \item \textbf{Exchange rate channel:} Higher Canadian interest rates
        attract foreign capital seeking better returns, increasing demand
        for Canadian dollars and appreciating the exchange rate. A stronger
        Canadian dollar makes imports cheaper (reducing inflation) and
        exports more expensive (reducing demand for Canadian goods abroad).
  \item \textbf{Expectations channel:} Perhaps most importantly, when
        the Bank credibly commits to fighting inflation, households and
        firms adjust their price-setting and wage demands downward,
        reducing inflationary momentum even before interest rates fully
        transmit through the system.
\end{enumerate}

%----------------------------------------------------------
\subsection{Fiscal Policy}
%----------------------------------------------------------

\begin{defbox}{Fiscal Policy}
\textbf{Fiscal policy} refers to the use of government spending and
taxation to influence macroeconomic conditions.

\textit{Economic intuition:} When aggregate demand is insufficient to
support full employment, the government can increase its own spending
(shifting the aggregate demand curve rightward) or cut taxes (increasing
household disposable income and thus consumption). Conversely, when
the economy is overheating, fiscal restraint --- cutting spending or
raising taxes --- can reduce inflationary pressure.

\textit{Key constraint --- the deficit and debt:} When the government
spends more than it collects in taxes, it runs a \term{budget deficit},
which must be financed by borrowing. The accumulation of past deficits
is the \term{national debt}. Sustained high deficits raise the stock
of debt relative to GDP, increasing interest payments and potentially
crowding out private investment. The fiscal arithmetic of debt
sustainability is explored fully in Chapter~\ref{ch:fiscal}.

\textit{Canadian example:} In 2020--21, the federal government ran a
deficit of \$327.7~billion --- 14.9\% of GDP --- as it deployed the
Canada Emergency Response Benefit (CERB), the Canada Emergency Wage
Subsidy (CEWS), and other programs to protect household incomes during
the pandemic shutdown. This was the largest single-year deficit in
Canadian history. It was deliberate and effective in preventing a
financial collapse, but it also raised Canada's federal debt-to-GDP
ratio from approximately 31\% to 49\% within two years.
\end{defbox}

%----------------------------------------------------------
\subsection{The Interaction Between Monetary and Fiscal Policy}
%----------------------------------------------------------

Monetary and fiscal policy are operated by different institutions with
different mandates, but they interact in important ways. When fiscal
policy is highly expansionary --- as it was in 2020--21 --- it can
increase aggregate demand and put upward pressure on inflation, requiring
a tighter monetary policy response. Conversely, when monetary policy is
very tight --- as it was in 2022--23 --- higher interest rates increase
the government's debt servicing costs, creating fiscal pressure.

The period 2020--2023 illustrates this interaction vividly. Emergency
fiscal spending in 2020--21 was appropriate and necessary. But the
combination of massive fiscal stimulus and near-zero interest rates
through 2021 contributed to the demand-side component of the inflation
surge that the Bank of Canada then had to combat aggressively in 2022--23.
Understanding this interaction is essential for evaluating both the
monetary policy response (Chapter~\ref{ch:inflation}) and the fiscal
policy choices (Chapter~\ref{ch:fiscal}).

\begin{figure}[H]
\centering
\begin{tikzpicture}[
  node distance=1.6cm and 2.8cm,
  mainbox/.style={rectangle, rounded corners=5pt, draw=unbcgreen,
                  fill=unbcgreenlight, minimum width=3.2cm,
                  minimum height=1.1cm, align=center, font=\small\bfseries},
  midbox/.style={rectangle, rounded corners=4pt, draw=policyblue,
                 fill=policyblulight, minimum width=3.0cm,
                 minimum height=0.95cm, align=center, font=\small},
  finalbox/.style={rectangle, rounded corners=5pt, draw=unbcgreen,
                   fill=unbcgreenmid, minimum width=7cm,
                   minimum height=1.0cm, align=center, font=\small\bfseries},
  arrow/.style={-Stealth, thick, color=unbcgreen!80}
]
\node[mainbox] (boc)     {Bank of Canada\\(Monetary Policy)};
\node[mainbox, right=3.6cm of boc] (fin) {Dept.\ of Finance\\(Fiscal Policy)};

\node[midbox, below=1.6cm of boc] (ir)   {Overnight Rate\\\& Credit Costs};
\node[midbox, below=1.6cm of fin] (gt)   {Spending, Taxes\\\& Transfers};

\node[midbox, below=1.6cm of ir]  (inv)  {Investment,\\ Consumption,\\ Exchange Rate};
\node[midbox, below=1.6cm of gt]  (ad2)  {Household\\ Disposable\\ Income};

\node[finalbox, below=4.8cm of $(boc)!0.5!(fin)$] (out)
     {Aggregate Demand $\longrightarrow$ GDP Growth, Employment, Inflation};

\draw[arrow] (boc) -- (ir);
\draw[arrow] (fin) -- (gt);
\draw[arrow] (ir)  -- (inv);
\draw[arrow] (gt)  -- (ad2);
\draw[arrow] (inv) -- (out);
\draw[arrow] (ad2) -- (out);

\draw[dashed, color=gray, thick, <->]
  ([xshift=0.5cm]boc.east) -- ([xshift=-0.5cm]fin.west)
  node[midway, above, font=\footnotesize, color=gray]
  {Coordination\,/\,Conflict};

\node[above=0.4cm of boc, font=\small\bfseries\color{unbcgreen}]
  {CANADA'S MACROECONOMIC POLICY ARCHITECTURE};
\end{tikzpicture}
\caption{Canada's Macroeconomic Policy Architecture.
Monetary policy (left branch) and fiscal policy (right branch) each
affect aggregate demand through distinct channels, ultimately influencing
GDP growth, employment, and inflation. The dashed bidirectional arrow
indicates that the two policy levers can work in the same direction
(as during the COVID stimulus) or in opposing directions (as in 2022--23,
when fiscal policy remained expansionary even as the Bank of Canada
was tightening). Understanding this interaction is one of the central
analytical tasks of contemporary macroeconomics.}
\label{fig:policy_framework}
\end{figure}

\begin{policydebate}{Should the Bank of Canada Be Independent from the Federal Government?}

\textbf{The Issue:} Central bank independence --- the insulation of
monetary policy decisions from direct political control --- is a cornerstone
of modern macroeconomic governance. But independence is not constitutionally
guaranteed in Canada. The \textit{Bank of Canada Act} allows the federal
government, in exceptional circumstances, to direct the Bank. Should this
arm's-length relationship be strengthened, weakened, or left as is?

\textbf{The Case for Independence:}
Economic theory and empirical evidence both support central bank independence.
The fundamental problem is \term{time inconsistency}: elected governments
have incentives to expand the economy before elections, even at the cost of
higher future inflation. If markets expect governments to exploit monetary
policy for electoral gain, they will build that expectation into wages and
prices, requiring permanently higher unemployment to achieve any given
inflation rate. An independent central bank can credibly commit to price
stability in a way that elected governments cannot --- and that credibility
is itself enormously valuable. Cross-country evidence consistently shows
that countries with more independent central banks have lower average
inflation without worse unemployment or growth outcomes
\citep[see][]{alesina1993,cukierman1992}.

\textbf{The Case Against Full Independence:}
Critics argue that central bank independence raises a \term{democratic
accountability} problem. When the Bank of Canada raises interest rates
by 475 basis points in 16 months --- dramatically increasing mortgage
costs for millions of Canadians --- this is a consequential distributional
decision that arguably deserves more democratic oversight than it currently
receives. Furthermore, the separation of monetary and fiscal policy can
make macroeconomic management incoherent: if the elected government wishes
to support employment and the central bank is tightening to control
inflation, the two policies work against each other, potentially prolonging
both the high interest rates and the high unemployment.

\textbf{The Canadian Position:}
Canada's current arrangement is a hybrid. The Bank of Canada is independent
in its day-to-day operations, but its inflation target is set jointly
with the Government of Canada in a periodic agreement (most recently renewed
in 2021). This ``joint mandate'' approach attempts to combine the credibility
benefits of independence with some degree of democratic input into the
overarching framework.

\textbf{Your Assessment:}
After completing Chapter~\ref{ch:inflation}, you will be equipped to evaluate
the Bank's rate-hike cycle on its merits. Did the Bank move too aggressively?
Too slowly? At the right pace? These are exactly the kinds of questions where
the distinction between positive analysis (what happened and what were the
effects?) and normative judgement (was this the right policy?) matters most.
\end{policydebate}

%==========================================================
\section{How Economists Think}
\label{sec:how_economists_think}
%==========================================================

Developing the analytical habits of economic reasoning takes time and
deliberate practice. This section introduces the core tools that you will
use throughout the course.

%----------------------------------------------------------
\subsection{Economic Models and Assumptions}
%----------------------------------------------------------

Economists use \term{models} --- carefully constructed simplifications of
reality --- to generate predictions and organize thinking about complex
phenomena. A model deliberately omits details in order to focus on the
mechanisms that most matter for a particular question.

The \term{supply and demand model}, for example, abstracts away from the
unique characteristics of millions of individual buyers and sellers and
focuses instead on the aggregate relationship between price and quantity.
This abstraction is a feature, not a limitation: it allows economists to
make clear, testable predictions about how markets respond to changes in
supply, demand, and policy.

Every model rests on assumptions. Evaluating a model requires understanding
which assumptions are load-bearing (if they are wrong, the model's
predictions fail) and which are harmless simplifications. The assumption
that economic agents ``respond to incentives'' is load-bearing for almost
all of economics. The assumption of perfect information is often harmless
as a first approximation but becomes crucial in analyses of insurance,
financial markets, and healthcare.

\begin{thinkbox}
A map is a model. A road map of Canada omits trees, buildings, and
terrain because those details are irrelevant to the task of navigating
from Prince George to Edmonton. A topographic map of the same region
omits roads but captures elevation, because that information matters
for a hiker. Neither map is ``more correct'' than the other --- they are
optimized for different questions. Economic models work the same way.
A model that is useful for understanding inflation may be useless for
understanding income inequality. The skill is in choosing the right
model for the question at hand.
\end{thinkbox}

%----------------------------------------------------------
\subsection{Positive and Normative Economics}
%----------------------------------------------------------

One of the most important distinctions in economic methodology is between
\term{positive economics} and \term{normative economics}.

\begin{defbox}{Positive and Normative Economics}
\textbf{Positive economics} describes and explains economic phenomena
as they are, without making value judgements. It answers ``what is''
or ``what will happen if.'' Positive claims are, in principle, testable
against empirical evidence.

\textbf{Normative economics} makes value judgements about what
\textit{should} happen. It answers ``what ought to be.'' Normative
claims cannot be resolved by evidence alone --- they also depend on
values and priorities.

\textbf{Examples in the Canadian context:}
\begin{itemize}
  \item \textit{Positive:} ``Canada's CPI inflation peaked at 8.1\% in
        June 2022.'' [Observable fact.]
  \item \textit{Positive:} ``A 25\% US tariff on Canadian auto exports
        would reduce Canadian GDP by an estimated 0.3--0.7 percentage
        points.'' [Testable empirical claim, though uncertain.]
  \item \textit{Normative:} ``The federal carbon tax should be increased
        to \$250 per tonne.'' [Value judgement about the right balance
        between climate and economic costs.]
  \item \textit{Normative:} ``Canada should prioritize reducing inequality
        over maximizing GDP growth.'' [Value judgement about priorities.]
\end{itemize}

Note: many policy debates conflate positive and normative claims. Part
of the discipline of economics is keeping them separate.
\end{defbox}

Policy debates are often confusing precisely because positive and normative
claims are mixed together. When a politician argues that the carbon tax
``kills jobs'' (positive claim --- testable) and ``should be abolished''
(normative claim --- based on values), they are making two distinct types
of argument that require two distinct types of response. Throughout this
textbook, we will be explicit about which claims are positive (and what
the evidence says) and which are normative (and what the competing values are).

%----------------------------------------------------------
\subsection{Opportunity Cost and Trade-offs}
%----------------------------------------------------------

A fundamental principle of economic reasoning is that \textit{all choices
involve trade-offs}. Resources --- time, money, productive capacity, policy
attention --- allocated to one purpose cannot simultaneously be used for
another. The \term{opportunity cost} of any choice is the value of the
next-best alternative foregone.

\begin{equation}
  \text{Opportunity Cost of Choice A} = \text{Value of the best foregone alternative to A}
  \label{eq:opportunity_cost}
\end{equation}

Opportunity cost is not the same as monetary cost. When a university student
works part-time during the school year, the opportunity cost of those hours
worked is not just the forgone leisure but the forgone study time and the
potential impact on grades. When the federal government allocates
\$15~billion to the National Housing Strategy, the opportunity cost is not
just the \$15~billion in budget space but whatever the next-best use of
that money would have been --- whether tax cuts, healthcare investment,
debt repayment, or climate adaptation spending.

\begin{canexample}{The Opportunity Cost of Canada's COVID Fiscal Response}
In 2020--21, the federal government spent approximately \$350 billion in
emergency COVID-19 support. This was financed by borrowing, raising the
federal debt-to-GDP ratio by nearly 18 percentage points. The opportunity
cost of this spending was not only the future interest payments
(approximately \$43 billion per year by 2024--25), but also the fiscal
space foregone for other priorities: a national pharmacare program, a
substantial climate investment fund, or a faster post-pandemic deficit
reduction path. The emergency spending was arguably necessary and worth
its cost --- but a complete analysis must acknowledge the trade-offs.
Chapter~\ref{ch:fiscal} revisits this question with full fiscal data.
\end{canexample}

%----------------------------------------------------------
\subsection{Marginal Analysis}
%----------------------------------------------------------

Economists typically analyse decisions not as binary choices (do this or do
not do this) but at the \term{margin}: how much of something should be done?
\term{Marginal cost} (MC) is the additional cost of one more unit of
an activity. \term{Marginal benefit} (MB) is the additional benefit.
The economically optimal decision equates the two:

\begin{equation}
  \text{Optimal decision:} \quad MB = MC
  \label{eq:mb_mc}
\end{equation}

The marginal analysis framework applies broadly. Should the Bank of Canada
raise interest rates another 25 basis points? The marginal benefit is a
further reduction in inflation; the marginal cost is a further increase
in mortgage payments and a risk of over-tightening into recession. Should
Canada raise the carbon tax from \$65 to \$80 per tonne? The marginal
benefit is additional emissions reduction; the marginal cost is additional
burden on households and firms. In both cases, the right question is not
``should we do more?'' in the abstract, but ``does the marginal benefit
exceed the marginal cost at this particular point?''

%----------------------------------------------------------
\subsection{Correlation and Causation}
%----------------------------------------------------------

Perhaps the most important methodological lesson in empirical economics is
the distinction between \term{correlation} and \term{causation}. Two
variables are correlated if they tend to move together. But correlation
does not establish causation: variable A may cause B, B may cause A, or
a third variable C may cause both A and B.

\begin{thinkbox}
Ice cream sales and drowning deaths are positively correlated --- both rise
in summer. Does eating ice cream cause drowning? Of course not: the third
variable, warm weather, drives both. Identifying this kind of spurious
correlation is easy. In economics, the confounding variables are often
far less obvious.

Consider: provinces with higher minimum wages tend to have higher
unemployment rates. Does this prove that minimum wages cause unemployment?
Not necessarily --- provinces with historically weaker labour markets may
have adopted higher minimum wages as a compensatory policy, or other
factors (industry mix, demographics) may explain both. Establishing
\textit{causal} effects from observational data requires careful
research design, and the results are often contested. We will encounter
this challenge repeatedly throughout the course.
\end{thinkbox}

%----------------------------------------------------------
\subsection{Efficiency and Equity}
%----------------------------------------------------------

A final pair of concepts that runs through every chapter is the distinction
between \term{economic efficiency} and \term{economic equity} (or fairness).
An allocation of resources is \term{efficient} if it is impossible to make
someone better off without making someone else worse off --- a
\term{Pareto optimal} state. Markets, under the right conditions, tend to
produce efficient outcomes. But efficient outcomes need not be equitable:
a distribution in which one person owns everything and everyone else
owns nothing is technically Pareto optimal (you cannot help the poor
without taking from the rich).

Many of the most important policy debates in this course --- should Canada
have a carbon tax? should it have rent control? should it have a wealth tax?
--- involve precisely this tension. A policy that increases efficiency
(total GDP) may worsen equity (distribution of that GDP), or vice versa.
Evaluating these trade-offs requires clear thinking about both what the
evidence says (positive) and what values should guide our priorities
(normative).

%==========================================================
\section{Canada in International Context}
\label{sec:international_context}
%==========================================================

Understanding Canada's contemporary economic situation requires situating
it within the broader international environment. Canada's economic performance
is shaped by global forces --- commodity prices, US monetary policy, global
trade flows, and climate commitments --- that no Canadian institution can
fully control.

\begin{figure}[H]
\centering
\begin{tikzpicture}
\begin{axis}[
  width=14cm, height=7.5cm,
  xlabel={Year},
  ylabel={Real GDP Growth Rate (\%)},
  xlabel style={font=\small},
  ylabel style={font=\small},
  xmin=1999, xmax=2027,
  ymin=-7, ymax=7,
  xtick={2000,2002,2004,2006,2008,2010,2012,2014,2016,2018,2020,2022,2024,2026},
  xticklabel style={rotate=40, anchor=east, font=\footnotesize},
  ytick={-6,-4,-2,0,2,4,6},
  yticklabel style={font=\small},
  ymajorgrids=true,
  grid style={dashed, gray!25},
  axis line style={color=gray!60},
  tick style={color=gray!60},
  legend pos=north east,
  legend style={font=\small},
  clip=false
]
% Reference line at 0%
\draw[dashed, gray, thick] (axis cs:1999,0) -- (axis cs:2027,0);
% Shaded recession bands
\fill[gray!15] (axis cs:2008.5,-7) rectangle (axis cs:2009.5,7);
\fill[gray!15] (axis cs:2019.8,-7) rectangle (axis cs:2020.8,7);
% Real GDP growth line
\addplot[unbcgreen, thick, mark=*, mark size=1.5pt] coordinates {
  (2000, 5.2)(2001, 1.8)(2002, 3.0)(2003, 1.9)(2004, 3.1)
  (2005, 3.2)(2006, 2.6)(2007, 2.1)(2008, 1.0)(2009,-2.9)
  (2010, 3.1)(2011, 3.1)(2012, 1.7)(2013, 2.5)(2014, 2.9)
  (2015, 0.7)(2016, 1.0)(2017, 3.0)(2018, 2.4)(2019, 1.9)
  (2020,-5.2)(2021, 5.0)(2022, 3.4)(2023, 1.2)(2024, 1.5)
  (2025, 1.3)(2026, 1.8)
};
% Potential growth band
\addplot[dataorange, dashed, thick] coordinates {
  (2000,2.0)(2026,2.0)
} node[right, font=\footnotesize, xshift=2pt] {Potential growth ($\approx$2\%)};

% Annotations
\node[font=\footnotesize, align=center, text=gray] at (axis cs:2009,-4.5)
  {Financial\\Crisis\\$-$2.9\%};
\node[font=\footnotesize, align=center, text=gray] at (axis cs:2015.5,-1.5)
  {Oil price\\shock};
\node[font=\footnotesize, align=center, text=gray] at (axis cs:2020,-6.3)
  {COVID-19\\$-$5.2\%};
\node[font=\footnotesize, align=center, text=unbcgreen] at (axis cs:2021,6.2)
  {Rebound\\+5.0\%};
\addlegendentry{Canada Real GDP Growth}
\end{axis}
\end{tikzpicture}
\caption{Canada's Real GDP Growth Rate, 2000--2026. The shaded bands mark
recessionary years. Three episodes stand out: the 2009 financial crisis
contraction ($-$2.9\%), driven by the global banking collapse and the
near-simultaneous US recession; the 2015--16 slowdown driven by the oil
price collapse to below \$30/barrel (US); and the 2020 COVID-19 collapse
($-$5.2\%), the steepest in peacetime Canadian history. The dashed orange
line at 2\% represents Canada's approximate \textit{potential} growth rate ---
the rate consistent with stable unemployment and inflation. Years above this
line indicate excess demand; years below it indicate slack.}
\label{fig:gdp_growth}
\source{Statistics Canada, Table 36-10-0104-01, Chained (2017) dollars.}
\end{figure}

\begin{table}[H]
\centering
\caption{Key Canadian Macroeconomic Indicators, 2015--2026}
\label{tab:macro_indicators}
\renewcommand{\arraystretch}{1.30}
\begin{tabular}{lrrrrrr}
\toprule
\textbf{Year} & \textbf{Real GDP} & \textbf{CPI} & \textbf{Unemp.}
              & \textbf{Bank} & \textbf{Fed.\ Balance} & \textbf{Fed.\ Debt} \\
              & \textbf{Growth (\%)} & \textbf{Infl.\ (\%)} & \textbf{Rate (\%)}
              & \textbf{Rate (\%)} & \textbf{(\$B)} & \textbf{(\% GDP)} \\
\midrule
2015 & 0.7  & 1.1  & 6.9 & 0.50 & $-$2.9   & 31.0 \\
2016 & 1.0  & 1.4  & 7.0 & 0.50 & $-$17.8  & 31.5 \\
2017 & 3.0  & 1.6  & 6.3 & 1.00 & $-$19.0  & 30.8 \\
2018 & 2.4  & 2.3  & 5.8 & 1.75 & $-$14.0  & 30.1 \\
2019 & 1.9  & 1.9  & 5.7 & 1.75 & $-$39.4  & 30.8 \\
2020 & $-$5.2 & 0.7 & 9.5 & 0.25 & $-$327.7 & 49.1 \\
2021 & 5.0  & 3.4  & 7.5 & 0.25 & $-$113.8 & 45.2 \\
2022 & 3.4  & 6.8  & 5.3 & 4.25 & $-$36.0  & 42.4 \\
2023 & 1.2  & 3.9  & 5.4 & 5.00 & $-$40.0  & 42.0 \\
2024 & 1.5  & 2.6  & 6.3 & 4.25 & $-$46.8  & 42.3 \\
2025 & 1.3  & 2.1  & 6.7 & 3.00 & $-$48.0  & 42.5 \\
2026 & 1.8  & 2.0  & 6.5 & 2.75 & $-$30.0  & 42.0 \\
\bottomrule
\end{tabular}
\source{Statistics Canada, Tables 36-10-0104-01, 18-10-0004-01,
14-10-0287-01; Bank of Canada Key Interest Rates; Department of Finance
Canada, Fiscal Reference Tables, 2026. Values for 2025--2026 are
preliminary or projected. Federal balance figures in nominal Canadian
dollars; negative values indicate a deficit.}
\end{table}

\noindent Table~\ref{tab:macro_indicators} is the single most important
reference table in this textbook. It captures the macroeconomic story of
Canada across 12 years in six dimensions. Study it carefully. Notice that
2020 is the anomaly year on every dimension simultaneously: growth collapses,
inflation vanishes, unemployment spikes, the Bank rate hits its floor, and
the budget deficit reaches a historic extreme. Notice that 2022 is the
year of the great inflation --- 6.8\%, the highest since 1991. Notice that
the Bank rate and CPI inflation move together through 2022--2024, reflecting
the monetary policy response. And notice that the federal debt-to-GDP ratio
rose by nearly 18 percentage points in a single year (2019 to 2020) and
has not fully recovered --- a fiscal legacy that Chapter~\ref{ch:fiscal}
examines in depth.

%==========================================================
\section{Case Studies}
\label{sec:ch1_casestudies}
%==========================================================

%----------------------------------------------------------
\begin{casestudy}{COVID-19's Economic Shock and Canada's Recovery}

\textbf{Background:} On March 11, 2020, the World Health Organization
declared COVID-19 a global pandemic. Within days, Canadian provinces
ordered the closure of non-essential businesses, schools, and public
spaces. Economic activity collapsed at an unprecedented speed.

\textbf{The Economic Shock:} In April 2020, Canada lost approximately
2~million jobs --- more than were lost in the entire 2008--09 recession.
GDP fell at an annualized rate exceeding 40\% in the second quarter of 2020.
The unemployment rate reached 13.7\% in May 2020. These numbers, while
startling, understate the disruption: millions of additional Canadians
were employed but working zero hours (classified as employed but absent
by the LFS methodology), and the labour force participation rate fell
sharply as many stopped looking for work.

\textbf{Policy Response:} The federal government deployed the Canada Emergency
Response Benefit (CERB), which paid \$2,000 per month to approximately
9 million workers who lost employment income --- more than one in four
Canadian workers. The Canada Emergency Wage Subsidy (CEWS) paid up to
75\% of the wages of employees at firms whose revenues had fallen, at a
total cost exceeding \$100 billion. The Bank of Canada cut its policy
rate to 0.25\% and committed to holding it there until the recovery was
``well underway.'' Combined, the fiscal and monetary response was the
most aggressive in Canadian peacetime history.

\textbf{Recovery:} Canada's recovery was faster than most projections.
By early 2022, employment had fully recovered to pre-pandemic levels.
GDP recovered its lost ground by early 2022 as well. But the recovery
was uneven: high-contact service sectors (restaurants, hotels, live
entertainment) recovered more slowly; women, racialized Canadians, and
low-wage workers experienced longer-lasting earnings impacts; and the
combination of massive fiscal stimulus, low interest rates, and supply
chain disruptions created the inflationary surge of 2022 that reshaped
the economic environment for years afterward.

\textbf{Evaluation:} The COVID episode illustrates both the power and the
limits of macroeconomic policy. The fiscal and monetary response prevented
a financial collapse and protected millions of households. But it also
contributed to inflationary dynamics that the Bank of Canada then had to
combat at significant cost to mortgage-holding households. Understanding
this sequence of events --- pandemic shock, policy response, inflationary
consequence, monetary tightening, housing market correction --- provides
the narrative architecture for Chapters~\ref{ch:inflation} through
\ref{ch:housing}.

\textit{Data reference:} Statistics Canada, Table 14-10-0287-01 (employment);
Table 36-10-0104-01 (GDP); Parliamentary Budget Officer, COVID-19
Fiscal Analysis, 2021.
\end{casestudy}

%----------------------------------------------------------
\begin{casestudy}{Northern British Columbia as a Regional Economy}

\textbf{Background:} The University of Northern British Columbia sits in
Prince George --- the largest city in northern BC, with a population of
approximately 75,000. The regional economy of northern BC illustrates,
in compressed form, many of the structural features and contemporary
challenges of the Canadian economy as a whole.

\textbf{Economic Structure:} Northern BC's economy is built on natural
resources: forestry, mining (gold, copper, coal), and increasingly, natural
gas from the Montney formation and the Peace River region. The completion
of the Trans Mountain Pipeline Expansion (2024) and the LNG Canada facility
near Kitimat (under construction) represent billion-dollar investments that
will shift the regional resource mix toward hydrocarbons. The forest sector,
historically the dominant employer, has been reshaped by mill closures,
automation, and the mountain pine beetle epidemic that devastated the lodgepole
pine forests of the 1990s and 2000s.

\textbf{Labour Market Features:} Northern BC labour markets are characterized
by: above-average wages in resource sectors (unionized sawmill and mine
workers earn substantially more than provincial minimum wage); persistent
labour shortages in healthcare and social services (geographic remoteness
makes recruitment difficult); a significant Indigenous population (First
Nations account for approximately 16\% of the regional population) whose
economic participation and outcomes reflect the historical legacy of
dispossession analysed in Chapter~\ref{ch:inequality}; and a mismatch
between the skills demanded by resource employers and those produced by
the local education system.

\textbf{Contemporary Challenges:} The three most acute economic challenges
in northern BC today mirror those of the national economy: a housing market
that has become significantly less affordable as in-migration from
Metro Vancouver and Alberta has driven up prices; a healthcare system
under fiscal pressure (Northern Health Authority); and the question of
how a resource-dependent regional economy navigates the clean energy
transition while maintaining employment and fiscal revenues. These regional
challenges appear in the course not as peripheral ``local colour'' but as
concrete examples through which national economic forces become tangible.
\end{casestudy}

%----------------------------------------------------------
\begin{casestudy}{Canada's 2022 Inflation Surge: A Turning Point}

\textbf{Background:} Canada had not experienced inflation above 5\% since
1991. A generation of economists, policymakers, and households had come
of age in a low-inflation environment. The 2021--2022 surge was therefore
not merely an economic event but a psychological and institutional stress test.

\textbf{What Happened:} Consumer price inflation rose from 0.7\% in 2020
(deflated by pandemic-related demand collapse) to 3.4\% in 2021 and then
to a peak of 8.1\% in June 2022. By that point, food prices were rising
at over 10\% year-over-year. Shelter costs were rising at over 7\%.
Energy prices --- driven by Russian sanctions following the Ukraine invasion ---
were rising at over 30\% for gasoline and over 20\% for natural gas.

\textbf{Why It Happened:} The inflation surge had multiple, simultaneous causes.
On the demand side: the COVID fiscal stimulus had placed over \$350 billion
into Canadian household bank accounts and business balance sheets, creating
excess demand when supply chains could not expand fast enough to meet it.
Low interest rates through 2021 added further fuel. On the supply side:
global shipping costs rose tenfold; semiconductor shortages constrained
automobile and electronics production; pandemic-related labour shortages
reduced output in food processing, transportation, and construction. The
Russian invasion of Ukraine (February 2022) added an energy and food
commodity price shock on top of already-elevated inflation.

\textbf{Policy Response:} The Bank of Canada raised its policy rate 10 times
between March 2022 and July 2023, from 0.25\% to 5.00\%. By end-2023,
inflation had fallen to 3.4\%; by end-2024, it was approaching 2\%.
But the rate hikes came at significant cost: mortgage payments for variable-rate
holders increased by hundreds of dollars per month; housing affordability
deteriorated further; and economic growth slowed sharply.

\textbf{Lesson:} The 2022 episode illustrates that macroeconomic stability
cannot be taken for granted. Institutional frameworks (like the inflation-targeting
regime) that functioned smoothly for three decades can be tested by the right
combination of shocks. It also illustrates the limits of monetary policy:
the Bank of Canada can raise rates to reduce demand, but it cannot directly
fix supply chain disruptions, energy prices determined in global markets,
or geopolitical shocks. These limits are a recurring theme in
Chapter~\ref{ch:inflation}.
\end{casestudy}

%==========================================================
\section{Economics in Everyday Life}
%==========================================================

\begin{everydaylife}{From the Gas Station to the Central Bank: Tracing Economic Forces in Your Daily Life}
The next time you fill up your car, consider the chain of economic
forces determining the number on the pump.

The \textbf{global oil price} (set in markets in Houston and Amsterdam)
determines the base cost of crude. Canada's \textbf{exchange rate} ---
the value of the Canadian dollar relative to the US dollar --- translates
that global price into Canadian dollars: a 10-cent drop in the loonie
raises the price of a barrel of oil (priced in USD) by approximately
7--8\% in Canadian terms. The \textbf{Bank of Canada's interest rate
decisions} affect the exchange rate by influencing capital flows into and
out of Canada. Federal and provincial \textbf{excise taxes} add
approximately 40--50 cents per litre at the pump --- of which about
17 cents is the federal carbon levy (at the 2024 rate of \$65 per tonne).
\textbf{Refinery margins}, influenced by the number and capacity of
Canadian refineries (a market structure question), add another layer of cost.

Every chapter in this textbook touches on one of these forces. Chapter~\ref{ch:inflation}
examines how energy prices drive CPI inflation. Chapter~\ref{ch:trade}
examines the exchange rate. Chapter~\ref{ch:fiscal} examines excise taxes.
Chapter~\ref{ch:climate} examines the carbon levy. By the end of this course,
the number on that gas pump will be transparent to you in a way it is not
for most Canadians --- not because you will agree with every policy that
produced it, but because you will understand the economic forces behind each
component.
\end{everydaylife}

%==========================================================
\section{Preview of Course Themes}
\label{sec:course_preview}
%==========================================================

The remaining nine chapters of this textbook address the major contemporary
economic issues in the order determined by the pedagogical logic explained
in the preface. Each chapter builds on the institutional and analytical
foundations established here.

\textbf{Chapter~\ref{ch:inflation}: Inflation, Monetary Policy, and the Cost of Living.}
The immediate crisis of 2022--2023 examined in full: CPI measurement and its
biases, the aggregate demand--aggregate supply framework, the Phillips Curve,
the Bank of Canada's transmission mechanism, and the distributional consequences
of both inflation and the interest rate response.

\textbf{Chapter~\ref{ch:housing}: Canada's Housing Crisis.}
The economics of housing markets, the user cost model, the supply constraints
that prevent new construction from responding to demand, the role of
immigration and interest rates on the demand side, and the full range of
federal and provincial policy responses. This chapter flows directly from
Chapter~\ref{ch:inflation}: the interest rate hikes needed to control
inflation were the same rate hikes that deepened the mortgage crisis.

\textbf{Chapter~\ref{ch:labour}: Labour Markets, Employment, and Demographic Change.}
Labour supply and demand, labour force accounting, the Beveridge Curve, the
wage-productivity gap, minimum wage economics, non-standard employment, and
the relationship between an aging workforce and labour market tightness.

\textbf{Chapter~\ref{ch:inequality}: Income Inequality, Poverty, and Indigenous Economic Development.}
The Lorenz Curve and Gini coefficient, trends in Canadian income distribution
since 1980, intergenerational mobility, the gender wage gap, poverty
measurement, and the economic consequences of the historical dispossession
of Indigenous peoples.

\textbf{Chapter~\ref{ch:immigration}: Immigration, Population Growth, and the Canadian Economy.}
Canada's immigration system, the demographic arithmetic of the dependency
ratio, the labour market effects of immigration, the housing demand effects,
immigrant economic integration, and the policy debate over Canada's
historically high immigration targets.

\textbf{Chapter~\ref{ch:fiscal}: Fiscal Policy, Government Debt, and Public Finance.}
The federal budget process, deficit and debt arithmetic, the fiscal multiplier,
the COVID-era fiscal expansion, healthcare spending and its demographic drivers,
fiscal federalism and equalization, and the question of long-run fiscal
sustainability.

\textbf{Chapter~\ref{ch:trade}: International Trade, Supply Chains, and Canada-US Relations.}
Comparative advantage, the gains from trade, Canada's merchandise trade
patterns, CUSMA and its predecessors, the exchange rate and its effects,
the costs of tariff protection, supply chain vulnerabilities, and the
catastrophic potential of a US-Canada trade conflict.

\textbf{Chapter~\ref{ch:technology}: Technology, Automation, Artificial Intelligence, and Productivity.}
The Solow growth model, Canada's productivity puzzle, the task model of
automation, which Canadian jobs are at risk, the economics of AI adoption,
and what productivity growth requires in terms of investment, competition
policy, and innovation incentives.

\textbf{Chapter~\ref{ch:climate}: Energy, Climate Change, and Canada's Green Transition.}
The carbon externality as a market failure, the Pigouvian tax rationale
for carbon pricing, Canada's carbon price schedule, the oil sands and
their fiscal importance, the electricity grid transition, the distributional
incidence of carbon pricing, and the economic challenge of achieving
net-zero by 2050 while managing a responsible transition for fossil-fuel
dependent workers and communities.

%==========================================================
\section*{Chapter Summary}
%==========================================================
\addcontentsline{toc}{section}{Chapter Summary}

\begin{itemize}
  \item Contemporary economic issues are actively unfolding, contested, and
        consequential problems that require both theoretical tools and
        empirical evidence to analyse rigorously.

  \item Canada is the world's ninth-largest economy with a GDP of
        approximately \$3.0~trillion (2024), a population of 41.5~million,
        and a GDP per capita of roughly \$72,000~CAD (\$53,000~USD
        PPP-adjusted) --- placing it above the OECD average but
        significantly below the United States.

  \item GDP is measured using the expenditure approach: $Y = C + I + G +
        (X-M)$. Real GDP adjusts nominal GDP for changes in the price level,
        measured by the GDP deflator. Real GDP growth is the standard measure
        of economic expansion or contraction.

  \item Canada's economy is predominantly service-based (77\% of GDP), but
        with a disproportionately important natural resource sector (oil and
        gas, mining, forestry). Economically, Canada is dominated by Ontario
        and Quebec (60\% of national GDP), with Alberta as the third-largest
        provincial economy due to oil sands.

  \item Canada is highly trade-dependent (trade openness $\approx$ 60.5\%
        of GDP) and extremely concentrated in the US market (75\% of
        merchandise exports go to the United States).

  \item Canada's major economic institutions include the Bank of Canada
        (monetary policy, inflation targeting), the Department of Finance
        (fiscal policy), Statistics Canada (data), CMHC (housing), and
        provincial governments (healthcare, education, labour, natural
        resources).

  \item Monetary policy works through the overnight rate and its transmission
        channels (credit cost, asset prices, exchange rate, expectations).
        Fiscal policy works through government spending and taxation.
        The two interact: loose fiscal policy may require tight monetary
        policy, and vice versa.

  \item Economic reasoning requires: models (deliberate simplifications),
        positive--normative distinctions (what is vs.\ what should be),
        opportunity cost thinking, marginal analysis, and careful attention
        to the difference between correlation and causation.

  \item Canada's macroeconomic history from 2015 to 2026 --- summarized in
        Table~\ref{tab:macro_indicators} --- encompasses a commodity price
        shock (2015--16), a historic pandemic shock (2020), a rapid fiscal
        and monetary policy response, and an inflationary surge and
        tightening cycle (2022--2024) that will shape Canada's economy
        for years.
\end{itemize}

%==========================================================
\section*{Key Terms}
%==========================================================
\addcontentsline{toc}{section}{Key Terms}

\begin{description}[style=nextline, leftmargin=3cm]
  \item[\term{Aggregate demand}] Total spending on domestically produced
        goods and services in the economy ($C + I + G + NX$).
  \item[\term{Bank of Canada}] Canada's central bank; sets the overnight
        rate to achieve the 2\% inflation target.
  \item[\term{Budget deficit}] The excess of government spending over
        government revenues in a given fiscal year.
  \item[\term{Comparative advantage}] The ability to produce a good at
        a lower opportunity cost than another producer.
  \item[\term{Contemporary economic issues}] Actively unfolding economic
        problems that are contested and consequential for social well-being.
  \item[\term{Correlation}] A statistical relationship in which two
        variables tend to move together; does not imply causation.
  \item[\term{Fiscal policy}] The use of government spending and taxation
        to influence macroeconomic outcomes.
  \item[\term{GDP deflator}] The ratio of nominal GDP to real GDP multiplied
        by 100; a broad measure of economy-wide price level changes.
  \item[\term{Gross Domestic Product (GDP)}] The total market value of all
        final goods and services produced within a country's borders in a
        given period.
  \item[\term{GDP per capita}] GDP divided by population; a rough measure
        of average material living standards.
  \item[\term{Inflation targeting}] A monetary policy framework in which
        a central bank commits to keeping inflation within a specified range.
  \item[\term{Marginal cost and benefit}] The additional cost or benefit
        of one more unit of an activity.
  \item[\term{Model}] A deliberate simplification of reality used to
        focus analysis on the mechanisms that matter for a specific question.
  \item[\term{Monetary policy}] Actions by a central bank to influence
        interest rates and credit conditions, aiming at price stability
        and sustainable growth.
  \item[\term{Net exports (NX)}] Exports minus imports; the trade balance
        component of the GDP expenditure identity.
  \item[\term{Nominal GDP}] GDP measured in current-year prices.
  \item[\term{Opportunity cost}] The value of the next-best alternative
        foregone when making a choice.
  \item[\term{Overnight rate}] The interest rate targeted by the Bank of
        Canada; the rate at which major financial institutions lend to each
        other overnight.
  \item[\term{Pareto efficiency}] An allocation in which no one can be made
        better off without making someone else worse off.
  \item[\term{Positive economics}] Description and explanation of economic
        phenomena without value judgements; testable against evidence.
  \item[\term{Normative economics}] Value judgements about what economic
        outcomes \textit{should} be.
  \item[\term{Real GDP}] GDP measured in the prices of a fixed base year;
        removes the effect of inflation from output comparisons.
  \item[\term{Trade openness}] $(X + M)/Y \times 100$; the ratio of total
        trade to GDP.
  \item[\term{Transfer payments}] Government payments to individuals that
        do not correspond to the production of new goods or services
        (e.g.\ EI, CPP, OAS).
\end{description}

%==========================================================
\section*{Review Questions}
%==========================================================
\addcontentsline{toc}{section}{Review Questions}

\begin{enumerate}

\item \textbf{(Concept)} Define GDP and explain what it measures and what
it does not measure. Give two examples of economically important activities
that GDP omits, and explain why each omission might matter for policy.

\item \textbf{(Concept)} Explain the difference between nominal GDP growth
and real GDP growth. If Canada's nominal GDP rose from \$2.72~trillion
in 2022 to \$2.88~trillion in 2023, but the GDP deflator rose from
110.3 to 115.2, what was Canada's real GDP growth rate in 2023?
Show all steps.

\item \textbf{(Institutions)} Explain the difference between the Bank of
Canada and the Department of Finance. Why does Canada have two separate
institutions managing monetary and fiscal policy rather than one?
What problem does this separation solve?

\item \textbf{(Concept)} Classify each of the following statements as
positive or normative, and explain your classification:
\begin{enumerate}[(a)]
  \item ``A 25\% US tariff on Canadian auto exports would cost approximately
        80,000 Canadian jobs.''
  \item ``Canada should implement a national childcare program to increase
        female labour force participation.''
  \item ``Canada's GDP per capita is approximately 37\% lower than that
        of the United States.''
  \item ``The Bank of Canada should have started raising interest rates
        in late 2021 rather than waiting until March 2022.''
\end{enumerate}

\item \textbf{(Application)} Use Table~\ref{tab:macro_indicators} to
answer the following:
\begin{enumerate}[(a)]
  \item In which year was Canada's unemployment rate highest?
        What event caused this?
  \item In which year was Canada's inflation rate highest?
        What was the Bank of Canada's policy rate in that year?
  \item What happened to the federal debt-to-GDP ratio between 2019
        and 2020? Express the change in percentage points and explain
        what drove it.
\end{enumerate}

\item \textbf{(Critical Thinking)} Canada's trade openness ratio is
approximately 60.5\%, while the United States' is approximately 27\%.
Using the concept of opportunity cost, explain why a high trade openness
ratio creates both benefits and risks. Is there a ``right'' level of
trade openness for a country like Canada?

\end{enumerate}

%==========================================================
\section*{Quantitative Problems}
%==========================================================
\addcontentsline{toc}{section}{Quantitative Problems}

\begin{enumerate}

\item \textbf{GDP Calculation.} You are given the following data for a
hypothetical Canadian province:
\begin{center}
\begin{tabular}{lr}
\toprule
Item & Value (\$B) \\
\midrule
Household consumption expenditure & 480 \\
Business investment in fixed capital & 120 \\
Change in business inventories & 8 \\
Provincial government spending on goods/services & 140 \\
Exports to other provinces and countries & 210 \\
Imports from other provinces and countries & 195 \\
\bottomrule
\end{tabular}
\end{center}
\begin{enumerate}[(a)]
  \item Calculate nominal GDP using the expenditure approach
        (Equation~\ref{eq:gdp_expenditure}). Show all steps.
  \item What is the trade balance for this province?
        What does its sign indicate?
  \item If the population is 4.8~million, calculate GDP per capita.
\end{enumerate}

\item \textbf{Real vs. Nominal GDP.} The following data describe a
simplified economy producing only two goods: wheat and lumber.
\begin{center}
\begin{tabular}{lcccc}
\toprule
& \multicolumn{2}{c}{\textbf{Year 1 (Base)}} & \multicolumn{2}{c}{\textbf{Year 2}} \\
\cmidrule(lr){2-3}\cmidrule(lr){4-5}
Good & Price & Quantity & Price & Quantity \\
\midrule
Wheat  & \$4.00/bushel & 500 & \$5.50/bushel & 520 \\
Lumber & \$8.00/board-ft & 300 & \$9.00/board-ft & 330 \\
\bottomrule
\end{tabular}
\end{center}
\begin{enumerate}[(a)]
  \item Calculate nominal GDP in Year 1 and Year 2.
  \item Calculate real GDP in Year 2 using Year 1 prices.
  \item Calculate the nominal GDP growth rate and the real GDP growth rate
        between Year 1 and Year 2.
  \item What is the GDP deflator in Year 2 (with Year 1 = 100)?
        What is the implied inflation rate between Year 1 and Year 2?
  \item If the nominal growth rate was 17.8\% but the real growth rate
        was only 6.5\%, what does the difference represent?
\end{enumerate}

\item \textbf{Trade Openness.} Using Table~\ref{tab:trading_partners}
and the following data:
\begin{itemize}
  \item Canada's total merchandise exports, 2024: \$897~billion CAD
  \item Canada's total merchandise imports, 2024: \$918~billion CAD
  \item Canada's nominal GDP, 2024: \$3,001~billion CAD
\end{itemize}
\begin{enumerate}[(a)]
  \item Calculate Canada's trade openness ratio
        (Equation~\ref{eq:trade_openness}).
  \item Calculate Canada's merchandise trade balance.
  \item What share of Canada's total exports goes to the United States?
        (Use Table~\ref{tab:trading_partners}.)
  \item If the United States were to impose a 25\% tariff on all
        Canadian goods, and this caused Canada's US-bound exports to
        fall by 30\%, what would be the new trade balance?
        (Assume all other trade is unchanged.)
\end{enumerate}

\item \textbf{GDP per Capita Growth.} The relationship between GDP per
capita growth, GDP growth, and population growth is approximately:
\[
  g_{\text{per capita}} \approx g_{\text{GDP}} - g_{\text{population}}
\]
\begin{enumerate}[(a)]
  \item In 2023, Canada's real GDP grew by 1.2\% and its population
        grew by approximately 3.2\% (driven largely by immigration).
        What was the approximate growth rate of real GDP \textit{per capita}?
        What does this imply about average living standards?
  \item If Canada's potential GDP growth rate is 2.0\% and population
        growth is 3.0\%, what is the implied trend in GDP per capita?
        Is this sustainable?
  \item Suggest two economic policies that could raise Canada's GDP
        growth rate so that it exceeds its population growth rate.
        For each policy, identify a potential trade-off.
\end{enumerate}

\item \textbf{Opportunity Cost.} The federal government has \$20~billion
in additional fiscal capacity. Three competing proposals are on the table:
\begin{itemize}
  \item Proposal A: Invest in new public housing construction.
  \item Proposal B: Cut the corporate income tax rate by 2 percentage points.
  \item Proposal C: Fund a national pharmacare program for uninsured drugs.
\end{itemize}
\begin{enumerate}[(a)]
  \item For each proposal, identify one major benefit and one major cost
        (the opportunity cost of choosing that proposal over the others).
  \item Which economic concept (efficiency, equity, or both) is most
        relevant for evaluating each proposal?
  \item Is the choice among these three proposals a positive economic
        question, a normative question, or both? Explain.
\end{enumerate}

\item \textbf{Data Exercise.} Visit the Statistics Canada website
(\url{www.statcan.gc.ca}) and locate Table 36-10-0104-01 (GDP by
expenditure, seasonally adjusted, at annual rates).
\begin{enumerate}[(a)]
  \item Report the real GDP growth rate (annualized) for the most recent
        available quarter.
  \item Break down GDP by component (C, I, G, NX) for the most recent quarter.
        Which component is largest? Which is most volatile over the past
        four quarters?
  \item Calculate the percentage change in each component over the past
        four quarters.
  \item Based on the current composition of growth, assess whether Canada's
        economic expansion is likely to be sustainable. Which components
        would you expect to change if interest rates rise further?
\end{enumerate}

\end{enumerate}

%==========================================================
\section*{Essay Questions}
%==========================================================
\addcontentsline{toc}{section}{Essay Questions}

\begin{enumerate}

\item \textbf{(500--700 words)} Using data from Table~\ref{tab:macro_indicators},
write an essay describing the macroeconomic history of Canada between 2019
and 2026. Your essay should identify the major shocks, explain the policy
responses, assess their effectiveness, and discuss the economic legacy of
this period for Canada going forward. Use at least three specific data
points from the table to support your argument.

\item \textbf{(400--600 words)} Canada runs an enormous merchandise trade
surplus with the United States but trade deficits with most other countries
(see Table~\ref{tab:trading_partners}). Using the concepts of comparative
advantage and opportunity cost, explain what this pattern tells us about
Canada's productive specialization. Discuss one risk and one benefit of
Canada's trade concentration in the US market.

\item \textbf{(400--500 words)} The opening vignette of this chapter
describes four people --- Diane, Marcus, the sawmill manager, and Priya
--- each facing an economic challenge shaped by macroeconomic forces.
Choose two of these four people. For each, identify the specific macroeconomic
force most relevant to their situation, explain the economic mechanism
at work, and identify one policy from Table~\ref{tab:institutions}
(or one of the institutions described in Section~\ref{sec:institutions})
that is most directly relevant to addressing their challenge.

\end{enumerate}

%==========================================================
\section*{Discussion Questions}
%==========================================================
\addcontentsline{toc}{section}{Discussion Questions}

\begin{enumerate}

\item Figure~\ref{fig:gdp_per_capita} shows that Canada's GDP per capita
is approximately 37\% lower than that of the United States, despite
Canada having a comparable education system, similar institutions, and
vast natural resource wealth. What do you think explains this gap?
Is it a problem that Canada should try to close? What would closing
it require?

\item Should the Bank of Canada be fully independent from the elected
federal government? What would you gain, and what would you lose,
by giving Parliament more direct control over monetary policy?
Draw on the Policy Debate box in Section~\ref{sec:policy_toolkit}
in your answer.

\item Canada's regional economic inequality --- illustrated in
Figure~\ref{fig:gdp_by_province} --- means that national economic
policies have unequal regional consequences. Is this a problem
unique to Canada, or a challenge faced by all federations?
How should the federal government balance national economic objectives
with regional equity?

\item The opening of this chapter describes four people whose lives
are shaped by macroeconomic forces they cannot individually control.
Do you think most Canadians understand how economics affects their
daily lives? What are the consequences --- political, social, and
economic --- of widespread economic illiteracy?

\end{enumerate}

%==========================================================
%  CHAPTER 2: Inflation, Monetary Policy, and the Cost of Living
%  ECON 204 — Contemporary Economic Issues: A Canadian Perspective
%==========================================================

\chapter{Inflation, Monetary Policy, and the Cost of Living}
\label{ch:inflation}

%----------------------------------------------------------
%  OPENING VIGNETTE
%----------------------------------------------------------

\noindent In January 2021, Fatima filled her grocery cart in Saskatoon and
paid \$183. By January 2022, the same cart cost \$206. By June 2022, it
cost \$228. In eighteen months, her weekly grocery bill had risen by
\$45 --- a 25\% increase on food alone. Gasoline at the pump had gone from
\$1.04 per litre to \$1.85. Her landlord had issued a rent increase notice
citing rising property taxes and maintenance costs. Her variable-rate
mortgage payment had climbed \$390 per month since March 2022, tracking the
Bank of Canada's overnight rate almost in lockstep.

Fatima had not changed her lifestyle. She had not moved to a more expensive
neighbourhood or started eating at restaurants. She had the same job, the
same apartment, and the same spending habits. But the purchasing power of
her income was shrinking, month by month, in a way she could not control
and could barely predict.

Statistics Canada would later record that Canada's Consumer Price Index rose
8.1\% in June 2022 --- the highest reading in 40 years. That number meant
something very specific: an average Canadian household needed \$108.10 in
June 2022 to buy what \$100 bought in June 2021. For Fatima, whose budget
had less slack than the average, the real erosion was larger.

What caused inflation this severe? Why did it emerge so suddenly after three
decades of low, stable prices? How does the Bank of Canada fight it, and
at what cost? And who bears the heaviest burden when prices rise faster
than wages? These are the questions this chapter answers.

%----------------------------------------------------------
%  LEARNING OBJECTIVES
%----------------------------------------------------------

\begin{learningobjectives}
  \item Define inflation, calculate the CPI inflation rate from price
        data, and identify the four main biases in CPI measurement.
  \item Distinguish headline CPI from core inflation measures and explain
        why the Bank of Canada monitors both.
  \item Apply the Fisher equation to calculate real interest rates from
        nominal rates and inflation, and derive real wage growth from
        nominal wage and CPI data.
  \item Use the aggregate demand--aggregate supply framework to explain
        demand-pull and cost-push inflation, and apply both to Canada's
        2021--2023 episode.
  \item Describe the short-run and long-run Phillips Curve, interpret
        Canadian data on the inflation--unemployment relationship, and
        define the sacrifice ratio.
  \item Explain how the Bank of Canada's inflation-targeting framework
        works, trace the channels of monetary policy transmission, and
        evaluate the 2022--2024 rate-hike cycle.
  \item Analyse the distributional consequences of inflation across
        income groups, between debtors and creditors, and between
        workers with indexed and non-indexed compensation.
\end{learningobjectives}

%==========================================================
\section{Why Inflation Matters}
%==========================================================

\term{Inflation} is a sustained increase in the general price level of goods
and services in an economy over time. Three words in that definition deserve
emphasis. \textbf{Sustained}: a single price spike --- gasoline surging one
month and falling the next --- is not inflation. \textbf{General}: inflation
refers to the average price level, not any specific good. A single price
rising while others fall is a \term{relative price} change, not inflation.
\textbf{Over time}: inflation is a dynamic phenomenon measured as a rate
of change, typically expressed as a percentage per year.

The opposite of inflation --- a sustained \textit{decrease} in the general
price level --- is \term{deflation}. Deflation may sound attractive
(prices falling means money buys more), but it is one of the most dangerous
macroeconomic conditions an economy can experience. When prices are expected
to fall, households postpone purchases (why buy today what will be cheaper
tomorrow?), firms delay investment, profits shrink, layoffs rise, and the
economy can fall into a self-reinforcing downward spiral. Japan's ``Lost
Decade'' of the 1990s is the modern cautionary example.

Low, stable, \textit{positive} inflation --- in Canada's case, the
2\% target --- is the policy goal. Why 2\% rather than 0\%? Three reasons:

\begin{enumerate}
  \item \textbf{Measurement bias:} The CPI tends to overstate true inflation
        by approximately 0.3--0.6 percentage points because it cannot
        fully capture quality improvements, new goods, and consumer
        substitution toward cheaper alternatives. A 2\% measured target
        therefore implies roughly zero true inflation.
  \item \textbf{Policy buffer:} If the equilibrium real interest rate is
        approximately 2\% and inflation is targeted at 2\%, the Bank of
        Canada can set a nominal rate of approximately 4\% at equilibrium.
        This leaves 400 basis points of room to cut rates before hitting
        the zero lower bound in a recession. A 0\% inflation target would
        leave only 200 basis points of room --- dangerously thin.
  \item \textbf{Wage adjustment:} Inflation allows real wages to fall
        without requiring nominal pay cuts, which workers strongly resist.
        In a world of 2\% inflation, a firm that needs to reduce labour
        costs can simply freeze wages --- real wages fall while nominal
        wages stay constant.
\end{enumerate}

\begin{thinkbox}
If 2\% inflation is the goal, should the Bank of Canada have started
raising rates in late 2021 when inflation was already at 4.7\%, rather than
waiting until March 2022 when it had reached 5.7\%? This is one of the
most debated questions in recent Canadian monetary policy history.
Your answer requires a theory of how quickly monetary policy works
(the transmission lags), an assessment of whether the inflation was
supply-side or demand-side, and a judgement about the risk of
over-tightening into recession. We return to this question in the
Policy Debate at the end of this chapter.
\end{thinkbox}

%==========================================================
\section{Measuring Inflation: The Consumer Price Index}
\label{sec:cpi_measurement}
%==========================================================

%----------------------------------------------------------
\subsection{The CPI Basket Approach}
%----------------------------------------------------------

Canada's primary measure of inflation is the \term{Consumer Price Index (CPI)},
produced monthly by Statistics Canada. The CPI measures the change in the
total cost of a fixed \term{representative basket} of goods and services
purchased by a typical Canadian urban household.

\begin{defbox}{Consumer Price Index (CPI)}
The \textbf{Consumer Price Index} measures the average change over time
in the prices paid by urban consumers for a fixed basket of goods and
services.

\textit{Economic intuition:} The CPI asks: ``How much does it cost today,
relative to a base period, to maintain a fixed standard of living?''
If the basket cost \$1{,}000 in the base period and costs \$1{,}081 today,
the CPI is 108.1, implying 8.1\% cumulative inflation since the base.

\textit{Formal definition:} Statistics Canada periodically determines the
basket weights from the Survey of Household Spending. Prices are collected
from approximately 30{,}000 retail outlets, service providers, and
landlords across Canada each month.

\textit{Canadian note:} Statistics Canada publishes the CPI monthly in
Table 18-10-0004-01. The June 2022 CPI reading was the highest year-over-year
change since September 1983, when Canada was in the final stages of
fighting the 1970s inflation surge.
\end{defbox}

Construction of the CPI proceeds in three steps:

\begin{enumerate}
  \item \textbf{Determine basket weights.} Statistics Canada's Survey of
        Household Spending identifies what Canadian urban households buy
        and in what proportions. These proportions are the
        \term{expenditure weights} $w_i$ for each category $i$.
  \item \textbf{Collect prices.} Each month, approximately 100{,}000 price
        observations are gathered from retailers, landlords, and service
        providers across Canada.
  \item \textbf{Construct the index.} The CPI compares current-period prices
        to base-period prices, weighted by expenditure shares.
\end{enumerate}

The formal CPI formula is a Laspeyres price index:

\begin{equation}
  \text{CPI}_t = \frac{\displaystyle\sum_{i} w_i \cdot P_{i,t}}
                      {\displaystyle\sum_{i} w_i \cdot P_{i,0}}
                 \times 100
  \label{eq:cpi_formula}
\end{equation}

\noindent where $P_{i,t}$ is the price of good $i$ in period $t$,
$P_{i,0}$ is the price of good $i$ in the base period, and $w_i$ is the
expenditure weight of good $i$.

The \term{inflation rate} between two consecutive periods is:

\begin{equation}
  \pi_t = \frac{\text{CPI}_t - \text{CPI}_{t-1}}{\text{CPI}_{t-1}} \times 100
  \label{eq:inflation_rate}
\end{equation}

\noindent\textbf{Worked Example 2.1 --- Calculating the CPI from a Simple Basket}

Suppose a household's basket consists of three goods:

\begin{center}
\renewcommand{\arraystretch}{1.3}
\begin{tabular}{lrrrr}
\toprule
Good & Weight $w_i$ & Base Price $P_{i,0}$ & Current Price $P_{i,t}$
     & Price Relative $P_{i,t}/P_{i,0}$ \\
\midrule
Food       & 0.16 & \$120 & \$131.40 & 1.095 \\
Shelter    & 0.30 & \$900 & \$963.00 & 1.070 \\
Transport  & 0.20 & \$300 & \$342.00 & 1.140 \\
\midrule
\multicolumn{5}{l}{\textit{(Other categories omitted for simplicity)}} \\
\bottomrule
\end{tabular}
\end{center}

\noindent\textit{Step 1:} Calculate the weighted sum of price relatives:
\[
  \text{CPI}_t = (0.16 \times 1.095) + (0.30 \times 1.070)
                + (0.20 \times 1.140) + \ldots
\]
\[
  = 0.1752 + 0.3210 + 0.2280 + \ldots = 0.7242 \text{ (for these three)}
\]

\noindent\textit{Full CPI} (with all categories at assumed price relatives):
Scaling up with remaining categories at an assumed average relative of 1.060:
\[
  \text{CPI}_t \approx 107.5 \implies \pi_t \approx 7.5\%
\]

\noindent\textit{Interpretation:} The basket now costs 7.5\% more than
in the base period. Because shelter (the largest weight) rose only 7.0\%
while transport (a smaller weight) rose 14.0\%, the shelter-weighted
average moderates what would otherwise be a higher measured rate.

\begin{table}[H]
\centering
\caption{Canada CPI Inflation by Component, 2020--2024 (Annual Average, \%)}
\label{tab:cpi_components}
\renewcommand{\arraystretch}{1.30}
\begin{tabular}{lrrrrr}
\toprule
\textbf{CPI Component} & \textbf{2020} & \textbf{2021} & \textbf{2022}
                       & \textbf{2023} & \textbf{2024} \\
\midrule
All-items (headline)  &  0.7 &  3.4 &  6.8 & 3.9 & 2.6 \\
\midrule
Food                  &  2.4 &  3.9 &  9.7 & 5.9 & 1.5 \\
\quad Food purchased from stores & 2.8 & 3.1 & 10.3 & 6.5 & 1.2 \\
\quad Food from restaurants & 1.5 & 5.0 & 8.3 & 5.1 & 2.1 \\
Shelter               &  2.4 &  4.6 &  6.9 & 6.0 & 5.7 \\
\quad Mortgage interest costs & 0.8 & 1.3 & 11.5 & 28.5 & 19.4 \\
\quad Rent             &  3.0 &  1.4 &  4.5 &  8.1 &  8.9 \\
Transportation        & $-$3.6 &  8.3 & 10.6 & 2.4 & 0.8 \\
\quad Gasoline        &$-$14.3 & 32.2 & 31.7 &$-$4.2 &$-$2.1 \\
Energy (total)        & $-$8.5 & 19.5 & 28.5 &$-$8.5 &$-$2.7 \\
Services              &  1.4 &  2.5 &  5.5 & 4.7 & 4.0 \\
Goods (excl.\ food/energy) & 0.3 & 4.3 & 7.2 & 2.7 & 0.5 \\
\bottomrule
\end{tabular}
\source{Statistics Canada, Table 18-10-0004-01, CPI by component. Values
are annual averages; 2024 values are estimates consistent with published data.}
\end{table}

\begin{figure}[H]
\centering
\begin{tikzpicture}
\begin{axis}[
  ybar,
  bar width=0.55cm,
  width=13cm, height=6.5cm,
  symbolic x coords={{Shelter},{Transport},{Food},{Rec./Educ.},
                     {Household},{Health},{Clothing},{Alc./Tob.}},
  xtick=data,
  x tick label style={rotate=35, anchor=east, font=\small},
  ylabel={Weight in CPI Basket (\%)},
  ylabel style={font=\small},
  ymin=0, ymax=34,
  ytick={0,5,10,15,20,25,30},
  yticklabel style={font=\small},
  ymajorgrids=true,
  grid style={dashed, gray!25},
  enlarge x limits=0.08,
  nodes near coords,
  nodes near coords style={font=\footnotesize},
  every axis plot/.append style={fill=unbcgreen!70, draw=unbcgreen!90}
]
\addplot coordinates {
  ({Shelter},    30.0)
  ({Transport},  20.0)
  ({Food},       16.0)
  ({Rec./Educ.}, 11.0)
  ({Household},  10.0)
  ({Health},      5.0)
  ({Clothing},    5.0)
  ({Alc./Tob.},   3.0)
};
\end{axis}
\end{tikzpicture}
\caption{Weights in Canada's CPI Basket, 2022. Shelter and transportation
together account for 50\% of the CPI basket. This structural weighting
means that even moderate increases in rent or mortgage costs --- both
classified under Shelter --- move the headline inflation rate substantially.
The dominance of Shelter also means that the Bank of Canada's rate hikes,
which raised mortgage costs sharply, paradoxically \textit{increased}
one component of the CPI (mortgage interest costs) even as they reduced
the underlying demand driving inflation.}
\label{fig:cpi_basket}
\source{Statistics Canada, CPI Basket Weights, 2022 reference basket.}
\end{figure}

%----------------------------------------------------------
\subsection{Biases in CPI Measurement}
%----------------------------------------------------------

The CPI is a Laspeyres index: it holds the basket fixed at its base-period
composition. This creates systematic upward biases in measured inflation:

\begin{itemize}
  \item \textbf{Substitution bias:} When the price of good A rises relative
        to good B, consumers substitute toward B. The fixed CPI basket
        does not reflect this --- it still assumes the old, pre-increase
        quantities of A are purchased. Result: the CPI overstates the
        true cost-of-living increase.

  \item \textbf{Quality bias:} If a new car costs \$5{,}000 more than last
        year's model but is also significantly better (safer, more
        fuel-efficient, with better technology), not all of the \$5{,}000
        increase is inflation. Some of it is more quality per dollar.
        Statistics Canada adjusts for some quality changes, but imperfectly.

  \item \textbf{New goods bias:} When a completely new product (say,
        smartphones in 2007) enters the market, it is not immediately
        included in the basket. The large initial price declines as new
        products diffuse through the market are missed.

  \item \textbf{Outlet substitution bias:} When consumers shift purchases
        from traditional retailers to lower-priced online sellers, the
        CPI may not fully capture the effective price reduction.
\end{itemize}

Studies by Statistics Canada and the Bank of Canada suggest the CPI
overstates true inflation by approximately 0.3--0.5 percentage points
per year on average. This is one reason the Bank of Canada's 2\% inflation
target is set above zero rather than at zero.

%----------------------------------------------------------
\subsection{Canada's Annual Inflation Record, 2010--2026}
%----------------------------------------------------------

\begin{figure}[H]
\centering
\begin{tikzpicture}
\begin{axis}[
  width=14cm, height=7cm,
  xlabel={Year},
  ylabel={CPI Inflation, Year-over-Year (\%)},
  xlabel style={font=\small},
  ylabel style={font=\small},
  xmin=2009.5, xmax=2026.5,
  ymin=-1, ymax=10,
  xtick={2010,2012,2014,2016,2018,2020,2022,2024,2026},
  xticklabel style={font=\small},
  ytick={-1,0,1,2,3,4,5,6,7,8,9,10},
  yticklabel style={font=\small},
  ymajorgrids=true,
  grid style={dashed, gray!20},
  axis line style={color=gray!60},
  tick style={color=gray!60},
  clip=false
]
% Control range band 1--3%
\fill[unbcgreen!12] (axis cs:2009.5,1) rectangle (axis cs:2026.5,3);
% 2% target line
\addplot[unbcgreen!60, dashed, thick] coordinates {(2009.5,2.0)(2026.5,2.0)}
  node[right, font=\footnotesize, xshift=2pt, color=unbcgreen!80] {2\% target};
% 1% lower bound
\addplot[unbcgreen!40, dotted] coordinates {(2009.5,1.0)(2026.5,1.0)};
% 3% upper bound
\addplot[unbcgreen!40, dotted] coordinates {(2009.5,3.0)(2026.5,3.0)};
% CPI data line
\addplot[policyblue, thick, mark=*, mark size=1.5pt] coordinates {
  (2010, 1.8)(2011, 2.9)(2012, 1.5)(2013, 0.9)(2014, 2.0)
  (2015, 1.1)(2016, 1.4)(2017, 1.6)(2018, 2.3)(2019, 1.9)
  (2020, 0.7)(2021, 3.4)(2022, 6.8)(2023, 3.9)(2024, 2.6)
  (2025, 2.1)(2026, 2.0)
};
% Annotations
\node[font=\footnotesize, align=center, text=policyblue, fill=white,
      inner sep=1pt] at (axis cs:2022,7.5)
  {June 2022 peak: 8.1\%\\(annual avg: 6.8\%)};
\draw[->, policyblue, thin] (axis cs:2022,7.3) -- (axis cs:2022,6.9);
\node[font=\footnotesize, color=gray] at (axis cs:2020,0.2) {COVID};
\node[font=\footnotesize, align=center, color=unbcgreen!70] at (axis cs:2017,0.2)
  {Control range\\(1--3\%)};
\end{axis}
\end{tikzpicture}
\caption{Canada Annual CPI Inflation Rate, 2010--2026. The shaded band
shows the Bank of Canada's 1--3\% control range; the dashed line marks the
2\% midpoint target. For most of 2012--2021, inflation remained within or
near the band, reflecting credible monetary policy. The 2020 drop to 0.7\%
(pandemic demand collapse) and the subsequent surge to a peak of 8.1\%
in June 2022 (annual average 6.8\%) represent the most dramatic deviation
from target in the Bank's 30-year history of inflation targeting. By 2024,
inflation had returned to the upper portion of the target range.}
\label{fig:cpi_annual}
\source{Statistics Canada, Table 18-10-0004-01; Bank of Canada. Annual averages;
2025--2026 are estimates consistent with Bank of Canada projections.}
\end{figure}

\begin{keytakeaway}
The CPI is the standard measure of inflation but contains systematic
upward biases of approximately 0.3--0.5 percentage points per year.
The Bank of Canada targets 2\% measured CPI inflation, which is
roughly consistent with near-zero \textit{true} inflation once biases
are accounted for. Understanding CPI biases is essential for interpreting
economic data and evaluating policy.
\end{keytakeaway}

%==========================================================
\section{Core Inflation: Seeing Through the Noise}
\label{sec:core_inflation}
%==========================================================

Headline CPI measures the price level experienced by the average household.
But for \textit{monetary policy}, headline CPI is a problematic guide. The
reason is that a central bank can only affect demand conditions --- it cannot
directly lower the price of crude oil, ease a shipping container shortage,
or replace a failed harvest. When inflation is driven by temporary supply
shocks (a spike in gasoline prices, a drought-induced food price jump), a
central bank that responds aggressively risks \term{policy error}: tightening
into a supply-driven recession.

\term{Core inflation} measures remove the most volatile components from
the headline CPI to reveal the underlying, persistent inflation trend that
monetary policy can actually address. The Bank of Canada monitors three
official core measures:

\begin{defbox}{Bank of Canada Core Inflation Measures}
\textbf{CPI-trim:} Removes the tails of the monthly price change
distribution --- the highest and lowest 20\% of components by price
change magnitude --- and computes the weighted average of the remaining
60\%.

\textbf{CPI-median:} The price change at the 50th percentile of the
weighted distribution of component price changes. Exactly half of the
CPI basket (by weight) is changing faster than this rate, and half
slower.

\textbf{CPI-common:} Extracts the common variation in price changes
across CPI components using statistical factor analysis. It captures
inflation that is economy-wide rather than sector-specific.

\textit{Why three measures?} Each measure has different statistical
properties. By monitoring all three and looking for consensus, the Bank
reduces the risk of any single measure sending a misleading signal.
When all three core measures are elevated, the Bank is confident that
inflation is broad-based, persistent, and amenable to a policy response.
\end{defbox}

\begin{figure}[H]
\centering
\begin{tikzpicture}
\begin{axis}[
  width=14cm, height=7cm,
  xlabel={Year},
  ylabel={Inflation Rate (\%, Year-over-Year)},
  xlabel style={font=\small},
  ylabel style={font=\small},
  xmin=2017.8, xmax=2026.5,
  ymin=0, ymax=9,
  xtick={2018,2019,2020,2021,2022,2023,2024,2025,2026},
  xticklabel style={font=\small},
  ytick={0,1,2,3,4,5,6,7,8,9},
  yticklabel style={font=\small},
  ymajorgrids=true,
  grid style={dashed, gray!20},
  legend pos=north west,
  legend style={font=\small, fill=white, fill opacity=0.9},
  clip=false
]
% Control range band
\fill[unbcgreen!10] (axis cs:2017.8,1) rectangle (axis cs:2026.5,3);
\addplot[unbcgreen!50, dashed] coordinates {(2017.8,2)(2026.5,2)};
% Headline CPI
\addplot[gray!70, thick, mark=none] coordinates {
  (2018,2.3)(2019,1.9)(2020,0.7)(2021,3.4)(2022,6.8)(2023,3.9)(2024,2.6)(2025,2.1)(2026,2.0)
};
\addlegendentry{Headline CPI}
% CPI-trim (smoother, peaks lower than headline)
\addplot[policyblue, thick, mark=none] coordinates {
  (2018,2.1)(2019,2.1)(2020,1.5)(2021,2.9)(2022,5.5)(2023,3.7)(2024,2.7)(2025,2.2)(2026,2.0)
};
\addlegendentry{CPI-trim}
% CPI-median
\addplot[casered, thick, mark=none, dashed] coordinates {
  (2018,2.0)(2019,2.2)(2020,1.9)(2021,2.7)(2022,5.0)(2023,3.5)(2024,2.6)(2025,2.1)(2026,2.0)
};
\addlegendentry{CPI-median}
% CPI-common
\addplot[dataorange, thick, mark=none, dotted] coordinates {
  (2018,1.9)(2019,1.9)(2020,1.4)(2021,2.2)(2022,4.6)(2023,3.1)(2024,2.5)(2025,2.1)(2026,2.0)
};
\addlegendentry{CPI-common}
% Annotations
\node[font=\footnotesize, text=gray!80] at (axis cs:2022.3,7.2)
  {Headline 8.1\% peak};
\node[font=\footnotesize, text=policyblue, align=center] at (axis cs:2022.8,6.0)
  {Core peaks:\\5--6\%};
\end{axis}
\end{tikzpicture}
\caption{Headline CPI vs.\ Bank of Canada Core Inflation Measures, 2018--2026.
The core measures (CPI-trim, CPI-median, CPI-common) peaked approximately
2--3 percentage points below the 8.1\% headline peak, reflecting the
contribution of volatile energy and some food prices to the headline figure.
Crucially, all three core measures rose substantially above the 3\% upper
bound of the control range in 2022 --- confirming that the inflation surge
was broad-based, not merely a transitory spike in a few volatile categories.
This broad-based elevation justified the Bank of Canada's aggressive tightening.}
\label{fig:core_inflation}
\source{Bank of Canada, Measures of Core Inflation database; Statistics Canada,
Table 18-10-0004-01. Values for 2025--2026 are estimates.}
\end{figure}

A critical insight from Figure~\ref{fig:core_inflation}: the gap between
headline and core inflation \textit{narrowed} in 2023--2024. As energy prices
fell back from their 2022 peaks, headline inflation declined sharply. But
core measures remained elevated well into 2024 because \term{services
inflation} --- wages, rents, and labour-intensive services --- proved
``sticky'' downward. This stickiness reflects the second-round effects of
inflation: when workers see prices rise, they demand higher wages; higher
wages increase firms' costs, pushing up prices further; and the cycle
reinforces itself. Breaking this cycle is the primary rationale for the
Bank of Canada's rate-hike strategy.

%==========================================================
\section{Real vs.\ Nominal Interest Rates: The Fisher Equation}
\label{sec:fisher_equation}
%==========================================================

When a bank offers you a savings account paying 5\% interest, is that a good
deal? The answer depends on inflation. If prices are rising at 3\% per year,
your 5\% nominal return only delivers 2\% more purchasing power --- the
real return is just 2\%. If inflation is running at 6\%, your ``5\%'' nominal
return is actually a loss of 1\% in purchasing power.

This insight is captured by the \term{Fisher equation}, one of the most
important relationships in monetary economics:

\begin{equation}
  r \approx i - \pi
  \label{eq:fisher_simple}
\end{equation}

\noindent where $r$ is the \term{real interest rate}, $i$ is the
\term{nominal interest rate}, and $\pi$ is the \term{inflation rate}.

The exact version is:

\begin{equation}
  (1 + r) = \frac{(1 + i)}{(1 + \pi)}
  \quad \Longleftrightarrow \quad
  r = \frac{1 + i}{1 + \pi} - 1
  \label{eq:fisher_exact}
\end{equation}

For the values typically encountered in Canada, the approximation
$r \approx i - \pi$ (Equation~\ref{eq:fisher_simple}) is sufficiently
accurate for analysis.

\begin{defbox}{Real vs.\ Nominal Interest Rates}
The \textbf{nominal interest rate} $i$ is the rate stated in the
loan contract or savings agreement --- the number on the Bank of Canada
website.

The \textbf{real interest rate} $r = i - \pi^e$ is the nominal rate
adjusted for (expected) inflation. It represents the true cost of
borrowing or the true return to saving in terms of purchasing power.

When making economic decisions, it is the \textit{real} interest rate
that matters. A firm deciding whether to invest in new equipment
compares the real borrowing cost against the real rate of return on
the equipment. A household deciding whether to take out a mortgage
compares the real cost of the loan against the real appreciation of
the home.

\textit{Ex ante vs.\ ex post:} Before the fact (\textit{ex ante}),
borrowers and lenders agree on a nominal rate based on their expectation
of future inflation $\pi^e$. After the fact (\textit{ex post}), the real
rate reflects actual inflation $\pi$. Unexpected inflation benefits
debtors (who repay in cheapened dollars) and hurts creditors (who receive
less purchasing power than expected).
\end{defbox}

\begin{table}[H]
\centering
\caption{Real Interest Rate Calculation, Canada: 2018--2026}
\label{tab:real_rates}
\renewcommand{\arraystretch}{1.30}
\begin{tabular}{lrrrrr}
\toprule
\textbf{Year} & \textbf{Nominal Rate} & \textbf{CPI Inflation} & \textbf{Real Rate}
              & \textbf{Real Rate} & \textbf{Policy} \\
              & $i$ (Bank Rate, \%) & $\pi$ (\%) & $r \approx i - \pi$ (\%)
              & Interpretation & Context \\
\midrule
2018 & 1.75 & 2.3 & $-$0.6 & Mildly stimulative & Late expansion \\
2019 & 1.75 & 1.9 & $-$0.2 & Roughly neutral    & Pre-COVID \\
2020 & 0.25 & 0.7 & $-$0.5 & Stimulative        & COVID emergency \\
2021 & 0.25 & 3.4 & $-$3.2 & Highly stimulative & Post-COVID stimulus \\
2022 & 4.25\rlap{$^*$} & 6.8 & $-$2.6\rlap{$^*$} & Still negative & Tightening cycle \\
2023 & 5.00 & 3.9 & +1.1   & Mildly restrictive & Tightening complete \\
2024 & 3.75\rlap{$^\dagger$} & 2.6 & +1.2 & Mildly restrictive & Easing begins \\
2025 & 3.00 & 2.1 & +0.9   & Near neutral       & Easing continues \\
2026 & 2.75 & 2.0 & +0.8   & Near neutral       & Soft landing \\
\bottomrule
\end{tabular}
\vspace{4pt}
\noindent\footnotesize{$^*$End-of-year Bank Rate; mid-year 2022 average rate was
approximately 2.0\%, making the real rate even more negative at $-4.8\%$.
$^\dagger$December 2024 rate after four cuts in the second half of 2024.}
\source{Bank of Canada key interest rates; Statistics Canada Table 18-10-0004-01;
author calculations.}
\end{table}

\noindent\textbf{The Key Lesson from Table~\ref{tab:real_rates}:}
In 2021 and through most of 2022, the real interest rate was deeply
\textit{negative} --- borrowers were effectively being paid to borrow in
real terms. At a mid-2022 Bank Rate of $i = 2.5\%$ and inflation of
$\pi = 7\%$, the real rate was approximately $r = -4.5\%$. This meant
that anyone borrowing at the prime rate to purchase real assets (housing,
stocks, commodities) was receiving a massive implicit subsidy from
inflation. Negative real rates were a powerful engine of the demand boom
that itself contributed to inflation.

%----------------------------------------------------------
\subsection{Real Wages: The Fisher Equation Applied to Labour}
%----------------------------------------------------------

The Fisher relationship applies equally to wages. The \term{nominal wage}
$W$ is the dollar amount a worker earns. The \term{real wage} is the nominal
wage deflated by the price level:

\begin{equation}
  w_{\text{real},t} = \frac{W_t}{\text{CPI}_t / 100}
  \label{eq:real_wage}
\end{equation}

The growth rate of real wages is approximately:

\begin{equation}
  g_{w,\text{real}} \approx g_{W,\text{nominal}} - \pi
  \label{eq:real_wage_growth}
\end{equation}

\noindent\textbf{Worked Example 2.2 --- Real Wage Growth, Canada 2021--2023}

Average hourly wages in Canada grew from \$28.57 in January 2021 to
\$32.36 in January 2023.
\begin{itemize}
  \item Nominal wage growth: $(32.36 - 28.57) / 28.57 \times 100 = 13.3\%$
  \item CPI in January 2021: 138.8; January 2023: 155.6
  \item CPI change: $(155.6 - 138.8)/138.8 \times 100 = 12.1\%$
  \item Real wage growth: $13.3\% - 12.1\% = +1.2\%$ over two years
\end{itemize}

Despite nominal wage gains that looked impressive, real wages grew by
only 0.6\% per year on average --- barely keeping pace with productivity.
For workers in the bottom two income quintiles, where wage gains were
smaller and inflation in food and shelter (which loom larger in their
budgets) was higher, real wages actually fell.

%==========================================================
\section{The Aggregate Demand--Aggregate Supply Framework}
\label{sec:adas}
%==========================================================

The \term{Aggregate Demand--Aggregate Supply (AD-AS) model} is the
workhorse framework for analysing macroeconomic fluctuations, including
inflation. It explains how the overall price level and real output are
jointly determined in the economy, and why they change over time.

%----------------------------------------------------------
\subsection{Aggregate Demand}
%----------------------------------------------------------

The \term{Aggregate Demand (AD) curve} shows the combinations of price
level and real output at which the goods market and money market are
simultaneously in equilibrium. It slopes downward for three reasons:

\begin{itemize}
  \item \textbf{Real balance effect:} A higher price level reduces the
        real value of household wealth (including cash and bonds held at
        face value), reducing consumption spending.
  \item \textbf{Interest rate effect:} A higher price level increases
        the demand for money (people need more nominal money to buy the
        same real quantities), pushing up interest rates and reducing
        investment.
  \item \textbf{Exchange rate effect:} Higher domestic prices make
        Canadian goods less competitive internationally, reducing exports
        and increasing imports --- reducing net exports $NX$.
\end{itemize}

A simplified aggregate demand relationship can be written:

\begin{equation}
  Y = \alpha_0 - \beta \cdot r + \gamma \cdot G - \delta \cdot P
  \label{eq:ad_simplified}
\end{equation}

\noindent where $Y$ is real output, $r$ is the real interest rate,
$G$ is government spending, and $P$ is the price level. The parameters
$\beta$ and $\delta$ capture the sensitivity of spending to interest
rates and prices, respectively.

\textbf{Shifts in AD:} The AD curve shifts rightward (increasing output
demand at every price level) when:
\begin{itemize}
  \item Government spending $G$ increases (fiscal expansion)
  \item The central bank cuts interest rates $r$ (monetary expansion)
  \item Household wealth or confidence rises
  \item Export demand increases (trading partner growth or currency depreciation)
\end{itemize}

%----------------------------------------------------------
\subsection{Aggregate Supply}
%----------------------------------------------------------

The \term{Short-Run Aggregate Supply (SRAS) curve} shows the total quantity
of goods and services firms are willing to produce at each price level,
given that input prices (especially wages) are fixed in the short run.
It slopes upward because higher output prices increase profit margins,
incentivizing firms to expand production.

A simple SRAS relationship:

\begin{equation}
  P = P^e + \frac{1}{\alpha}(Y - Y_f)
  \label{eq:sras}
\end{equation}

\noindent where $P$ is the current price level, $P^e$ is the expected
price level (determining the position of the SRAS curve), $Y_f$ is
potential output (full-employment output), and $\alpha$ is a positive
parameter governing how responsive prices are to output gaps.

The \term{Long-Run Aggregate Supply (LRAS) curve} is vertical at
$Y = Y_f$ --- potential output. In the long run, all prices are flexible,
the economy returns to $Y_f$, and the price level adjusts to whatever is
required for this equilibrium. Monetary policy cannot permanently raise
output above $Y_f$; it can only affect the price level in the long run.

%----------------------------------------------------------
\subsection{Two Types of Inflation}
%----------------------------------------------------------

\begin{figure}[H]
\centering
\begin{tikzpicture}[
  node distance=0cm,
  axisline/.style={-Stealth, thick, color=gray!70},
  curvep/.style={thick, color=policyblue},
  curvea/.style={thick, color=casered},
  curved/.style={thick, color=unbcgreen},
  curveas/.style={thick, color=dataorange},
  labelstyle/.style={font=\small}
]

% ---- PANEL A: DEMAND-PULL ----
\begin{scope}[xshift=0cm]
\node[font=\small\bfseries] at (3.2, 5.3) {Panel A: Demand-Pull Inflation};
% Axes
\draw[axisline] (0,0) -- (0,5) node[above, font=\small] {$P$};
\draw[axisline] (0,0) -- (6.2,0) node[right, font=\small] {$Y$};
% LRAS vertical
\draw[dashed, gray, thick] (3.5,0) -- (3.5,4.8)
  node[above, font=\footnotesize] {LRAS};
% SRAS curve (upward sloping)
\draw[curveas, thick] (0.5,0.5) .. controls (2,1.8) .. (5.5,4.5)
  node[right, font=\footnotesize, color=dataorange] {SRAS};
% AD1 (original)
\draw[curved] (0.5,4.2) .. controls (2.5,2.5) .. (5.5,0.5)
  node[right, font=\footnotesize, color=unbcgreen] {AD$_1$};
% AD2 (rightward shift)
\draw[curvep, dashed] (1.5,4.5) .. controls (3.5,2.8) .. (6.0,0.8)
  node[right, font=\footnotesize, color=policyblue] {AD$_2$};
% Equilibria
\fill[unbcgreen] (3.0,2.0) circle (3pt);
\fill[policyblue] (4.2,2.9) circle (3pt);
% Labels
\draw[dashed, gray!50, thin] (3.0,0) -- (3.0,2.0) -- (0,2.0);
\draw[dashed, gray!50, thin] (4.2,0) -- (4.2,2.9) -- (0,2.9);
\node[left, font=\footnotesize] at (0,2.0) {$P_1$};
\node[left, font=\footnotesize] at (0,2.9) {$P_2$};
\node[below, font=\footnotesize] at (3.0,0) {$Y_1$};
\node[below, font=\footnotesize] at (4.2,0) {$Y_2$};
\node[below, font=\footnotesize] at (3.5,-0.05) {$Y_f$};
% Arrow showing shift
\draw[->, thick, policyblue!60] (2.8,2.6) -- (3.6,2.6);
\end{scope}

% ---- PANEL B: COST-PUSH ----
\begin{scope}[xshift=7.5cm]
\node[font=\small\bfseries] at (3.2, 5.3) {Panel B: Cost-Push Inflation};
% Axes
\draw[axisline] (0,0) -- (0,5) node[above, font=\small] {$P$};
\draw[axisline] (0,0) -- (6.2,0) node[right, font=\small] {$Y$};
% LRAS vertical
\draw[dashed, gray, thick] (3.5,0) -- (3.5,4.8)
  node[above, font=\footnotesize] {LRAS};
% SRAS1 (original)
\draw[curveas] (0.5,0.5) .. controls (2,1.8) .. (5.5,4.5)
  node[right, font=\footnotesize, color=dataorange] {SRAS$_1$};
% SRAS2 (leftward shift --- cost shock)
\draw[casered, dashed, thick] (0.2,1.5) .. controls (1.5,2.8) .. (4.8,5.0)
  node[right, font=\footnotesize, color=casered] {SRAS$_2$};
% AD curve
\draw[curved] (0.5,4.2) .. controls (2.5,2.5) .. (5.5,0.5)
  node[right, font=\footnotesize, color=unbcgreen] {AD};
% Equilibria
\fill[unbcgreen] (3.0,2.0) circle (3pt);
\fill[casered] (2.1,2.9) circle (3pt);
% Labels
\draw[dashed, gray!50, thin] (3.0,0) -- (3.0,2.0) -- (0,2.0);
\draw[dashed, gray!50, thin] (2.1,0) -- (2.1,2.9) -- (0,2.9);
\node[left, font=\footnotesize] at (0,2.0) {$P_1$};
\node[left, font=\footnotesize] at (0,2.9) {$P_2$};
\node[below, font=\footnotesize] at (3.0,0) {$Y_1=Y_f$};
\node[below, font=\footnotesize] at (2.1,0) {$Y_2$};
% Arrow showing leftward SRAS shift
\draw[->, thick, casered!70] (2.0,2.3) -- (1.2,2.3);
% Stagflation label
\node[font=\footnotesize, casered] at (3.5,4.0)
  {\textit{Stagflation}};
\node[font=\footnotesize, casered] at (3.5,3.6)
  {($P\uparrow$, $Y\downarrow$)};
\end{scope}

\end{tikzpicture}
\caption{Demand-Pull vs.\ Cost-Push Inflation in the AD-AS Framework.
\textbf{Panel A (Demand-Pull):} A rightward shift of AD (from AD$_1$ to
AD$_2$) --- driven by fiscal stimulus, low interest rates, or rising
household wealth --- moves the economy from equilibrium $(Y_1, P_1)$ to
$(Y_2, P_2)$, raising both output and prices. If $Y_2 > Y_f$, the economy
is overheating. \textbf{Panel B (Cost-Push):} A leftward shift of SRAS ---
driven by higher energy prices, supply chain disruptions, or rising input
costs --- produces the combination of higher prices \textit{and} lower
output known as \textbf{stagflation}. Canada's 2021--2023 inflation episode
involved both mechanisms simultaneously: massive AD expansion from COVID
stimulus and near-zero rates, combined with supply chain bottlenecks and
the Russian-invasion energy shock.}
\label{fig:adas}
\end{figure}

\begin{canexample}{Canada's 2021--2023 Inflation: Both Demand and Supply}
Canada's inflation surge was unusual because both AD and AS shocks
occurred simultaneously and reinforced each other.

\textbf{Demand side:} The federal government injected approximately
\$350 billion in emergency spending between 2020 and 2022. The Bank of
Canada held its policy rate at 0.25\% until March 2022. Household savings
rates surged during lockdowns (\$250+ billion in excess savings accumulated),
creating pent-up demand that exploded when the economy reopened. These
forces shifted AD rightward continuously from mid-2020 through 2021.

\textbf{Supply side:} Global shipping container costs rose tenfold.
Semiconductor shortages idled automotive production lines --- used car
prices in Canada rose over 40\% in 2021. Russia's invasion of Ukraine
(February 2022) triggered an energy and food commodity shock, shifting
SRAS leftward. Labour shortages in food processing, transportation, and
construction added further upward cost pressure.

\textbf{Result:} In the AD-AS diagram, Canada experienced both a large
rightward AD shift and a leftward SRAS shift. Prices rose sharply. Output
expanded beyond potential in 2021 (generating an ``output gap'' of
approximately +2\% of potential GDP) before the rate hikes began
cooling demand. This combination --- demand exceeding supply on both
sides --- is precisely what generates the most acute inflationary episodes.
\end{canexample}

%==========================================================
\section{The Phillips Curve: Inflation and Unemployment}
\label{sec:phillips_curve}
%==========================================================

The \term{Phillips Curve} describes the empirical and theoretical
relationship between inflation and unemployment. The original version,
identified by economist A.W. Phillips (1958) using UK data from 1861--1957,
showed a stable negative relationship: when unemployment was low, wages
(and thus prices) tended to rise faster.

%----------------------------------------------------------
\subsection{The Short-Run Phillips Curve}
%----------------------------------------------------------

The short-run Phillips Curve can be expressed as:

\begin{equation}
  \pi_t = \pi^e_t - \alpha(u_t - u_n)
  \label{eq:phillips_curve}
\end{equation}

\noindent where:
\begin{itemize}
  \item $\pi_t$ is the actual inflation rate in period $t$
  \item $\pi^e_t$ is the expected inflation rate (which shifts the curve)
  \item $u_t$ is the actual unemployment rate
  \item $u_n$ is the \term{natural rate of unemployment} (the NAIRU ---
        Non-Accelerating Inflation Rate of Unemployment)
  \item $\alpha > 0$ is a parameter measuring how sensitive inflation is
        to labour market slack
\end{itemize}

The term $(u_t - u_n)$ is the \term{unemployment gap}. When $u_t < u_n$
(labour markets are tight, unemployment below natural rate), there is
upward pressure on wages and prices: $\pi_t > \pi^e_t$. When
$u_t > u_n$ (slack in labour markets), inflation tends to run below
expectations: $\pi_t < \pi^e_t$.

For Canada, the Bank of Canada estimates the natural rate of unemployment
$u_n$ at approximately 5.5--6.0\% in the mid-2020s (accounting for
structural factors such as the geographic mismatch between job vacancies
and unemployed workers, and skills mismatches). When unemployment fell
to 4.9\% in mid-2022, the gap $(u_t - u_n) \approx -1.0$ to $-1.5$
percentage points, contributing upward pressure on wages and prices.

%----------------------------------------------------------
\subsection{The Long-Run Phillips Curve}
%----------------------------------------------------------

If workers have \term{rational expectations} --- they correctly anticipate
inflation on average --- then any attempt by the central bank to exploit
the inflation--unemployment trade-off will be defeated. If the Bank
tolerates higher inflation to keep unemployment below the natural rate,
workers eventually demand higher nominal wages to compensate for the
higher expected inflation, shifting the SRAS curve upward and returning
the economy to $u_n$ --- but now at a permanently higher inflation rate.

The long-run Phillips Curve is therefore \textbf{vertical at $u = u_n$}:
there is no long-run trade-off between inflation and unemployment. This
result, developed independently by Edmund Phelps and Milton Friedman in
1968, is one of the most important insights in macroeconomics. It implies:
\begin{itemize}
  \item Monetary policy cannot permanently lower unemployment below $u_n$.
  \item Attempting to do so only generates accelerating inflation.
  \item The only permanent gains in employment come from structural
        policies that reduce $u_n$ itself (better job matching, education,
        reduced geographical barriers to labour mobility).
\end{itemize}

%----------------------------------------------------------
\subsection{The Sacrifice Ratio}
%----------------------------------------------------------

When a central bank wants to \textit{reduce} inflation that has become
embedded in expectations, it must push unemployment above the natural
rate --- accepting a period of higher unemployment. The \term{sacrifice
ratio} measures the cost of this disinflation:

\begin{equation}
  \text{Sacrifice Ratio} = \frac{\text{Cumulative excess unemployment (\%)}}
                                 {\text{Reduction in inflation (pp)}}
  \label{eq:sacrifice_ratio}
\end{equation}

Empirical estimates for Canada typically place the sacrifice ratio between
1.5 and 3.0: to permanently reduce inflation by 1 percentage point requires
approximately 1.5--3.0 percentage points of excess unemployment sustained
for one year. For the 2022--2024 disinflation of approximately 4 percentage
points (from 6.8\% to 2.6\%), the implied employment cost was significant.

\begin{figure}[H]
\centering
\begin{tikzpicture}
\begin{axis}[
  width=10cm, height=8cm,
  xlabel={Unemployment Rate (\%)},
  ylabel={CPI Inflation, Year-over-Year (\%)},
  xlabel style={font=\small},
  ylabel style={font=\small},
  xmin=4.5, xmax=10.5,
  ymin=-0.5, ymax=8.5,
  xtick={5,6,7,8,9,10},
  ytick={0,1,2,3,4,5,6,7,8},
  xticklabel style={font=\small},
  yticklabel style={font=\small},
  ymajorgrids=true, xmajorgrids=true,
  grid style={dashed, gray!20},
  clip=false
]
% 2% inflation target line
\addplot[unbcgreen!50, dashed, thin] coordinates {(4.5,2)(10.5,2)};
% Natural rate reference
\addplot[gray!50, dotted, thin] coordinates {(5.8,0)(5.8,8.5)}
  node[above, font=\footnotesize, text=gray] {$u_n \approx 5.8\%$};
% Data points with year labels
\addplot[only marks, mark=*, mark size=3pt, policyblue]
  coordinates {
    (6.9,1.1)(7.0,1.4)(6.3,1.6)(5.8,2.3)(5.7,1.9)
    (9.5,0.7)(7.5,3.4)(5.3,6.8)(5.4,3.9)(6.3,2.6)(6.7,2.1)(6.5,2.0)
  };
% Year labels
\node[font=\footnotesize, policyblue, above right] at (axis cs:6.9,1.1) {15};
\node[font=\footnotesize, policyblue, above right] at (axis cs:7.0,1.4) {16};
\node[font=\footnotesize, policyblue, above right] at (axis cs:6.3,1.6) {17};
\node[font=\footnotesize, policyblue, above right] at (axis cs:5.8,2.3) {18};
\node[font=\footnotesize, policyblue, above right] at (axis cs:5.7,1.9) {19};
\node[font=\footnotesize, policyblue, right] at (axis cs:9.5,0.7) {20};
\node[font=\footnotesize, policyblue, above right] at (axis cs:7.5,3.4) {21};
\node[font=\footnotesize, policyblue, above] at (axis cs:5.3,6.8) {22};
\node[font=\footnotesize, policyblue, right] at (axis cs:5.4,3.9) {23};
\node[font=\footnotesize, policyblue, below right] at (axis cs:6.3,2.6) {24};
\node[font=\footnotesize, policyblue, below right] at (axis cs:6.7,2.1) {25};
\node[font=\footnotesize, policyblue, below right] at (axis cs:6.5,2.0) {26};
% Arrow showing 2022 surge direction
\draw[->, casered, thick] (axis cs:7.5,3.4) -- (axis cs:5.4,6.5);
\node[font=\footnotesize, casered, rotate=45] at (axis cs:6.4,5.2) {rate surge};
% Arrow showing disinflation
\draw[->, unbcgreen!80, thick] (axis cs:5.4,6.7) to[bend right=15] (axis cs:6.7,2.2);
\node[font=\footnotesize, unbcgreen!80] at (axis cs:4.9,4.5) {disinflation};
\end{axis}
\end{tikzpicture}
\caption{Phillips Curve: Canada, 2015--2026. Each point plots the annual
unemployment rate against the annual average CPI inflation rate, labelled
by the last two digits of the year. The broad negative relationship
--- lower unemployment associated with higher inflation --- is visible
in 2015--2022, consistent with the short-run Phillips Curve. The 2020
point (year 20) is an outlier: unemployment spiked to 9.5\% while
inflation collapsed to 0.7\%, reflecting the demand destruction of
the COVID lockdowns. The 2022 point (year 22) sits at the extreme
top-left: very low unemployment (5.3\%) and very high inflation (6.8\%),
reflecting a severely overheated economy. The red arrow traces the
surge and the green arrow traces the subsequent disinflation path,
along which unemployment rose as inflation fell.}
\label{fig:phillips_curve}
\source{Statistics Canada, Tables 14-10-0287-01 and 18-10-0004-01.
Annual averages; 2025--2026 estimated.}
\end{figure}

%==========================================================
\section{Canada's Inflation-Targeting Framework}
\label{sec:inflation_targeting}
%==========================================================

Canada was the second country in the world (after New Zealand in 1990)
to adopt a formal \term{inflation target}. The framework was introduced
in February 1991 under a joint agreement between the Bank of Canada and
the federal government. It has been renewed multiple times, most recently
in 2021 with modifications that reflected the lessons of the COVID era.

The key features of the current framework:

\begin{itemize}
  \item \textbf{Target:} 2\% measured CPI inflation.
  \item \textbf{Control range:} 1--3\% (the Bank uses the range as an
        indicator of acceptable outcomes, not as hard bounds).
  \item \textbf{Horizon:} 18--24 months --- the time over which the Bank
        aims to return inflation to the 2\% target following a shock.
        This reflects the lags in monetary policy transmission.
  \item \textbf{Instrument:} The overnight rate (policy rate).
  \item \textbf{Communication:} Eight pre-announced Monetary Policy
        Announcement (MPA) dates per year, accompanied by a quarterly
        Monetary Policy Report (MPR) with detailed economic forecasts.
  \item \textbf{Accountability:} The Bank Governor must write an open
        letter to the Minister of Finance explaining why inflation has
        moved more than 1 percentage point above or below the 2\% target.
\end{itemize}

\textbf{Why targeting works --- the role of credibility:}
Inflation targeting is more than a rule; it is a commitment device.
When firms, workers, and financial markets believe the Bank will keep
inflation at 2\%, they set wages and prices consistent with 2\% inflation.
This self-fulfilling property --- ``credibility'' --- makes disinflation
less costly (smaller sacrifice ratio) and makes it easier to keep
inflation at target even when shocks occur. If the Bank loses credibility,
inflation expectations become de-anchored, requiring more aggressive
tightening to achieve the same disinflationary effect.

The 2021--2022 episode tested this credibility severely. By maintaining
the policy rate at 0.25\% while inflation surpassed 5\% --- citing
concerns about premature withdrawal of pandemic support --- the Bank
allowed long-term inflation expectations to creep above 2\%. By the time
rates rose aggressively, some of the benefit of credibility had been
partially eroded, requiring more tightening than if the Bank had acted
earlier.

%==========================================================
\section{Monetary Policy Transmission}
\label{sec:transmission}
%==========================================================

When the Bank of Canada changes its overnight rate, the effects on inflation
and output do not occur immediately. There are multiple channels through
which the policy rate change propagates through the economy, and each
channel operates with different lags. The full effect on inflation is
typically estimated to take \textbf{18--24 months}.

\begin{defbox}{The Four Channels of Monetary Policy Transmission}
\textbf{1. Credit Cost Channel:} Higher overnight rates raise the cost
of borrowing for households and firms. Mortgage rates, car loan rates,
business credit lines, and corporate bond yields all rise. This reduces
investment spending (firms borrow less to expand) and consumer spending
(households reduce debt-financed purchases). This is the fastest and
most direct channel.

\textbf{2. Asset Price Channel:} Higher interest rates reduce the
present value of future cash flows, lowering stock prices, bond prices,
and real estate valuations. Lower household wealth reduces consumption
through the \term{wealth effect}. A fall in equity prices also raises
firms' cost of capital, further reducing investment.

\textbf{3. Exchange Rate Channel:} Higher Canadian interest rates attract
foreign capital seeking superior returns, increasing demand for Canadian
dollars and appreciating the exchange rate. A stronger Canadian dollar
makes imports cheaper (reducing the domestic price of imported goods)
and makes Canadian exports less competitive internationally (reducing
net exports). Both effects reduce domestic price pressure.

\textbf{4. Expectations Channel:} When the Bank credibly signals
its inflation intentions, households and firms adjust their price- and
wage-setting behaviour immediately, \textit{before} the other channels
fully operate. If workers and firms believe the Bank will achieve 2\%
inflation, they set wages and prices consistent with 2\% --- effectively
doing the Bank's job for it. This is why central bank communication
(\term{forward guidance}) is itself a monetary policy instrument.
\end{defbox}

The relative importance of these channels in Canada reflects the
structure of the Canadian economy. The credit cost channel is particularly
powerful because:

\begin{enumerate}
  \item Canada has one of the highest household debt-to-income ratios
        among OECD countries (~180\% as of 2023), making households
        extremely sensitive to interest rate changes.
  \item Approximately 35\% of Canadian mortgages were in variable-rate
        products as of 2022, meaning rate increases transmitted to
        mortgage payments with minimal lag.
  \item Housing wealth constitutes approximately 40\% of total Canadian
        household net worth, making the asset price channel powerful.
\end{enumerate}

%==========================================================
\section{Case Studies}
\label{sec:ch2_casestudies}
%==========================================================

%----------------------------------------------------------
\begin{casestudy}{The Bank of Canada's 2022--2024 Rate-Hike Cycle}

\textbf{Background:} By early 2022, Canada's inflation had risen to levels
not seen since the early 1980s. The Bank of Canada, which had kept its
policy rate at the emergency floor of 0.25\% since March 2020, faced the
challenge of rapidly tightening monetary policy without triggering a
severe recession.

\textbf{The Timeline:} The tightening cycle began on March 2, 2022, with
a 25-basis-point increase --- the Bank's first rate hike since October 2018.
What followed was the most rapid tightening cycle in the Bank's modern
history.

\begin{table}[H]
\centering
\caption{Bank of Canada Rate Decisions, March 2022 -- December 2024}
\label{tab:rate_decisions}
\renewcommand{\arraystretch}{1.25}
\begin{tabular}{lrrl}
\toprule
\textbf{Date} & \textbf{Change (bps)} & \textbf{New Rate (\%)}
              & \textbf{Bank's Primary Rationale} \\
\midrule
Mar 2, 2022  & +25  & 0.50 & Inflation exceeding target; recovery solid \\
Apr 13, 2022 & +50  & 1.00 & Inflation broad-based; front-loading needed \\
Jun 1, 2022  & +50  & 1.50 & CPI 7.7\% in Apr; labour markets tight \\
Jul 13, 2022 & +100 & 2.50 & CPI 8.1\% June peak; unconventional move \\
Sep 7, 2022  & +75  & 3.25 & Core measures still rising \\
Oct 26, 2022 & +50  & 3.75 & Moderate; economy showing some cooling \\
Dec 7, 2022  & +50  & 4.25 & Inflation still above target \\
Jan 25, 2023 & +25  & 4.50 & Conditional pause signalled \\
Mar 8, 2023  & +25  & 4.75 & Core inflation re-accelerated \\
Jun 7, 2023  & +25  & 5.00 & GDP and CPI stronger than expected \\
\midrule
\textit{Jul--Dec 2023} & \textit{Held} & \textit{5.00} & \textit{Assessment period} \\
\midrule
Jun 5, 2024  & $-$25 & 4.75 & Inflation returning to target \\
Jul 24, 2024 & $-$25 & 4.50 & Labour market easing \\
Sep 4, 2024  & $-$25 & 4.25 & CPI at 2.5\%; output gap opened \\
Oct 23, 2024 & $-$50 & 3.75 & Downside risks growing \\
Dec 11, 2024 & $-$50 & 3.25 & Inflation at target; growth below potential \\
\bottomrule
\end{tabular}
\source{Bank of Canada, Monetary Policy Announcements, 2022--2024.
``bps'' = basis points (1 basis point = 0.01 percentage point).}
\end{table}

\textbf{The July 2022 Anomaly:} The July 13, 2022 decision --- a 100-basis-point
increase, the largest single-meeting rate hike since 1998 --- was remarkable.
The Bank cited the June 2022 CPI print of 8.1\% and the risk that
inflation expectations would become ``de-anchored'' from the 2\% target.
At that point, the Bank was willing to accept a sharp slowdown rather than
allow inflation to become entrenched.

\textbf{Outcomes:} By June 2024, inflation had returned to 2.7\% --- near
the 2\% target. Canada avoided a technical recession (two consecutive
quarters of negative real GDP growth) in 2022--2023, though real GDP per
capita did decline. The unemployment rate rose from a 50-year low of 4.9\%
in mid-2022 to 6.5--6.8\% by late 2024.

\textbf{Evaluation:} Did the Bank of Canada succeed? On price stability,
yes: inflation returned to target within two years of the June 2022 peak.
On the real economy, the outcome was relatively benign by historical
standards --- no deep recession, modest unemployment increase, and nominal
wage growth that largely kept pace with declining inflation. However,
mortgage holders who renewed into 5\% rates experienced the most severe
payment shock in a generation, and some economists argue the Bank held
rates too high for too long in 2023--24, unnecessarily slowing growth.
\end{casestudy}

%----------------------------------------------------------
\begin{casestudy}{Shelter Inflation: The CPI's Most Powerful Component}

\textbf{The Issue:} Table~\ref{tab:cpi_components} shows that shelter
inflation remained elevated --- and in some measures \textit{accelerating}
--- into 2023 and 2024, even as energy and goods prices eased. Mortgage
interest costs rose 28.5\% in 2023 and 19.4\% in 2024 year-over-year,
making shelter the single largest driver of measured inflation in the
post-peak period.

\textbf{The Paradox:} The Bank of Canada raised interest rates to fight
inflation. But because mortgage interest costs are included in the CPI
shelter component (they are part of the cost of ``owning'' a home),
the rate hikes directly increased the CPI they were supposed to reduce.
This creates a short-run paradox: tighter monetary policy mechanically
raises measured inflation through the mortgage channel, even as it slows
the economy and reduces demand-side inflation pressure.

\textbf{Why Does This Happen in Canada?}
Canada's CPI shelter methodology, unlike that of many other countries,
includes mortgage interest cost (MIC) directly. The US CPI, by contrast,
uses \term{owners' equivalent rent} --- an estimate of what homeowners
would pay if they rented their homes --- which is more insulated from
short-run interest rate movements.

\textbf{Policy Implication:} This is one reason the Bank of Canada
monitors core inflation measures (particularly CPI-trim and CPI-median)
that tend to dampen the direct mortgage cost effect. It is also why
economists sometimes argue that the headline CPI in Canada gives a
misleading picture of ``true'' cost-of-living inflation during monetary
tightening cycles.

\textbf{Quantitative impact:} In 2023, mortgage interest costs contributed
approximately +1.4 percentage points to Canada's headline inflation of 3.9\%
--- meaning that without the mortgage component, headline inflation would
have been closer to 2.5\%. The same rate hikes that were cooling inflation
were adding roughly 1.4 points to the inflation number, creating both
an economic paradox and a political communications challenge for the Bank.
\end{casestudy}

%----------------------------------------------------------
\begin{casestudy}{Who Pays for Inflation? Distributional Consequences Across Income Groups}

\textbf{Background:} The CPI measures inflation for the ``average'' household.
But average households do not exist --- spending patterns, asset holdings,
debt levels, and income sources vary enormously across the income distribution.
The 2022 inflation surge was not experienced equally.

\textbf{The Regressivity of Inflation:} Low-income households devote larger
shares of their budgets to food, shelter, and energy --- the three categories
that rose fastest in 2022. Statistics Canada analysis estimates that the
effective inflation rate experienced by the bottom income quintile in
2022 was approximately 8.5\%, compared to approximately 6.8\% for the top
quintile. A gap of 1.7 percentage points may seem small in absolute terms,
but for a household spending 95\% of its income, the difference between
an 8.5\% and 6.8\% inflation rate is the difference between severe hardship
and moderate difficulty.

\textbf{Debtors and Creditors:} Unexpected inflation redistributes wealth
from creditors to debtors. Anyone holding a fixed-rate loan from before
the inflation surge found that the real value of their debt was being
eroded by inflation. Mortgage holders with 5-year fixed rates locked in
at 2\% in 2020 enjoyed negative real borrowing costs throughout 2022 ---
their 2\% nominal rate minus 6\%+ inflation implied a real borrowing cost
of approximately $-4\%$. Retirees and pension funds holding fixed-income
securities saw the real value of their assets eroded.

\textbf{Wage-earners vs.\ Capital owners:} Workers whose wages did not
fully adjust to inflation experienced real pay cuts. Workers in unionized
sectors with cost-of-living adjustment (COLA) clauses in their contracts
were partially protected. Non-unionized, part-time, and gig workers ---
disproportionately low-income, racialized, and young --- had the least
wage protection. Capital owners who held real assets (housing, stocks,
commodities) saw their nominal values rise with inflation, preserving
real wealth while nominal wealth grew.

\textbf{Policy Implications:} The distributional consequences of inflation
strengthen the case for price stability as a distributive goal, not just
an efficiency goal. Inflation is not a politically neutral phenomenon:
it systematically favours debtors over creditors, real asset holders
over cash savers, and the flexible-wage over the fixed-wage workforce.
\end{casestudy}

%==========================================================
\section{Distributional Effects of Inflation}
\label{sec:distributional}
%==========================================================

The case study above is worth formalising. Figure~\ref{fig:inflation_burden}
presents estimated effective inflation rates by income quintile for Canada
in 2022, constructed by reweighting CPI components using the expenditure
patterns of each quintile from the Survey of Household Spending.

\begin{figure}[H]
\centering
\begin{tikzpicture}
\begin{axis}[
  ybar,
  bar width=0.8cm,
  width=10cm, height=6.5cm,
  symbolic x coords={{Q1 (Lowest)},{Q2},{Q3},{Q4},{Q5 (Highest)}},
  xtick=data,
  x tick label style={font=\small, rotate=20, anchor=east},
  ylabel={Effective Inflation Rate (\%), 2022},
  ylabel style={font=\small},
  ymin=5.5, ymax=9.5,
  ytick={6.0,6.5,7.0,7.5,8.0,8.5,9.0},
  yticklabel style={font=\small},
  ymajorgrids=true,
  grid style={dashed, gray!25},
  enlarge x limits=0.15,
  nodes near coords,
  nodes near coords style={font=\small\bfseries},
  every axis plot/.append style={fill=casered!65, draw=casered!85}
]
\addplot coordinates {
  ({Q1 (Lowest)}, 8.5)
  ({Q2},          7.8)
  ({Q3},          7.5)
  ({Q4},          7.2)
  ({Q5 (Highest)},6.8)
};
% Reference line for official headline CPI
\addplot[unbcgreen, dashed, thick, forget plot] coordinates
  {({Q1 (Lowest)},6.8)({Q5 (Highest)},6.8)}
  node[right, font=\footnotesize, color=unbcgreen, xshift=2pt]
  {Headline CPI (6.8\%)};
\end{axis}
\end{tikzpicture}
\caption{Estimated Effective Inflation Rate by Income Quintile, Canada, 2022.
Low-income households (Q1) experienced approximately 8.5\% effective inflation
--- 1.7 percentage points above the headline CPI of 6.8\% --- because
they spend a larger share of their budgets on food, shelter, and energy,
the three categories with the sharpest price increases in 2022. The
dashed line shows the official headline CPI for reference. Inflation
is not a neutral tax: it falls disproportionately on those least able to
absorb it.}
\label{fig:inflation_burden}
\source{Author calculations using Statistics Canada Survey of Household
Spending and CPI component data. Methodology follows Argente and Lee (2021)
and Statistics Canada analytical studies.}
\end{figure}

%==========================================================
\section{The Rate-Hike Cycle: Nominal and Real Rates}
\label{sec:rate_cycle}
%==========================================================

\begin{figure}[H]
\centering
\begin{tikzpicture}
\begin{axis}[
  width=14cm, height=7cm,
  xlabel={Year},
  ylabel={Interest Rate / Inflation Rate (\%)},
  xlabel style={font=\small},
  ylabel style={font=\small},
  xmin=2017.8, xmax=2026.5,
  ymin=-6, ymax=7,
  xtick={2018,2019,2020,2021,2022,2023,2024,2025,2026},
  xticklabel style={font=\small},
  ytick={-6,-4,-2,0,2,4,6},
  yticklabel style={font=\small},
  ymajorgrids=true,
  grid style={dashed, gray!20},
  legend pos=north west,
  legend style={font=\small},
  clip=false
]
% Zero line
\draw[gray!50, thin] (axis cs:2017.8,0) -- (axis cs:2026.5,0);
% Neutral real rate reference
\addplot[gray!40, dotted] coordinates {(2017.8,1.0)(2026.5,1.0)}
  node[right, font=\footnotesize, color=gray!60] {$r^* \approx 1\%$};
% Overnight rate (nominal)
\addplot[policyblue, thick, mark=square*, mark size=2pt] coordinates {
  (2018,1.75)(2019,1.75)(2020,0.25)(2021,0.25)
  (2022,4.25)(2023,5.00)(2024,3.25)(2025,3.00)(2026,2.75)
};
\addlegendentry{Nominal overnight rate $i$}
% Real interest rate (i - pi)
\addplot[casered, thick, mark=triangle*, mark size=2.5pt] coordinates {
  (2018,-0.55)(2019,-0.15)(2020,-0.45)(2021,-3.15)
  (2022,-2.55)(2023,1.10)(2024,0.65)(2025,0.90)(2026,0.75)
};
\addlegendentry{Real rate $r = i - \pi$}
% Shaded negative real rate area
\fill[casered!10] (axis cs:2017.8,0) -- (axis cs:2018,-0.55)
  -- (axis cs:2019,-0.15) -- (axis cs:2020,-0.45) -- (axis cs:2021,-3.15)
  -- (axis cs:2022,-2.55) -- (axis cs:2022.8,0) -- cycle;
\node[font=\footnotesize, casered, rotate=0] at (axis cs:2020.5,-2.0)
  {Negative real rates};
\node[font=\footnotesize, casered] at (axis cs:2021,-3.5)
  {(fuel for inflation)};
\end{axis}
\end{tikzpicture}
\caption{Nominal Overnight Rate vs.\ Real Interest Rate, Canada, 2018--2026.
The shaded region marks the period of negative real interest rates (2018--2022).
In 2021, the real rate fell to approximately $-3.2\%$ --- deeply negative ---
as the nominal rate was held at 0.25\% while inflation rose to 3.4\%.
By mid-2022, the combination of rapid nominal rate increases and peak inflation
meant the real rate was still close to $-4.5\%$ even as nominal rates rose
to 2.5\%. Only in 2023, when nominal rates reached 5.00\% and inflation fell
to 3.9\%, did the real rate finally turn positive. The return to positive
real rates was the mechanism through which monetary policy finally cooled
aggregate demand.}
\label{fig:real_nominal_rates}
\source{Bank of Canada key interest rates; Statistics Canada Table 18-10-0004-01;
author calculations using Fisher approximation $r \approx i - \pi$.}
\end{figure}

%==========================================================
\section*{Policy Debate}
%==========================================================

\begin{policydebate}{Did the Bank of Canada Hike Too Far, Too Fast --- or Not Far Enough, Too Late?}

\textbf{The Issue:} The Bank of Canada's 2022--2024 rate-hike cycle was
the most aggressive in its modern history. Whether it was the right response
depends critically on the diagnosis of what caused the inflation and on
judgements about the costs of different types of policy error.

\textbf{Case 1 --- Too Slow to Act (The Hawkish Critique):}
Critics argue the Bank should have raised rates earlier --- potentially
as early as the summer of 2021, when CPI was already at 3.7\%. By
maintaining emergency-level rates through 2021, the Bank:
\begin{itemize}
  \item Allowed real interest rates to fall to $-3\%$, massively stimulating
        borrowing and spending.
  \item Failed to ``lean against'' surging housing prices, which added
        to the wealth effect and the inflation spiral.
  \item Let inflation expectations drift above 2\%, requiring more
        aggressive tightening later.
\end{itemize}
The implicit sacrifice ratio of this delayed response was higher than
it needed to be: more unemployment was ultimately required to achieve
the same disinflation because expectations had partially de-anchored.

\textbf{Case 2 --- Too Aggressive (The Dovish Critique):}
Others argue the Bank moved too fast and too far, particularly after
mid-2022. By the time the July 2022 100-basis-point hike was made,
global supply chains were already improving, energy prices were
beginning to fall, and the inflation peak was likely already past.
The aggressive tightening:
\begin{itemize}
  \item Imposed mortgage payment shocks on millions of Canadians who had
        no role in causing inflation.
  \item Slowed real GDP per capita growth below zero for five consecutive
        quarters in 2023--24.
  \item Risked pushing unemployment unnecessarily high, creating real
        human costs (job loss, bankruptcy, reduced business investment)
        to achieve a disinflation that was partly happening on its own
        as supply chains normalized.
\end{itemize}

\textbf{The Evidence:}
\begin{itemize}
  \item Canada's peak inflation (8.1\%) was lower than the UK (11.1\%)
        and comparable to the US (9.1\%).
  \item The Bank achieved its inflation target by mid-2024 without a
        technical recession.
  \item Real GDP per capita nonetheless fell for several quarters ---
        a ``per capita recession'' even without the formal definition
        being triggered.
  \item Variable-rate mortgage holders saw monthly payments increase
        by an average of approximately \$700--\$900.
  \item Business insolvencies rose to near-record levels in 2024.
\end{itemize}

\textbf{The Analytical Framework:}
Evaluating central bank decisions requires distinguishing what the Bank
\textit{knew} at the time from what we know now. Using today's data to
critique 2021--2022 decisions is ``Monday morning quarterbacking.'' The
correct standard is: given the information available in real time, were
the Bank's decisions within the range of reasonable policy choices?
Most economists conclude the answer is yes --- though many would have
preferred earlier action in 2021 and possibly a slightly lower terminal
rate in 2023. This is a case where the difference between a good outcome
and a great outcome was perhaps 50--75 basis points of peak rate ---
a reminder of how consequential seemingly small numerical differences
in monetary policy can be.

\textbf{Your Position:} After reading this chapter, where do you stand?
Use the Fisher equation, the AS-AD model, and the Phillips Curve to
structure your argument.
\end{policydebate}

%==========================================================
\begin{everydaylife}{How Inflation Affects Your Grocery Bill --- and What You Can Do About It}

When Statistics Canada reports that food prices rose 9.7\% in 2022,
that average conceals enormous variation. Fats and oils: +23.3\%.
Pasta and other grain products: +19.6\%. Fresh vegetables: +9.5\%.
Restaurant meals: +8.3\%.

As a consumer, you can partially escape measured food inflation through
\textit{substitution} --- exactly the behaviour that creates the CPI's
substitution bias. Switching from branded cereal to store-brand, choosing
chicken over beef, buying a larger format item or catching a sale --- all
of these reduce your \textit{personal} inflation rate below the measured
CPI for food.

This is not just a survival tip. It illustrates a deep economic point:
the CPI measures the cost of a \textit{fixed} basket. Real households
are not fixed-basket consumers --- they respond to price signals. Your
ability to substitute gives you more real purchasing power than the
headline inflation rate implies. The gap between your personal inflation
rate and the CPI, averaged across all consumers, is precisely the
\textit{substitution bias} in CPI measurement discussed in
Section~\ref{sec:cpi_measurement}. Estimated at 0.3--0.5 percentage
points per year on average, this bias was likely larger in 2022 when
relative price changes were unusually large --- giving households more
scope for substitution than normal.
\end{everydaylife}

%==========================================================
\section*{Chapter Summary}
%==========================================================
\addcontentsline{toc}{section}{Chapter Summary}

\begin{itemize}
  \item \textbf{Inflation} is a sustained increase in the general price
        level; \textbf{deflation} is a sustained decrease. Low positive
        inflation (2\% in Canada) is preferred to zero because it provides
        a policy buffer, accommodates relative wage adjustments, and
        partially offsets CPI measurement biases.

  \item The \textbf{CPI} is a Laspeyres index measuring the cost of a
        fixed basket of goods and services. It contains upward biases
        (substitution, quality, new goods, outlet) of approximately
        0.3--0.5 percentage points per year.

  \item \textbf{Core inflation measures} (CPI-trim, CPI-median, CPI-common)
        remove volatile components to reveal the underlying trend. All three
        rose above 5\% in 2022, confirming that the inflation surge was
        broad-based, not merely energy-driven.

  \item The \textbf{Fisher equation} ($r \approx i - \pi$) distinguishes
        nominal from real interest rates. In 2021, the real overnight
        rate in Canada was approximately $-3.2\%$ --- the most stimulative
        monetary stance in modern Canadian history.

  \item The \textbf{AD-AS model} explains two types of inflation:
        \textit{demand-pull} (rightward AD shift, rising prices and output)
        and \textit{cost-push} (leftward SRAS shift, stagflation). Canada's
        2021--2023 surge involved both simultaneously.

  \item The \textbf{short-run Phillips Curve} ($\pi = \pi^e - \alpha(u-u_n)$)
        shows a negative relationship between inflation and unemployment.
        The \textbf{long-run Phillips Curve} is vertical at the natural
        rate $u_n$: monetary policy cannot permanently reduce unemployment
        below $u_n$.

  \item Canada's \textbf{inflation-targeting framework} (since 1991) targets
        2\% CPI inflation within a 1--3\% control range. The Bank of Canada
        conducts policy primarily through the overnight rate.

  \item Monetary policy transmits through four channels: credit cost, asset
        prices, exchange rate, and expectations. Full effects take
        18--24 months.

  \item The 2022--2024 \textbf{rate-hike cycle} raised the overnight rate
        475 basis points to 5.00\% in 16 months --- the fastest tightening
        in modern Canadian history. Inflation returned to target by 2024
        without a technical recession, but at a significant cost to
        mortgage holders.

  \item \textbf{Inflation is regressive}: low-income households experienced
        approximately 8.5\% effective inflation in 2022 vs.\ 6.8\% for
        the top quintile, because they spend disproportionately more on
        food, shelter, and energy.
\end{itemize}

%==========================================================
\section*{Key Terms}
%==========================================================
\addcontentsline{toc}{section}{Key Terms}

\begin{description}[style=nextline, leftmargin=3.5cm]
  \item[\term{CPI (Consumer Price Index)}] A Laspeyres price index measuring
        the cost of a fixed basket of goods; Canada's primary inflation measure.
  \item[\term{CPI-common}] Bank of Canada core measure using factor analysis
        to extract economy-wide price movements.
  \item[\term{CPI-median}] The price change at the 50th percentile of the
        weighted component distribution.
  \item[\term{CPI-trim}] Core inflation measure that removes the top and
        bottom 20\% of price-change outliers.
  \item[\term{Core inflation}] Inflation measures that exclude the most
        volatile CPI components to reveal the underlying trend.
  \item[\term{Cost-push inflation}] Inflation arising from a leftward shift
        of the SRAS curve (rising input costs, supply shocks).
  \item[\term{Deflation}] A sustained fall in the general price level.
  \item[\term{Demand-pull inflation}] Inflation arising from a rightward
        shift of the AD curve (excess demand).
  \item[\term{Ex ante real rate}] The real interest rate based on expected
        future inflation; $r^e = i - \pi^e$.
  \item[\term{Ex post real rate}] The real interest rate based on actual
        realised inflation; $r = i - \pi$.
  \item[\term{Fisher equation}] The relationship $r \approx i - \pi$
        linking real rates, nominal rates, and inflation.
  \item[\term{Inflation targeting}] A monetary policy framework committing
        to keep inflation near a numerical target.
  \item[\term{Laspeyres index}] A price index using base-period quantity
        weights; the methodology underlying Canada's CPI.
  \item[\term{Long-run Phillips Curve}] The vertical relationship between
        inflation and unemployment at the natural rate $u_n$.
  \item[\term{NAIRU}] Non-Accelerating Inflation Rate of Unemployment;
        the unemployment rate consistent with stable inflation.
  \item[\term{Nominal interest rate}] The stated rate on a financial
        instrument; not adjusted for inflation.
  \item[\term{Overnight rate}] The Bank of Canada's policy rate; the rate
        at which major financial institutions lend to each other overnight.
  \item[\term{Phillips Curve}] The relationship between inflation (or wage
        inflation) and unemployment.
  \item[\term{Real interest rate}] The nominal interest rate adjusted for
        inflation; $r = i - \pi$.
  \item[\term{Real wage}] Nominal wage deflated by the price level;
        measures purchasing power of labour income.
  \item[\term{Sacrifice ratio}] Cumulative excess unemployment required
        per percentage-point reduction in inflation.
  \item[\term{Short-run Phillips Curve}] The negative relationship between
        inflation and unemployment at a given inflation expectation.
  \item[\term{Stagflation}] Simultaneous high inflation and high unemployment,
        typically arising from an adverse supply shock.
  \item[\term{Substitution bias}] CPI overstatement due to ignoring
        consumer substitution away from goods with rising relative prices.
  \item[\term{Zero lower bound}] The practical floor on nominal interest
        rates (approximately 0\%); limits the downward scope of monetary
        stimulus.
\end{description}

%==========================================================
\section*{Review Questions}
%==========================================================
\addcontentsline{toc}{section}{Review Questions}

\begin{enumerate}

\item \textbf{(Concept)} Define inflation and explain why a 2\% inflation
target is preferred to 0\% by most central banks. In your answer, discuss
the zero lower bound and the measurement bias arguments.

\item \textbf{(Measurement)} Canada's CPI basket assigns a 30\% weight
to Shelter and 20\% to Transportation. In 2022, Shelter rose 6.9\% and
Transportation rose 10.6\% year-over-year.
\begin{enumerate}[(a)]
  \item What is the combined contribution of these two categories to
        headline inflation, assuming the remaining categories average
        6.5\%?
  \item If all other categories fell to zero growth, what would headline
        inflation be?
  \item What does this illustrate about the sensitivity of CPI to
        its largest components?
\end{enumerate}

\item \textbf{(Fisher Equation)} In January 2021, the Bank of Canada
overnight rate was 0.25\% and CPI inflation was 1.0\%. By December 2022,
the overnight rate had risen to 4.25\% and inflation was 6.3\%.
\begin{enumerate}[(a)]
  \item Calculate the real interest rate in January 2021 and December 2022.
  \item In which period was monetary policy more accommodative (stimulative)?
        Explain using the Fisher equation.
  \item How does your answer change the interpretation of the rate hike
        cycle as ``aggressive''?
\end{enumerate}

\item \textbf{(AD-AS)} Use the AD-AS model to explain Canada's 2022
inflation surge. Draw a diagram showing:
\begin{enumerate}[(a)]
  \item The starting equilibrium at potential output $Y_f$.
  \item The rightward shift of AD resulting from COVID fiscal stimulus.
  \item The leftward shift of SRAS resulting from supply chain disruptions.
  \item The new equilibrium showing both higher prices and a different
        level of output.
\end{enumerate}
Label all curves and equilibria clearly.

\item \textbf{(Phillips Curve)} Using Equation~\ref{eq:phillips_curve}
and estimates $\alpha = 0.5$ and $u_n = 5.8\%$:
\begin{enumerate}[(a)]
  \item What does the model predict for inflation if the unemployment
        rate is 4.8\% and expected inflation is 2\%?
  \item What does the model predict if unemployment is 7.3\% and
        expected inflation is 3\%?
  \item If the Bank of Canada wants to reduce inflation from 5\% to 2\%
        and the sacrifice ratio is 2.0, how much cumulative excess
        unemployment (above $u_n$) would be required?
\end{enumerate}

\item \textbf{(Policy)} Why does the Bank of Canada monitor core inflation
measures (CPI-trim, CPI-median) rather than simply responding to headline
CPI? Under what circumstances would a supply shock cause the Bank to
\textit{tolerate} higher headline inflation without raising rates?

\end{enumerate}

%==========================================================
\section*{Quantitative Problems}
%==========================================================
\addcontentsline{toc}{section}{Quantitative Problems}

\begin{enumerate}

\item \textbf{CPI Calculation.} A household's monthly expenditure basket
consists of four categories:

\begin{center}
\begin{tabular}{lrrr}
\toprule
Category & Weight (\%) & Price (Base Year) & Price (Current Year) \\
\midrule
Food     & 20 & \$600 & \$672 \\
Housing  & 35 & \$1{,}200 & \$1{,}272 \\
Transport& 25 & \$400 & \$460 \\
Other    & 20 & \$300 & \$318 \\
\bottomrule
\end{tabular}
\end{center}

\begin{enumerate}[(a)]
  \item Calculate the Laspeyres CPI for the current year
        (base year = 100) using Equation~\ref{eq:cpi_formula}.
        Show all steps.
  \item Calculate the inflation rate $\pi$ from base year to current year.
  \item Which category contributed most to the measured inflation rate?
        Calculate each category's contribution to $\pi$.
  \item Suppose the household switches entirely from Transport to a
        category priced at base-year levels. What is the household's
        \textit{true} cost-of-living increase? How does it compare
        to the Laspeyres CPI answer? What bias does this illustrate?
\end{enumerate}

\item \textbf{Fisher Equation Applications.}
\begin{enumerate}[(a)]
  \item A 5-year Government of Canada bond yields 4.8\% nominal.
        If expected inflation over the next 5 years is 2.2\%, what
        is the expected real yield? Is this above or below the
        historical average neutral real rate of approximately 1\%?
  \item A student takes a variable-rate student loan at prime + 1.5\%.
        In January 2020, the prime rate is 3.45\% (overnight + 2.2\%).
        In December 2022, the prime rate is 6.45\%. The student's income
        is \$42{,}000 per year. CPI in 2020 was 137.0; in 2022 it was 151.9.
        Calculate the real interest rate on the student loan in each year
        and the real value of \$42{,}000 income in each year (in 2020 dollars).
  \item Canada runs a current account deficit when nominal interest rates
        are below trading partners' rates (capital flows out). In 2020, the
        Bank of Canada cut rates to 0.25\% while the ECB kept rates at
        $-0.50\%$ and the Fed cut to 0--0.25\%. Who had the relatively
        highest real rates given the respective inflation differentials?
        What does this imply for capital flows?
\end{enumerate}

\item \textbf{Real Wages and Cost of Living.}
Average weekly earnings in Canada were \$1{,}020 in January 2020
and \$1{,}147 in December 2022. CPI (all-items, 2002=100) was
137.9 in January 2020 and 156.0 in December 2022.
\begin{enumerate}[(a)]
  \item Calculate real average weekly earnings in each period
        (in January 2020 dollars) using Equation~\ref{eq:real_wage}.
  \item Calculate the nominal and real growth rates of weekly earnings
        over the period.
  \item If low-income workers (bottom quintile) earned \$720/week in
        January 2020 and \$780/week in December 2022, but experienced
        an effective inflation rate of 13.8\% over that period, what
        was the change in their real wages?
  \item What economic term describes a situation where nominal wages
        rise but real wages fall?
\end{enumerate}

\item \textbf{Phillips Curve and Sacrifice Ratio.}
Suppose the short-run Phillips Curve for Canada is:
$\pi_t = \pi^e_t - 0.6(u_t - 5.8\%)$
\begin{enumerate}[(a)]
  \item If expected inflation is 2\% and actual inflation is 4.5\%,
        what does the Phillips Curve imply about the unemployment rate
        at that time?
  \item The Bank of Canada wishes to reduce inflation from 4.5\% to 2\%
        over three years. In each year, the Bank accepts unemployment
        1.2 percentage points above the natural rate. What is the implied
        sacrifice ratio? Is this within the range of estimates in the
        literature (1.5--3.0)?
  \item If inflation expectations become fully rational (people correctly
        anticipate the Bank's actions), what happens to the sacrifice ratio?
        What does this tell us about the value of central bank credibility?
\end{enumerate}

\item \textbf{Distributional Analysis.}
Use the following simplified data to analyse distributional inflation:

\begin{center}
\begin{tabular}{lrrrr}
\toprule
Category & CPI Weight & 2022 Rise & Q1 Budget Share & Q5 Budget Share \\
\midrule
Food       & 16\% & 9.7\% & 22\% & 10\% \\
Shelter    & 30\% & 6.9\% & 38\% & 22\% \\
Energy     & 8\%  &28.5\% & 12\% &  4\% \\
Transport  & 12\% & 7.8\% &  8\% & 16\% \\
Services   & 20\% & 5.5\% & 12\% & 30\% \\
Other      & 14\% & 5.2\% &  8\% & 18\% \\
\bottomrule
\end{tabular}
\end{center}

\begin{enumerate}[(a)]
  \item Calculate the headline CPI inflation rate using official
        CPI weights and 2022 price changes.
  \item Calculate the effective inflation rate for Q1 (lowest-income)
        households using their budget shares as weights.
  \item Calculate the effective inflation rate for Q5 (highest-income)
        households.
  \item By how many percentage points did Q1 experience higher inflation
        than Q5? What drives this difference?
  \item Propose one policy instrument that could specifically compensate
        low-income households for this burden. What would be its fiscal cost
        if Canada has approximately 4 million households in Q1,
        each spending \$40{,}000 per year?
\end{enumerate}

\item \textbf{Data Exercise.} Visit the Bank of Canada website
(\url{www.bankofcanada.ca}) and find:
\begin{enumerate}[(a)]
  \item The current overnight rate target.
  \item The most recent annual CPI inflation reading.
  \item Using these two numbers, calculate the current real overnight rate.
        Is monetary policy currently stimulative, neutral, or restrictive
        (assume the neutral real rate is approximately 1\%)?
  \item Download the CPI-trim and CPI-median series from the Bank of Canada's
        core inflation page. Report whether core measures are above or below
        the 2\% target. What does this imply about the appropriate monetary
        policy stance?
\end{enumerate}

\end{enumerate}

%==========================================================
\section*{Essay Questions}
%==========================================================
\addcontentsline{toc}{section}{Essay Questions}

\begin{enumerate}

\item \textbf{(600--800 words)} ``The Bank of Canada's aggressive rate-hike
cycle in 2022--2023 was necessary and ultimately successful.'' Using the
AD-AS model, the Fisher equation, and the Phillips Curve, evaluate this
claim. Your essay should discuss: (a) the causes of the 2022 inflation surge;
(b) the mechanism through which rate hikes reduce inflation; (c) the costs
of the tightening cycle; and (d) your overall assessment.

\item \textbf{(500--700 words)} The CPI reports a single number for inflation,
but inflation is experienced differently across income groups, age cohorts,
homeowners vs.\ renters, and variable-rate vs.\ fixed-rate mortgage holders.
Using data from this chapter, write an essay arguing that ``the CPI is a
measure of average inflation, but Canada needs complementary distributional
inflation measures to make good policy.'' What additional measures would you
recommend Statistics Canada publish?

\item \textbf{(400--600 words)} Japan maintained near-zero inflation (and
frequently deflation) for much of the 1990s--2010s despite aggressive monetary
stimulus. Canada experienced 8.1\% peak inflation in 2022 despite similar
initial monetary conditions (near-zero rates). Using the AS-AD framework,
propose at least two hypotheses that explain why the same ultra-loose monetary
policy produced such different inflation outcomes in the two countries. What
does this imply about the limits of the quantity theory of money?

\end{enumerate}

%==========================================================
\section*{Discussion Questions}
%==========================================================
\addcontentsline{toc}{section}{Discussion Questions}

\begin{enumerate}

\item Fatima, from our opening vignette, is a variable-rate mortgage holder
who experienced both high inflation (eroding her real income) and sharply
higher mortgage payments (from the rate hikes). Is she a net winner or loser
from the sequence of events? What if she had locked into a 5-year fixed rate
at 2\% in 2020 instead? What if she were a renter with no mortgage?

\item The Bank of Canada held rates at emergency levels (0.25\%) throughout
2021, even as inflation rose above 4\%, citing concerns about pandemic
economic scarring. Was this the right call? What information would you need
to evaluate the Bank's decision in real time, rather than in hindsight?

\item ``Inflation is a hidden tax on savings, a transfer from creditors to
debtors, and a regressive burden on low-income households.'' Do you agree?
Who are the winners and who are the losers from an unexpected inflation
surge? Does your answer change if the inflation is anticipated vs.\ unexpected?

\item Canada's CPI Shelter component includes mortgage interest costs, which
rose sharply as the Bank of Canada raised rates. Should mortgage interest
costs be in the CPI? What are the arguments for and against the current
Canadian methodology? How would switching to ``owners' equivalent rent''
(the US approach) have changed the reported inflation numbers in 2023?

\end{enumerate}

%==========================================================
%  CHAPTER 3: Canada's Housing Crisis — Markets, Affordability,
%              and Policy
%  ECON 204 — Contemporary Economic Issues: A Canadian Perspective
%==========================================================

\chapter{Canada's Housing Crisis: Markets, Affordability, and Policy}
\label{ch:housing}

%----------------------------------------------------------
%  OPENING VIGNETTE
%----------------------------------------------------------

\noindent In February 2022, a detached house in East Vancouver sold for
\$2.1 million. The winning offer was \$400{,}000 above the asking price,
waived all conditions, and was accepted over 22 competing bids. The
buyers --- two engineers in their mid-thirties with combined household
income of approximately \$220{,}000 --- put down \$420{,}000 (20\%) and
took on a mortgage of \$1{,}680{,}000. At that month's prevailing
5-year fixed rate of 3.2\%, their monthly payment was \$8{,}067.

They were the lucky ones. They could afford it.

Across the city, a hospital nurse earning \$72{,}000 per year and a
teacher earning \$65{,}000 per year were renting a two-bedroom apartment
for \$2{,}800 per month --- 40\% of their combined take-home pay. They
had been saving for a down payment for six years, but house prices had
risen faster than their savings every single year. The down payment
they needed was now larger than the house price they had originally
targeted.

And in Prince George --- 800 kilometres to the north --- a single father
working at the mill was paying \$1{,}350 per month for a two-bedroom
apartment in a town where the average household income was \$78{,}000.
He was not in a ``hot market.'' But the housing shortage had spread.

These three households inhabit the same economy, regulated by the same
Bank of Canada, subject to the same federal tax law, and served by the
same CMHC mortgage insurance program. Yet their housing experiences are
radically different --- determined by the intersection of global capital
flows, local zoning bylaws, demographic waves, and macroeconomic policy.

Understanding why Canadian housing became so unaffordable --- and what
could actually fix it --- is one of the most important applied questions
in contemporary Canadian economics.

%----------------------------------------------------------
%  LEARNING OBJECTIVES
%----------------------------------------------------------

\begin{learningobjectives}
  \item Apply the supply-and-demand framework to housing markets,
        explain why the short-run housing supply curve is nearly vertical,
        and predict the price effects of demand shifts.
  \item Derive the user cost of housing formula and use it to evaluate
        the rent-versus-buy decision under different interest rate and
        price appreciation assumptions.
  \item Calculate standard housing affordability measures (price-to-income
        ratio, shelter cost burden, CMHC affordability index) and apply
        the standard mortgage payment formula to scenarios with different
        interest rates.
  \item Explain the supply-side constraints on Canadian housing
        production --- including zoning, land costs, construction labour,
        and permitting --- and evaluate their relative importance.
  \item Analyse the demand-side forces driving Canadian house prices,
        including immigration, financialization, low interest rates, and
        demographic trends.
  \item Describe Canada's rental market dynamics, define the rental
        vacancy rate, and evaluate the economic arguments for and against
        rent control.
  \item Evaluate major federal and provincial housing policy instruments
        using supply-and-demand analysis, and identify which interventions
        address the root cause of unaffordability.
\end{learningobjectives}

%==========================================================
\section{Housing as an Economic Good}
\label{sec:housing_economic_good}
%==========================================================

Housing is an unusual commodity. It is simultaneously a \textbf{consumption
good} (shelter is a basic human need), a \textbf{durable asset} (houses last
50--100 years), a \textbf{financial investment} (households hold most of their
net worth in housing), and a \textbf{location-specific good} (a house in
Vancouver cannot substitute for a house in Saskatoon).

These characteristics give housing markets properties that distinguish them
from markets for ordinary goods:

\begin{itemize}
  \item \textbf{Highly inelastic short-run supply.} It takes 18--36 months
        to plan, permit, and build new housing. The existing stock of homes
        changes slowly. A sudden increase in demand therefore causes prices
        to rise sharply in the short run, even if the long-run supply
        response would eventually moderate prices.

  \item \textbf{Spatially fixed.} Housing supply cannot respond to demand
        by shipping inventory from low-demand to high-demand regions. Land
        in downtown Toronto is permanently scarce in a way that land in
        Kansas is not. This creates permanent scarcity rents in desirable
        locations that compound over time.

  \item \textbf{Heterogeneous.} No two homes are identical. Differences
        in location, size, vintage, condition, and neighbourhood make
        ``the housing market'' really an overlapping set of segmented
        submarkets that interact but behave differently.

  \item \textbf{Highly leveraged.} Most housing purchases are debt-financed.
        A household making a 10\% down payment is borrowing 90 cents of
        every dollar of housing value. This leverage amplifies both gains
        and losses, makes housing demand extremely sensitive to interest
        rates, and connects housing markets tightly to the credit system.

  \item \textbf{Regulated at every level.} Zoning bylaws (municipal),
        building codes (provincial), mortgage insurance rules (federal),
        rent regulations (provincial), environmental assessments
        (federal/provincial), and development charges (municipal) all
        constrain the supply of, and demand for, housing simultaneously.
\end{itemize}

\begin{defbox}{The Rental Market and the Ownership Market}
The housing market has two interconnected segments.

The \textbf{ownership market} prices houses as assets. The equilibrium
price reflects the expected present value of future housing services,
including the option value of residing in a particular location.

The \textbf{rental market} prices housing as a flow of services per unit
of time. Renters pay for the right to occupy a dwelling without acquiring
the asset.

The two markets are linked: when owning becomes more expensive relative
to renting (e.g.\ because mortgage rates rise), demand shifts from
ownership to rental, pushing rental prices up. Conversely, when rents
are very high, the yield on owning rental property rises, incentivising
investors to buy.

In Canada's current crisis, both markets have tightened simultaneously ---
a phenomenon explained by the coincidence of record immigration, low rental
vacancy, high mortgage rates, and constrained new supply.
\end{defbox}

%==========================================================
\section{Supply and Demand in Housing Markets}
\label{sec:housing_supply_demand}
%==========================================================

%----------------------------------------------------------
\subsection{The Demand for Housing}
%----------------------------------------------------------

The demand for housing is driven by:

\begin{itemize}
  \item \textbf{Population growth and household formation:} More people
        require more dwellings. Canada's record immigration flows
        (431{,}645 permanent residents in 2022; over 1 million people
        added to the population including non-permanent residents) have
        been the single largest driver of housing demand in the mid-2020s.

  \item \textbf{Income:} Higher household income increases both the
        ability and willingness to pay for housing. Housing is a
        \term{normal good} with an income elasticity of demand
        estimated between 0.5 and 1.0.

  \item \textbf{Interest rates:} The monthly cost of servicing a mortgage
        depends directly on the interest rate. Lower rates reduce the
        flow cost of owning a given quantity of housing, increasing demand.
        The 2020--2021 rate cuts to 0.25\% dramatically reduced mortgage
        payments and turbocharged demand.

  \item \textbf{Expected capital gains:} If buyers expect house prices
        to rise, the expected capital gain reduces the user cost of owning
        and shifts demand rightward. Extrapolative expectations can generate
        self-fulfilling price bubbles.

  \item \textbf{Investor demand:} A significant share of housing demand
        comes from investors who do not occupy the units. In some markets,
        investor demand competes directly with owner-occupiers and pushes
        prices above what the fundamental rental income stream would justify.
\end{itemize}

%----------------------------------------------------------
\subsection{The Supply of Housing}
%----------------------------------------------------------

The supply of housing in any given market is determined by:

\begin{itemize}
  \item \textbf{Existing stock:} The vast majority of housing supply in
        any year is the existing stock of homes. In Canada, approximately
        15 million private dwellings are occupied. New construction in a
        typical year (200{,}000--260{,}000 units) adds only 1.3--1.7\%
        to the stock.

  \item \textbf{New construction:} \term{Housing starts} measure the number
        of new residential units on which construction has begun. Their
        determinants include land prices, labour costs, material costs,
        developer financing costs, development charges, permitting timelines,
        and zoning regulations.

  \item \textbf{Conversion and densification:} Basement suites, laneway
        houses, and conversions of single-family homes to multiplexes
        can increase supply without requiring greenfield land.

  \item \textbf{Demolition and obsolescence:} A small fraction of the
        existing stock is demolished or becomes uninhabitable each year.
        Net new supply = New construction $-$ Demolitions.
\end{itemize}

%----------------------------------------------------------
\subsection{Supply Elasticity and Its Consequences}
%----------------------------------------------------------

The \term{housing supply elasticity} measures how responsive new construction
is to price increases:

\begin{equation}
  \varepsilon_s = \frac{\%\Delta Q_s^{\text{housing}}}{\%\Delta P_{\text{housing}}}
  \label{eq:supply_elasticity}
\end{equation}

Studies for major Canadian cities consistently find supply elasticities
well below 1.0 --- typically between 0.2 and 0.6 for the short run.
This means that a 10\% increase in house prices produces only a 2--6\%
increase in housing supply. The reasons for this inelasticity are explored
in Section~\ref{sec:supply_constraints}.

Figure~\ref{fig:housing_sd} illustrates the consequences of low supply
elasticity for housing price dynamics.

\begin{figure}[H]
\centering
\begin{tikzpicture}[
  axisline/.style={-Stealth, thick, color=gray!70},
  lbl/.style={font=\small}
]

%---- PANEL A: INELASTIC SUPPLY ----
\begin{scope}[xshift=0cm]
\node[font=\small\bfseries, align=center] at (3.2,5.8)
  {Panel A: Inelastic Supply (Canada)};
\draw[axisline] (0,0) -- (0,5.5) node[above, lbl] {$P$};
\draw[axisline] (0,0) -- (6.0,0) node[right, lbl] {$Q$};
% Inelastic supply
\draw[policyblue, thick] (2.5,0.3) -- (3.0,5.0)
  node[right, font=\footnotesize, color=policyblue] {$S_{\text{inelastic}}$};
% D1
\draw[unbcgreen, thick] (0.5,4.5) -- (5.5,0.3)
  node[right, font=\footnotesize, color=unbcgreen] {$D_1$};
% D2
\draw[casered, thick, dashed] (1.5,4.8) -- (6.0,0.7)
  node[right, font=\footnotesize, color=casered] {$D_2$};
\fill[unbcgreen] (2.68,2.2) circle (2.8pt);
\fill[casered]   (2.83,3.9) circle (2.8pt);
\draw[dashed, gray!50, thin] (0,2.2) -- (2.68,2.2);
\draw[dashed, gray!50, thin] (0,3.9) -- (2.83,3.9);
\draw[dashed, gray!50, thin] (2.68,0) -- (2.68,2.2);
\draw[dashed, gray!50, thin] (2.83,0) -- (2.83,3.9);
\node[left, font=\footnotesize] at (0,2.2) {$P_1$};
\node[left, font=\footnotesize] at (0,3.9) {$P_2^A$};
\node[below, font=\footnotesize] at (2.68,0) {$Q_1$};
\node[below, font=\footnotesize] at (2.83,0) {$Q_2^A$};
\draw[<->, casered!70, thick] (0.5,2.2) -- (0.5,3.9)
  node[midway, left, font=\footnotesize, color=casered] {$\Delta P$ large};
\draw[<->, policyblue!70, thick] (2.68,-0.5) -- (2.83,-0.5)
  node[midway, below, font=\footnotesize, color=policyblue] {$\Delta Q$ small};
\end{scope}

%---- PANEL B: ELASTIC SUPPLY ----
\begin{scope}[xshift=7.5cm]
\node[font=\small\bfseries, align=center] at (3.2,5.8)
  {Panel B: Elastic Supply (Policy Goal)};
\draw[axisline] (0,0) -- (0,5.5) node[above, lbl] {$P$};
\draw[axisline] (0,0) -- (6.0,0) node[right, lbl] {$Q$};
\draw[policyblue, thick] (0.5,1.8) -- (5.5,2.8)
  node[right, font=\footnotesize, color=policyblue] {$S_{\text{elastic}}$};
\draw[unbcgreen, thick] (0.5,4.5) -- (5.5,0.3)
  node[right, font=\footnotesize, color=unbcgreen] {$D_1$};
\draw[casered, thick, dashed] (1.5,4.8) -- (6.0,0.7)
  node[right, font=\footnotesize, color=casered] {$D_2$};
\fill[unbcgreen] (2.9,2.12) circle (2.8pt);
\fill[casered]   (4.15,2.38) circle (2.8pt);
\draw[dashed, gray!50, thin] (0,2.12) -- (2.9,2.12);
\draw[dashed, gray!50, thin] (0,2.38) -- (4.15,2.38);
\draw[dashed, gray!50, thin] (2.9,0) -- (2.9,2.12);
\draw[dashed, gray!50, thin] (4.15,0) -- (4.15,2.38);
\node[left, font=\footnotesize] at (0,2.12) {$P_1$};
\node[left, font=\footnotesize] at (0,2.38) {$P_2^B$};
\node[below, font=\footnotesize] at (2.9,0) {$Q_1$};
\node[below, font=\footnotesize] at (4.15,0) {$Q_2^B$};
\draw[<->, casered!70, thick] (0.5,2.12) -- (0.5,2.38)
  node[midway, left, font=\footnotesize, color=casered] {$\Delta P$ small};
\draw[<->, policyblue!70, thick] (2.9,-0.5) -- (4.15,-0.5)
  node[midway, below, font=\footnotesize, color=policyblue] {$\Delta Q$ large};
\end{scope}
\end{tikzpicture}
\caption{Housing Supply Elasticity and Price Dynamics.
\textbf{Panel A} shows Canada's current reality: the supply curve is nearly
vertical (inelastic) because zoning, land scarcity, and slow permitting
prevent new housing from responding quickly to demand. When demand shifts
from $D_1$ to $D_2$, prices rise sharply while quantity barely changes.
\textbf{Panel B} shows the policy goal: a more elastic supply curve absorbs
the same demand shift with a much smaller price increase and a much larger
quantity response. The fundamental insight is that housing affordability is
primarily a supply problem, and the supply curve's slope is the critical
policy parameter.}
\label{fig:housing_sd}
\end{figure}

%==========================================================
\section{The User Cost of Housing: Buy or Rent?}
\label{sec:user_cost}
%==========================================================

The decision to buy or rent a home is analysed using the \term{user cost
of housing} framework. The user cost is the annual economic cost of occupying
a home you own, expressed as a fraction of its market value:

\begin{equation}
  \text{UC} = (r + \delta + \tau - g) \times P
  \label{eq:user_cost}
\end{equation}

\noindent where:
\begin{itemize}
  \item $r$ = real mortgage interest rate
  \item $\delta$ = annual depreciation rate (approximately 1.5--2.5\%)
  \item $\tau$ = property tax rate (approximately 0.5--1.5\%)
  \item $g$ = expected real capital gains rate (annual price appreciation above inflation)
  \item $P$ = current price of the home
\end{itemize}

Renting is more economical when the annual rent is less than UC;
buying is more economical when annual rent exceeds UC.

\noindent\textbf{Worked Example 3.1 --- User Cost Under Different Conditions}

\noindent Consider a home priced at \$700{,}000 in Ottawa under three scenarios:

\begin{center}
\renewcommand{\arraystretch}{1.35}
\begin{tabular}{lrrr}
\toprule
\textbf{Parameter} & \textbf{2021 (Low Rate)} & \textbf{2023 (High Rate)}
                   & \textbf{2026 (Normal)} \\
\midrule
Mortgage rate $r$ (\%)         & 2.0 & 5.8 & 4.5 \\
Depreciation $\delta$ (\%)     & 2.0 & 2.0 & 2.0 \\
Property tax $\tau$ (\%)       & 1.0 & 1.0 & 1.0 \\
Expected real gain $g$ (\%)    & 4.0 & 0.0 & 1.5 \\
\midrule
\textbf{User cost rate (\%)}   & \textbf{1.0} & \textbf{8.8} & \textbf{6.0} \\
\textbf{Annual user cost (\$)} & \textbf{\$7{,}000} & \textbf{\$61{,}600}
                                & \textbf{\$42{,}000} \\
\textbf{Monthly equivalent}    & \textbf{\$583} & \textbf{\$5{,}133}
                                & \textbf{\$3{,}500} \\
\bottomrule
\end{tabular}
\end{center}

\noindent\textit{Interpretation:} In 2021, the monthly user cost of owning
a \$700{,}000 home was only \$583 --- far below any comparable rent.
This explains the demand surge of 2020--2022: owning was far cheaper than
renting in user cost terms. By 2023, the monthly user cost had risen to
\$5{,}133, making owning economically irrational relative to renting
in many markets.

\begin{thinkbox}
Why did house prices continue rising in 2021 even though prices were already
``unaffordably high'' by price-to-income measures? The user cost framework
answers this. When $g > r$ (expected capital gains exceed the interest rate),
the user cost can approach zero or become negative. The more expensive the
house, the more capital gain the owner earns --- so higher prices make buying
\textit{more} attractive, not less. This creates a feedback loop that sustains
price bubbles: rising prices increase expected future gains, reduce user cost,
attract more buyers, push prices higher, and so on. The loop breaks when
interest rates rise sharply (as in 2022), expected capital gains collapse,
and user costs suddenly jump.
\end{thinkbox}

%==========================================================
\section{Measuring Housing Affordability}
\label{sec:affordability_measurement}
%==========================================================

%----------------------------------------------------------
\subsection{The Price-to-Income Ratio}
%----------------------------------------------------------

The \term{price-to-income ratio (PTI)} is the most widely reported
affordability benchmark:

\begin{equation}
  \text{PTI} = \frac{\text{Median House Price}}{\text{Median Annual Household Income}}
  \label{eq:pti}
\end{equation}

A PTI of 3.0 is traditionally considered ``affordable.''
Above 5.0 is ``seriously unaffordable.'' Above 9.0 is ``severely unaffordable.''

\begin{figure}[H]
\centering
\begin{tikzpicture}
\begin{axis}[
  ybar,
  bar width=0.48cm,
  width=13cm, height=7cm,
  symbolic x coords={{Halifax},{Winnipeg},{Calgary},{Montreal},{Ottawa},
                     {CAN Avg},{Houston},{Berlin},{Toronto},{NYC},{Sydney},{Vancouver}},
  xtick=data,
  x tick label style={rotate=40, anchor=east, font=\small},
  ylabel={Price-to-Income Ratio, 2024},
  ylabel style={font=\small},
  ymin=0, ymax=15,
  ytick={0,2,4,6,8,10,12,14},
  yticklabel style={font=\small},
  ymajorgrids=true,
  grid style={dashed, gray!25},
  enlarge x limits=0.05,
  nodes near coords,
  nodes near coords style={font=\footnotesize},
  every axis plot/.append style={fill=policyblue!60, draw=policyblue!80}
]
\addplot coordinates {
  ({Halifax},    6.8)
  ({Winnipeg},   4.8)
  ({Calgary},    6.5)
  ({Montreal},   6.2)
  ({Ottawa},     7.2)
  ({CAN Avg},    8.1)
  ({Houston},    3.8)
  ({Berlin},     8.9)
  ({Toronto},   10.2)
  ({NYC},        9.8)
  ({Sydney},    13.3)
  ({Vancouver}, 12.8)
};
\addplot[unbcgreen, dashed, thick, forget plot]
  coordinates {({Halifax},3.0)({Vancouver},3.0)};
\addplot[casered, dashed, thick, forget plot]
  coordinates {({Halifax},5.0)({Vancouver},5.0)};
\node[font=\footnotesize, unbcgreen] at (axis cs:{Winnipeg},3.45)
  {``Affordable'' (3.0)};
\node[font=\footnotesize, casered] at (axis cs:{Winnipeg},5.45)
  {``Unaffordable'' (5.0)};
\end{axis}
\end{tikzpicture}
\caption{Price-to-Income Ratio, Selected Canadian and International Cities, 2024.
Vancouver (12.8) and Toronto (10.2) rank among the least affordable cities
globally. Only Winnipeg (4.8) of the major Canadian cities approaches the
historic ``affordable'' norm of 3.0 (green dashed line). The international
comparison is instructive: Houston (3.8) is a major North American city with
permissive land-use rules; Berlin (8.9) has high prices but also a large
social housing sector that insulates lower-income households from full market
exposure. Canada's affordability problem is real, severe, and primarily
concentrated in its two largest metropolitan areas.}
\label{fig:pti_comparison}
\source{Demographia International Housing Affordability Survey, 2024;
CMHC Housing Market Outlook, 2024.}
\end{figure}

%----------------------------------------------------------
\subsection{The Mortgage Payment Formula}
%----------------------------------------------------------

A second affordability measure is the \term{shelter cost burden}: the fraction
of income required to service housing costs. The conventional threshold is
30\% of gross income.

The standard Canadian mortgage payment formula is:

\begin{equation}
  M = P_0 \times \frac{r_m (1 + r_m)^n}{(1 + r_m)^n - 1}
  \label{eq:mortgage_payment}
\end{equation}

\noindent where $M$ is the monthly payment, $P_0$ is the principal (price
minus down payment), $r_m$ is the effective monthly rate derived from the
annual rate under the Canadian semi-annual compounding convention
($r_m = (1 + r_{\text{annual}}/2)^{1/6} - 1$), and $n$ is total monthly
payments (amortization years $\times 12$).

\begin{table}[H]
\centering
\caption{Monthly Mortgage Payment and Affordability, \$700{,}000 Home,
20\% Down (\$140{,}000), \$560{,}000 Principal, 25-Year Amortization}
\label{tab:mortgage_payments}
\renewcommand{\arraystretch}{1.35}
\begin{tabular}{lrrrr}
\toprule
\textbf{Rate (\%)} & \textbf{Monthly Payment} & \textbf{Annual Cost}
                   & \textbf{Income at 30\% Rule} & \textbf{Total Interest} \\
\midrule
2.0 & \$2{,}370 & \$28{,}440 & \$94{,}800  & \$150{,}800 \\
3.0 & \$2{,}646 & \$31{,}752 & \$105{,}840 & \$233{,}600 \\
4.0 & \$2{,}936 & \$35{,}232 & \$117{,}440 & \$320{,}800 \\
5.0 & \$3{,}263 & \$39{,}156 & \$130{,}520 & \$418{,}900 \\
5.8 & \$3{,}531 & \$42{,}372 & \$141{,}240 & \$499{,}300 \\
7.0 & \$3{,}921 & \$47{,}052 & \$156{,}840 & \$616{,}600 \\
\bottomrule
\end{tabular}
\source{Author calculations using Equation~\ref{eq:mortgage_payment}
under Canadian semi-annual compounding. ``Income at 30\% rule'' is the
annual gross income required for the mortgage payment to represent exactly
30\% of income. Total interest assumes full 25-year amortization.}
\end{table}

At 5.8\% --- the approximate 5-year fixed rate prevailing through 2023 ---
a household buying a \$700{,}000 home needed gross income of at least
\$141{,}240 per year to meet the 30\% threshold. In Vancouver, where the
benchmark price exceeded \$1.1 million through 2024, the required income
at 5.8\% exceeded \$225{,}000.

%----------------------------------------------------------
\subsection{The CMHC Affordability Index}
%----------------------------------------------------------

CMHC publishes a \term{Housing Affordability Index} measuring mortgage
payment, property taxes, and utilities as a percentage of median household
income for a median-priced home.

\begin{table}[H]
\centering
\caption{CMHC Housing Affordability Index, Selected Cities, 2015--2024}
\label{tab:affordability_index}
\renewcommand{\arraystretch}{1.30}
\begin{tabular}{lrrrrrr}
\toprule
\textbf{City} & \textbf{2015} & \textbf{2018} & \textbf{2020}
              & \textbf{2022} & \textbf{2023} & \textbf{2024} \\
\midrule
Vancouver  & 80.2 & 87.6 & 61.4 & 93.8 & 90.2 & 92.1 \\
Toronto    & 56.8 & 74.2 & 67.3 & 89.4 & 82.7 & 75.3 \\
Ottawa     & 38.2 & 42.1 & 44.5 & 62.3 & 58.7 & 48.2 \\
Calgary    & 35.6 & 32.1 & 31.8 & 42.5 & 45.1 & 38.4 \\
Montr\'{e}al & 34.0 & 38.4 & 41.2 & 55.2 & 49.8 & 42.3 \\
Halifax    & 26.1 & 28.3 & 30.2 & 49.3 & 46.5 & 44.2 \\
Winnipeg   & 25.8 & 26.5 & 27.3 & 33.2 & 31.5 & 29.4 \\
\textbf{National avg} & \textbf{40.2} & \textbf{50.1} & \textbf{46.8}
                      & \textbf{65.3} & \textbf{58.4} & \textbf{53.2} \\
\bottomrule
\end{tabular}
\source{Canada Mortgage and Housing Corporation, Housing Market Outlook,
2024. Index = (mortgage payment + taxes + utilities) as \% of median
household income. Values above 30\% indicate unaffordable housing.}
\end{table}

Vancouver's 2024 index of 92.1\% means the median Vancouver household
buying the median-priced home would spend 92\% of gross income on housing
costs --- mathematically impossible for most families. The national average
peaked at 65.3\% in 2022 and has partially improved since, but remains
far above the pre-crisis levels of the early 2010s.

%==========================================================
\section{The Historical Price Record}
\label{sec:house_price_history}
%==========================================================

\begin{figure}[H]
\centering
\begin{tikzpicture}
\begin{axis}[
  width=14cm, height=7.5cm,
  xlabel={Year},
  ylabel={MLS Benchmark Price (\$CAD thousands)},
  xlabel style={font=\small},
  ylabel style={font=\small},
  xmin=1999, xmax=2027,
  ymin=0, ymax=1600,
  xtick={2000,2004,2008,2012,2016,2020,2024},
  xticklabel style={font=\small},
  ytick={0,200,400,600,800,1000,1200,1400,1600},
  yticklabel style={font=\small},
  ymajorgrids=true,
  grid style={dashed, gray!20},
  legend pos=north west,
  legend style={font=\small},
  clip=false
]
\addplot[gray!70, thick] coordinates {
  (2000,163)(2004,233)(2008,302)(2010,340)(2012,365)(2014,400)
  (2016,480)(2018,488)(2020,530)(2021,716)(2022,796)
  (2023,650)(2024,700)(2025,710)(2026,720)
};
\addlegendentry{Canada national}
\addplot[policyblue, thick] coordinates {
  (2000,270)(2004,370)(2008,520)(2010,598)(2012,683)(2014,720)
  (2016,1050)(2018,870)(2020,1004)(2021,1200)(2022,1370)
  (2023,1170)(2024,1250)(2025,1270)(2026,1290)
};
\addlegendentry{Vancouver}
\addplot[casered, thick, dashed] coordinates {
  (2000,250)(2004,310)(2008,362)(2010,431)(2012,497)(2014,566)
  (2016,730)(2018,736)(2020,930)(2021,1163)(2022,1340)
  (2023,1100)(2024,1148)(2025,1160)(2026,1175)
};
\addlegendentry{Toronto}
\addplot[dataorange, thick, dotted] coordinates {
  (2000,160)(2006,340)(2008,380)(2010,365)(2016,420)
  (2020,420)(2021,475)(2022,560)(2023,570)(2024,600)(2025,615)(2026,625)
};
\addlegendentry{Calgary}
\node[font=\footnotesize, align=center, color=gray] at (axis cs:2008.5,160)
  {GFC};
\node[font=\footnotesize, color=unbcgreen, align=left] at (axis cs:2020.2,420)
  {COVID\\low rates};
\node[font=\footnotesize, color=casered, align=right] at (axis cs:2022.5,1450)
  {2022 peak};
\end{axis}
\end{tikzpicture}
\caption{MLS Composite Benchmark House Prices, Canada and Selected Cities,
2000--2026. Canadian house prices have risen more than four-fold in nominal
terms since 2000; Vancouver and Toronto five- to six-fold. Key episodes:
the 2008 Global Financial Crisis caused only a brief pause; the 2016 BC
Foreign Buyer Tax caused a temporary Vancouver correction; COVID-era
ultra-low rates (2020--2021) triggered the most explosive price surge in
history; and the 2022--2023 rate hikes caused the first significant correction,
though prices remain far above pre-pandemic levels. Calgary's flatter profile
reflects oil price cycle sensitivity.}
\label{fig:house_prices}
\source{Canadian Real Estate Association (CREA), MLS Home Price Index.
2025--2026 estimated.}
\end{figure}

%==========================================================
\section{Supply-Side Constraints on Canadian Housing}
\label{sec:supply_constraints}
%==========================================================

The central analytical question in housing economics is whether price
increases are primarily driven by demand or supply. The academic evidence
increasingly points to \textbf{supply constraints as the primary long-run
driver}. The same demand pressures exist in many countries without producing
the same price outcomes --- what differs is how quickly supply responds.

%----------------------------------------------------------
\subsection{Zoning and Land Use Regulation}
%----------------------------------------------------------

The most important supply-side constraint is \term{zoning} --- the municipal
regulation of how land may be used. In most Canadian cities, the majority
of residential land is zoned exclusively for single-family detached housing.
This ``Euclidean'' zoning model prevents the construction of apartments,
townhouses, or duplexes on most residential land regardless of market demand.

The economic consequences are profound:

\begin{itemize}
  \item \textbf{Land value inflation:} When zoning prevents densification,
        the value of a single-family lot captures the entire surplus from
        excluding all alternative, potentially more productive uses.

  \item \textbf{Reduced supply elasticity:} Zoning directly reduces
        $\varepsilon_s$ (Equation~\ref{eq:supply_elasticity}). Developers
        who want to build apartments cannot do so without lengthy, costly,
        and uncertain rezoning processes.

  \item \textbf{Parking minimums and setback requirements:} Many municipalities
        require 1--2 parking spaces per unit, minimum lot sizes, setbacks,
        and maximum lot coverage. Each constraint reduces the number of
        units buildable per parcel.
\end{itemize}

\begin{canexample}{British Columbia's 2023 Zoning Reform}
In December 2023, BC passed legislation requiring all municipalities in
Metro Vancouver and most municipalities with populations over 5{,}000 to
permit ``small scale multi-unit housing'' --- duplexes, triplexes,
fourplexes, and sixplexes --- on lots previously zoned exclusively for
single-family homes. This \term{upzoning} legislation was the most
significant provincial override of municipal zoning in Canadian history.

Economists estimate the reform could permit hundreds of thousands of
additional units over the next decade if complemented by appropriate
infrastructure investment and streamlined permitting. Early evidence
from 2025 suggests modest but positive effects on building permit
applications in affected areas.
\end{canexample}

%----------------------------------------------------------
\subsection{Development Costs and Construction Capacity}
%----------------------------------------------------------

Even where zoning permits higher density, additional cost factors constrain supply:

\begin{itemize}
  \item \textbf{Development charges:} In Toronto, development charges
        on a new condo unit reached approximately \$150{,}000--\$200{,}000
        in 2023 --- representing 15--25\% of the unit's total cost, ultimately
        borne by buyers through higher prices.

  \item \textbf{Construction labour shortages:} Canada's construction sector
        faces significant skilled labour shortages --- carpenters, electricians,
        plumbers, ironworkers. Construction wages have risen faster than
        overall wages since 2020.

  \item \textbf{Permitting timelines:} The time from a development application
        to a building permit in major Canadian cities averages 12--36 months.
        Each month of delay costs developers carrying charges on their land.

  \item \textbf{High-rise construction costs:} Tall concrete and steel
        construction costs \$350--\$600 per square foot in Canada, compared
        to \$150--\$250 for wood-frame low-rise. The economics only work
        at very high price levels.
\end{itemize}

%----------------------------------------------------------
\subsection{The Supply-Demand Gap: Housing Starts vs. Population Growth}
%----------------------------------------------------------

\begin{figure}[H]
\centering
\begin{tikzpicture}
\begin{axis}[
  name=ax1,
  width=14cm, height=7cm,
  xlabel={Year},
  ylabel={Housing Starts (thousands of units)},
  ylabel style={font=\small, color=policyblue},
  xlabel style={font=\small},
  xmin=1999, xmax=2025.5,
  ymin=100, ymax=320,
  xtick={2000,2004,2008,2012,2016,2020,2024},
  xticklabel style={font=\small},
  ytick={100,150,200,250,300},
  yticklabel style={font=\small, color=policyblue},
  axis y line*=left,
  axis x line*=bottom,
  ymajorgrids=true,
  grid style={dashed, gray!20},
  clip=false,
  legend style={font=\small, at={(0.02,0.98)}, anchor=north west}
]
\addplot[policyblue, thick, mark=*, mark size=1.5pt] coordinates {
  (2000,151)(2002,205)(2004,233)(2006,228)(2008,211)(2009,149)
  (2010,190)(2012,215)(2014,190)(2016,198)(2018,212)(2020,218)
  (2021,271)(2022,263)(2023,239)(2024,245)
};
\addlegendentry{Housing Starts (left)}
\node[font=\footnotesize, policyblue, align=center] at (axis cs:2022.5,295)
  {Need $\approx$450K+ units/yr};
\draw[->, policyblue!60, thin] (axis cs:2022,283) -- (axis cs:2022,265);
\end{axis}
\begin{axis}[
  width=14cm, height=7cm,
  ylabel={Annual Population Growth (thousands)},
  ylabel style={font=\small, color=casered},
  xmin=1999, xmax=2025.5,
  ymin=200, ymax=1400,
  ytick={200,400,600,800,1000,1200,1400},
  yticklabel style={font=\small, color=casered},
  axis y line*=right,
  axis x line=none,
  legend style={font=\small, at={(0.98,0.98)}, anchor=north east},
  clip=false
]
\addplot[casered, thick, dashed, mark=square*, mark size=1.5pt] coordinates {
  (2000,362)(2002,339)(2004,337)(2006,355)(2008,393)(2010,338)
  (2012,357)(2014,388)(2016,339)(2018,528)(2020,329)(2021,507)
  (2022,1053)(2023,1159)(2024,1200)
};
\addlegendentry{Population Growth (right)}
\end{axis}
\end{tikzpicture}
\caption{Housing Starts vs.\ Annual Population Growth, Canada, 2000--2024.
Until approximately 2017, housing starts and population growth moved in rough
proportion. From 2022 onwards, Canada's record immigration levels produced
population growth of over 1 million persons per year --- requiring approximately
400{,}000--480{,}000 new housing units annually (at 2.5 persons per household)
--- while actual construction remained at 240{,}000--265{,}000 starts.
This structural supply shortfall is the primary driver of the post-2021
rental vacancy crisis and rent inflation.}
\label{fig:starts_vs_population}
\source{Statistics Canada, Tables 34-10-0035-01 (housing starts) and
17-10-0008-01 (population growth estimates).}
\end{figure}

%==========================================================
\section{Demand-Side Forces: Immigration, Finance, and Demographics}
\label{sec:demand_side}
%==========================================================

While supply constraints are the primary structural explanation for high
prices, several demand-side forces have amplified and accelerated the
affordability crisis.

%----------------------------------------------------------
\subsection{Immigration and Population Growth}
%----------------------------------------------------------

Canada's record immigration of 431{,}645 permanent residents in 2022 (and
approximately 1 million total population additions including non-permanent
residents) represents the single largest year-over-year demand shock to
housing in Canadian history. The housing consequences are significant because:

\begin{enumerate}
  \item Newcomers must immediately find accommodation.
  \item Newcomers concentrate geographically: 57\% of 2022 permanent
        residents settled in Ontario or British Columbia --- precisely the
        markets with lowest vacancy rates.
  \item Non-permanent residents (international students, temporary foreign
        workers) are disproportionately renters, adding pressure to already
        strained rental markets.
\end{enumerate}

Immigration does not make immigration bad policy --- Chapter~\ref{ch:immigration}
presents the strong evidence for its economic benefits. But it does require
that housing supply policy be coordinated with immigration policy. Canada
set record immigration targets without simultaneously implementing the
supply-side reforms needed to house the newcomers. This policy coordination
failure contributed significantly to the 2022--2024 rental crisis.

%----------------------------------------------------------
\subsection{The Interest Rate Channel}
%----------------------------------------------------------

As established in Chapter~\ref{ch:inflation}, the Bank of Canada's decision
to hold the overnight rate at 0.25\% through 2020--2021 created deeply
negative real interest rates. Through the user cost mechanism
(Equation~\ref{eq:user_cost}), this dramatically reduced the annual cost
of owning housing:

\begin{itemize}
  \item At $r = 2\%$ and $g = 5\%$ expected capital gains, the user cost
        rate was approximately $-0.5\%$ to $+1.0\%$.
  \item This meant the annual cost of owning a \$1{,}000{,}000 home was
        \$0--\$10{,}000, while comparable rental costs ran \$36{,}000--\$48{,}000
        per year in Vancouver or Toronto.
  \item The rent-vs-buy calculation was decisively in favour of buying,
        generating enormous ownership demand.
\end{itemize}

\begin{figure}[H]
\centering
\begin{tikzpicture}
\begin{axis}[
  width=14cm, height=6.5cm,
  xlabel={Year},
  ylabel={Rate (\%)},
  xlabel style={font=\small},
  ylabel style={font=\small},
  xmin=2014.5, xmax=2026.5,
  ymin=0, ymax=8,
  xtick={2015,2016,2017,2018,2019,2020,2021,2022,2023,2024,2025,2026},
  xticklabel style={rotate=35, anchor=east, font=\small},
  ytick={0,1,2,3,4,5,6,7,8},
  yticklabel style={font=\small},
  ymajorgrids=true,
  grid style={dashed, gray!20},
  legend pos=north west,
  legend style={font=\small},
  clip=false
]
\addplot[policyblue, thick, mark=*, mark size=1.5pt] coordinates {
  (2015,4.6)(2016,4.6)(2017,4.6)(2018,5.3)(2019,3.5)
  (2020,2.2)(2021,2.0)(2022,5.1)(2023,6.1)(2024,5.2)(2025,4.6)(2026,4.3)
};
\addlegendentry{5-year fixed mortgage rate}
\addplot[casered, thick, dashed, mark=square*, mark size=1.5pt] coordinates {
  (2015,2.7)(2016,2.7)(2017,3.2)(2018,3.9)(2019,3.5)
  (2020,1.5)(2021,1.4)(2022,4.8)(2023,6.9)(2024,5.4)(2025,4.0)(2026,3.6)
};
\addlegendentry{Variable rate (approx.)}
\addplot[unbcgreen!70, dotted, thick] coordinates {
  (2015,0.5)(2016,0.5)(2017,1.0)(2018,1.75)(2019,1.75)
  (2020,0.25)(2021,0.25)(2022,4.25)(2023,5.0)(2024,3.25)(2025,3.0)(2026,2.75)
};
\addlegendentry{Bank of Canada overnight rate}
\node[font=\footnotesize, policyblue, align=center] at (axis cs:2020.5,0.8)
  {Historic\\rate lows};
\node[font=\footnotesize, casered] at (axis cs:2023,7.3)
  {Post-hike peak};
\end{axis}
\end{tikzpicture}
\caption{Mortgage Rates vs.\ Bank of Canada Overnight Rate, 2015--2026.
The 5-year fixed mortgage rate fell to a historic low of approximately 2.0\%
in 2021, driving a massive surge in housing demand. The rate-hike cycle
drove the 5-year fixed rate above 6\% by mid-2023 --- more than tripling
the cost of new mortgage borrowing. Variable-rate holders saw their payments
rise almost instantly, tracking the overnight rate with minimal lag.}
\label{fig:mortgage_rates}
\source{Bank of Canada; Statistics Canada. 2025--2026 estimated.}
\end{figure}

%----------------------------------------------------------
\subsection{Financialization of Housing}
%----------------------------------------------------------

A third demand-side factor is the \term{financialization of housing} ---
the increasing treatment of residential real estate as a financial asset class:

\begin{itemize}
  \item \textbf{Investor purchasing:} CMHC estimates that investors
        accounted for 20--30\% of condominium purchases in Toronto and
        Vancouver in recent years, competing with first-time buyers.

  \item \textbf{Short-term rentals:} Airbnb and similar platforms converted
        approximately 31{,}000 long-term rental units into tourist
        accommodation in Canada's 10 largest cities.

  \item \textbf{Vacancy and speculation:} Some buyers hold units vacant
        as stores of value. Several Canadian cities have introduced vacancy
        taxes (Vancouver Empty Homes Tax, BC Speculation and Vacancy Tax)
        that address the problem directly and generate data to measure its scale.
\end{itemize}

%==========================================================
\section{The Rental Crisis}
\label{sec:rental_crisis}
%==========================================================

As buying became impossible for an increasing share of Canadians, demand
for rental housing surged --- but rental supply grew slowly, pushing vacancy
rates to historic lows and rents to historic highs.

\begin{figure}[H]
\centering
\begin{tikzpicture}
\begin{axis}[
  width=14cm, height=6.5cm,
  xlabel={Year},
  ylabel={Rental Vacancy Rate (\%)},
  xlabel style={font=\small},
  ylabel style={font=\small},
  xmin=2013.5, xmax=2025,
  ymin=0, ymax=7,
  xtick={2014,2015,2016,2017,2018,2019,2020,2021,2022,2023,2024},
  xticklabel style={font=\small},
  ytick={0,1,2,3,4,5,6,7},
  yticklabel style={font=\small},
  ymajorgrids=true,
  grid style={dashed, gray!20},
  legend pos=north east,
  legend style={font=\small},
  clip=false
]
\addplot[gray!50, dotted, thick, forget plot]
  coordinates {(2013.5,3.0)(2025,3.0)};
\node[right, font=\footnotesize, color=gray!70]
  at (axis cs:2024.7,3.0) {};
\node[font=\footnotesize, align=left, color=gray!70]
  at (axis cs:2015.5,3.3) {Healthy (3\%)};
\addplot[policyblue, thick, mark=*, mark size=1.5pt] coordinates {
  (2014,0.6)(2015,0.6)(2016,0.7)(2017,0.9)(2018,1.0)(2019,1.1)
  (2020,2.5)(2021,1.7)(2022,0.9)(2023,0.8)(2024,0.9)
};
\addlegendentry{Vancouver}
\addplot[casered, thick, dashed, mark=square*, mark size=1.5pt] coordinates {
  (2014,1.2)(2015,1.3)(2016,1.1)(2017,1.0)(2018,1.1)(2019,1.1)
  (2020,3.4)(2021,4.3)(2022,1.7)(2023,1.5)(2024,1.5)
};
\addlegendentry{Toronto}
\addplot[dataorange, thick, dotted, mark=triangle*, mark size=1.5pt] coordinates {
  (2014,4.7)(2015,5.3)(2016,7.0)(2017,6.4)(2018,3.8)(2019,3.6)
  (2020,5.2)(2021,4.5)(2022,2.7)(2023,1.4)(2024,1.4)
};
\addlegendentry{Calgary}
\addplot[unbcgreen!80, thick, mark=diamond*, mark size=1.5pt] coordinates {
  (2014,2.7)(2015,2.8)(2016,3.0)(2017,3.0)(2018,2.4)(2019,2.2)
  (2020,3.0)(2021,3.1)(2022,1.9)(2023,1.5)(2024,1.5)
};
\addlegendentry{National average}
\end{axis}
\end{tikzpicture}
\caption{Rental Market Vacancy Rates, Selected Canadian Cities, 2014--2024.
The ``healthy'' vacancy rate of 3\% provides a benchmark for a functioning
rental market. Vancouver has been below 3\% for all but one year since at
least 2014. The COVID-19 pandemic temporarily elevated Toronto's vacancy to
4.3\% in 2021, but this relief was transitory: by 2022, all major markets
had returned to near-zero vacancy. Calgary's trajectory is distinctive ---
the oil price collapse of 2015--2016 created a temporary glut, but the
city subsequently tightened as energy employment recovered and immigration
accelerated post-2021.}
\label{fig:vacancy_rates}
\source{Canada Mortgage and Housing Corporation, Rental Market Report, 2024.}
\end{figure}

When vacancy rates fall below 2\%, the rental market becomes a landlords'
market. Average asking rents for 2-bedroom apartments reached \$3{,}200/month
in Vancouver and \$2{,}750/month in Toronto in late 2024. For a household
earning the provincial median income, these rents consume 40--55\% of
after-tax income.

\begin{figure}[H]
\centering
\begin{tikzpicture}
\begin{axis}[
  ybar,
  bar width=0.40cm,
  width=13cm, height=6.5cm,
  symbolic x coords={{Vancouver},{Toronto},{Ottawa},{Halifax},
                     {Calgary},{Montreal},{Winnipeg},{National Avg}},
  xtick=data,
  x tick label style={rotate=38, anchor=east, font=\small},
  ylabel={Shelter Cost Burden (\% of After-Tax Income)},
  ylabel style={font=\small},
  ymin=0, ymax=72,
  ytick={0,10,20,30,40,50,60,70},
  yticklabel style={font=\small},
  ymajorgrids=true,
  grid style={dashed, gray!25},
  enlarge x limits=0.08,
  legend pos=north east,
  legend style={font=\small},
  clip=false
]
\addplot[fill=casered!60, draw=casered!80] coordinates {
  ({Vancouver},67)({Toronto},58)({Ottawa},46)({Halifax},48)
  ({Calgary},36)({Montreal},42)({Winnipeg},28)({National Avg},50)
};
\addlegendentry{Renter, bottom income quintile}
\addplot[fill=policyblue!50, draw=policyblue!70] coordinates {
  ({Vancouver},48)({Toronto},42)({Ottawa},32)({Halifax},35)
  ({Calgary},24)({Montreal},28)({Winnipeg},22)({National Avg},36)
};
\addlegendentry{Renter, median income}
\addplot[fill=unbcgreen!50, draw=unbcgreen!70] coordinates {
  ({Vancouver},24)({Toronto},21)({Ottawa},17)({Halifax},16)
  ({Calgary},14)({Montreal},15)({Winnipeg},12)({National Avg},18)
};
\addlegendentry{Owner, median income (existing mortgage)}
\addplot[gray!60, dashed, thick, forget plot]
  coordinates {({Vancouver},30)({National Avg},30)};
\node[font=\footnotesize, color=gray!70, align=right]
  at (axis cs:{Montreal},31.5) {30\% threshold};
\end{axis}
\end{tikzpicture}
\caption{Shelter Cost Burden by Tenure and Income Group, Major Canadian Cities, 2024.
Shelter cost burden is annual shelter costs as a percentage of annual after-tax
income. The 30\% threshold (grey dashed line) is the conventional affordability
standard. Low-income renters (bottom quintile) spend over two-thirds of
after-tax income on shelter in Vancouver. Even median-income renters exceed
the 30\% threshold in Vancouver and Toronto. By contrast, median-income
homeowners with existing mortgages --- often locked in at pre-2022 rates ---
face a far smaller shelter cost share, illustrating the growing economic divide
between housing ``haves'' (established owners) and ``have-nots''
(renters and prospective first-time buyers).}
\label{fig:shelter_burden}
\source{Statistics Canada, Survey of Household Spending, 2023--2024;
CMHC Rental Market Report, 2024. Author calculations.}
\end{figure}

%==========================================================
\section{Case Studies}
\label{sec:ch3_casestudies}
%==========================================================

\begin{casestudy}{Vancouver: From Affordable to Global Outlier, 1980--2026}

\textbf{Background:} In 1980, the average house price in Metro Vancouver
was approximately \$120{,}000. Average household income was approximately
\$28{,}000. The price-to-income ratio was 4.3. By 2022, the average price
had reached \$1.37 million while average household income was approximately
\$107{,}000. The PTI had reached 12.8 --- among the highest in the world.

\textbf{What Changed?}

\begin{itemize}
  \item \textbf{Geographic constraints:} Metro Vancouver is squeezed between
        mountains, the US border, the ocean, and the Agricultural Land Reserve.
        The ALR, established in 1973 to protect farmland, constrains the urban
        growth boundary and reduces the supply of developable land.

  \item \textbf{Low supply elasticity:} Single-family zoning covers
        approximately 70\% of Metro Vancouver's residential land. Height limits
        restrict building up; the ALR restricts building out.

  \item \textbf{Capital flows and international demand:} Beginning in the
        late 1980s, Metro Vancouver became a destination for substantial capital
        inflows from Asia --- Hong Kong in the 1990s, mainland China in the
        2000s--2010s. The correlation between capital account openness and
        Vancouver price acceleration is visible in Figure~\ref{fig:house_prices}.

  \item \textbf{Interest rate cycles:} Vancouver prices are highly sensitive
        to real interest rates via the user cost mechanism, amplifying each
        low-rate period.
\end{itemize}

\textbf{Policy Responses:} BC and Vancouver have deployed virtually every
policy tool: 15--20\% Foreign Buyer Tax (2016), BC Speculation and Vacancy
Tax (2018), mortgage stress test (federal, 2018), transit-oriented development
zoning (2023), Small Scale Multi-Unit Housing legislation (2023), and
short-term rental restrictions (2024).

\textbf{Outcome:} Despite these interventions, Vancouver remains among the
three least affordable cities in the world. The 2022--2023 rate-hike correction
reduced prices approximately 14\%, but prices remained 25--30\% above 2019
levels. The fundamental supply constraint has not been resolved. The lesson:
demand-side interventions alone cannot solve a problem rooted in supply
inelasticity and decades of underbuilding.
\end{casestudy}

%----------------------------------------------------------
\begin{casestudy}{Canada vs. Germany: What Different Housing Policies Produce}

\textbf{Background:} Germany provides one of the most instructive international
comparisons. Both countries are wealthy, democratic, urbanised societies with
significant immigrant populations. Yet their housing outcomes differ dramatically.

\textbf{German Housing Outcomes (2024):}
\begin{itemize}
  \item Berlin price-to-income ratio: approximately 8.9
  \item National homeownership rate: approximately 50\% (the lowest in the EU)
  \item Long-term rental tenure: approximately 60\% of households rent,
        many in the same unit for 10--20+ years
  \item Purpose-built rental stock: approximately 55\% of total dwellings
\end{itemize}

\textbf{What Germany Does Differently:}
\begin{enumerate}
  \item \textbf{Tenant protection:} Strong legal protection against eviction
        allows households to plan long-term without fear of displacement,
        reducing the ``premium'' on homeownership as the only path to
        residential stability.

  \item \textbf{Large social housing sector:} Germany's social housing stock
        is 5--6 times larger (as a share of dwellings) than Canada's.

  \item \textbf{Cooperative housing:} Cooperatives own approximately 7\%
        of German housing stock; members pay below-market cost rents and
        have tenure security without speculative price exposure.

  \item \textbf{Transfer taxes on resale:} Germany charges 3.5--6.5\%
        land transfer taxes on property sales, reducing speculative trading.
\end{enumerate}

\textbf{Implication for Canada:} Canada's high homeownership aspiration and
underdeveloped rental sector are not inevitable features of a wealthy economy ---
they are the product of specific policy choices. The German comparison suggests
that sustained, multi-decade commitments to building non-market housing and
protecting tenant tenure can produce meaningfully better affordability outcomes.
\end{casestudy}

%----------------------------------------------------------
\begin{casestudy}{Rent Control: The Economics of a Popular Policy}

\textbf{Background:} Several Canadian provinces impose rent control ---
limits on rent increases for existing tenants. Ontario's rent increase
guideline caps annual increases at the Ontario CPI rate for most units.
BC has similar provisions.

\textbf{The Economic Case Against Rent Control:}
A binding rent ceiling $\bar{R}$ set below market rent $R^*$ creates:
\begin{itemize}
  \item \textbf{Excess demand:} At $\bar{R} < R^*$, quantity demanded
        exceeds quantity supplied. Non-price allocation mechanisms
        (personal connections, waiting lists) replace the price mechanism.
  \item \textbf{Reduced supply:} If developers anticipate future capped
        rents, expected returns on purpose-built rental projects fall,
        reducing new investment. Ontario explicitly exempts units built
        after November 2018 from rent control to address this.
  \item \textbf{Misallocation:} Tenants have strong incentives to remain
        in controlled units even when their circumstances change, reducing
        labour market mobility.
  \item \textbf{Reduced maintenance:} Landlords facing controlled rents
        may reduce maintenance spending to cut costs.
\end{itemize}

\textbf{The Economic Case For Rent Control:}
When vacancy rates are near zero, landlords have market power and can
extract rents above competitive levels. Rent control can correct a market
power problem, provide tenure security, and prevent displacement of
established low-income communities.

\textbf{The Evidence:}
A 2019 Stanford study found that San Francisco's rent control reduced
rental housing supply by 15\% as landlords converted units to condos
or redeveloped to escape controls. Ontario's pattern of low purpose-built
rental construction throughout the 1990s--2000s (when rent control was
broadly applied), followed by a surge in purpose-built starts after the
2018 exemption, is consistent with the supply-reduction effect.

\textbf{The Policy Synthesis:} Rent control protects \textit{existing}
tenants at the expense of \textit{future} tenants and new supply. Where
tenant protection and tenure security are the goals, alternatives such
as vacancy controls, tenant right-of-first-refusal, and anti-eviction
legislation may achieve those goals with less supply distortion. Where
affordability is the goal, housing allowances and social housing expand
supply rather than constrain it.
\end{casestudy}

%==========================================================
\section{Social Housing: Canada in International Context}
\label{sec:social_housing}
%==========================================================

\begin{figure}[H]
\centering
\begin{tikzpicture}
\begin{axis}[
  ybar,
  bar width=0.50cm,
  width=13cm, height=6.5cm,
  symbolic x coords={{USA},{Canada},{Australia},{Germany},
                     {France},{UK},{Denmark},{Netherlands}},
  xtick=data,
  x tick label style={font=\small},
  ylabel={Social Housing Stock (\% of Total Dwellings)},
  ylabel style={font=\small},
  ymin=0, ymax=35,
  ytick={0,5,10,15,20,25,30},
  yticklabel style={font=\small},
  ymajorgrids=true,
  grid style={dashed, gray!25},
  enlarge x limits=0.10,
  nodes near coords,
  nodes near coords style={font=\small},
  every axis plot/.append style={fill=exampleteal!60, draw=exampleteal!80}
]
\addplot coordinates {
  ({USA},     1.0)
  ({Canada},  3.5)
  ({Australia},4.8)
  ({Germany}, 12.5)
  ({France},  17.0)
  ({UK},      17.5)
  ({Denmark}, 21.0)
  ({Netherlands},30.5)
};
\end{axis}
\end{tikzpicture}
\caption{Social Housing as a Share of Total Dwellings, Selected Countries.
Canada's social housing stock (approximately 3.5\%) is among the smallest
of comparable wealthy nations. The Netherlands (30.5\%) and Denmark (21\%)
maintain large non-market housing sectors that insulate a significant share
of their populations from private market volatility. The OECD recommends
that countries targeting housing affordability increase non-market supply
to at least 10--15\% of total dwellings. Achieving this in Canada would
require building approximately 750{,}000--1{,}000{,}000 additional social
and affordable units --- a multi-decade project requiring sustained political
commitment and federal investment.}
\label{fig:social_housing}
\source{OECD Affordable Housing Database, 2023; CMHC National Housing Strategy, 2024.}
\end{figure}

%==========================================================
\section{Federal Housing Policy: What Has Been Tried}
\label{sec:housing_policy}
%==========================================================

\begin{table}[H]
\centering
\caption{Major Federal Housing Policy Measures, 2015--2026}
\label{tab:housing_policies}
\renewcommand{\arraystretch}{1.30}
\begin{tabular}{p{3.3cm}p{1.0cm}p{2.0cm}p{4.6cm}p{2.6cm}}
\toprule
\textbf{Policy} & \textbf{Year} & \textbf{Type} & \textbf{Mechanism}
                & \textbf{Assessment} \\
\midrule
Mortgage stress test
  & 2018 & Demand (regulatory)
  & Borrowers qualify at contract rate + 2\%, reducing borrowing capacity
  & Effective at limiting leverage; worsens access for marginal buyers \\
\addlinespace
First Home Savings Account
  & 2023 & Demand (tax)
  & \$40K tax-free savings for first-time buyers; contributions tax-deductible
  & Modest benefit; may capitalise into higher prices \\
\addlinespace
FTHBI (Shared Equity)
  & 2019 & Demand (subsidy)
  & CMHC takes 5--10\% equity stake; disbanded 2024
  & Limited uptake; demand subsidy in supply-constrained market \\
\addlinespace
National Housing Co-Investment Fund
  & 2018 & Supply (subsidy)
  & Low-interest loans for affordable rental construction
  & Positive but underfunded; units produced below target \\
\addlinespace
GST Removal on Purpose-Built Rentals
  & 2023 & Supply (tax)
  & Eliminates 5\% GST on new purpose-built rental construction
  & Positive; reduces effective construction cost by 3--5\% \\
\addlinespace
Housing Accelerator Fund
  & 2023 & Supply (incentive)
  & \$4B fund to municipalities contingent on zoning reform
  & Promising; targets root constraint; early results positive \\
\addlinespace
Foreign Buyer Ban
  & 2023 & Demand (regulatory)
  & Prohibits non-residents from buying residential property
  & Marginal effect; foreign buyers were small share of transactions \\
\bottomrule
\end{tabular}
\source{Department of Finance Canada; CMHC; Parliamentary Budget Officer,
Housing Affordability Assessment, 2024.}
\end{table}

The dominant pattern in Table~\ref{tab:housing_policies} is instructive:
most federal housing policies through 2022 were \textbf{demand-side}
measures --- tax credits, subsidies, and regulatory restrictions. With a
near-vertical supply curve (Panel A of Figure~\ref{fig:housing_sd}),
demand-side interventions cannot improve affordability. A demand subsidy
to a fixed-supply market raises prices; a demand restriction may lower
prices briefly but also reduces ownership opportunity for the households
the policy was meant to help.

The shift from 2022 onwards toward supply-side measures --- the Housing
Accelerator Fund, the GST removal on purpose-built rentals, federal pressure
on municipalities to upzone --- reflects a belated recognition of this basic
supply-and-demand reality.

%==========================================================
\begin{policydebate}{Is Foreign Ownership Driving Canadian Housing Unaffordability?}

\textbf{The Issue:} Few housing policy questions generate more political
heat and less empirical consensus than the role of foreign ownership. Advocates
argue that offshore capital has driven Vancouver and Toronto prices beyond
what domestic incomes could support. Sceptics argue foreign ownership is a
convenient scapegoat for domestic policy failures.

\textbf{Evidence Supporting a Significant Foreign Effect:}
\begin{itemize}
  \item CMHC's non-resident ownership data showed that non-residents owned
        approximately 2.4\% of residential properties in Ontario and 4.6\%
        in British Columbia, with concentrations in condominium markets.
  \item The 2016 BC Foreign Buyer Tax caused an immediate and dramatic
        (though temporary) cooling of Vancouver prices, particularly in
        the detached segment most favoured by offshore buyers.
\end{itemize}

\textbf{Evidence Against a Large Foreign Effect:}
\begin{itemize}
  \item Aggregate non-resident ownership of 2--5\% of total housing
        stock is insufficient to explain prices 3--5 times higher
        than domestic income fundamentals alone would support. If foreign
        ownership were the primary driver, the 2023 foreign buyer ban
        should have made housing affordable --- it did not.
  \item Canadian housing has appreciated across \textit{all} markets,
        including cities (Halifax, Saskatoon, Prince George) with essentially
        zero foreign ownership activity. A supply-constrained explanation
        applies everywhere; a foreign-buyer explanation applies only to
        a few markets.
  \item Many studies find that domestic investors and domestic demographic
        pressure are quantitatively more important drivers.
\end{itemize}

\textbf{Policy Implication:} The evidence is clear: Canada cannot tax its
way to housing affordability. Foreign buyer restrictions are not irrelevant,
but they address a marginal contributor to a problem driven primarily by
supply constraints and domestic demand. Supply must grow.
\end{policydebate}

%==========================================================
\begin{everydaylife}{The Mathematics of Your First Home}

Suppose you graduate from UNBC at age 22 earning \$58{,}000 per year in
Prince George, with student debt of \$28{,}000. The current benchmark price
for a two-bedroom home in Prince George is approximately \$380{,}000.

\textbf{Down payment:} A 5\% minimum down payment requires \$19{,}000,
plus CMHC insurance of approximately \$7{,}600 added to your loan.
A 20\% down payment (\$76{,}000) eliminates the insurance.

\textbf{Saving timeline:} After tax, rent (\$1{,}500/month), loan payments
(\$400/month), and living costs (\$1{,}800/month), you might save
\$600--\$1{,}000/month. At \$800/month, you accumulate \$76{,}000 in about
8 years --- but if prices appreciate at 3\% annually while you save, the
20\% down payment target grows from \$76{,}000 to \$96{,}200. You are on
a treadmill.

\textbf{At purchase (age 30):} A 25-year mortgage on a \$304{,}000 principal
at 4.5\% costs approximately \$1{,}655/month --- 28.6\% of your gross monthly
income of \$5{,}800. This is within the 30\% guideline, but tight.

\textbf{Key insight from the user cost formula:} Whether buying is rational
depends on the gap between $r$ (your mortgage rate) and $g$ (expected future
appreciation). If prices are flat or falling, the user cost rises sharply.
If rates fall and prices appreciate, the user cost makes owning attractive.
You cannot control either variable --- they are set by the Bank of Canada
and the housing market. What you can control is your savings rate and your
First Home Savings Account contributions. The mathematics of homeownership
in modern Canada is more uncertain, more leveraged, and more dependent on
external forces than at any previous point in the country's history.
\end{everydaylife}

%==========================================================
\section*{Chapter Summary}
%==========================================================
\addcontentsline{toc}{section}{Chapter Summary}

\begin{itemize}
  \item Housing is simultaneously a consumption good, a durable asset,
        a financial investment, and a location-specific good. These
        characteristics make housing markets unusually sensitive to
        interest rates, supply constraints, and demographic shifts.

  \item The short-run housing supply curve is nearly vertical in Canadian
        cities (supply elasticity $\varepsilon_s \approx 0.2$--$0.6$) due
        to zoning restrictions, slow permitting, development charges, and
        construction labour shortages. When supply is inelastic, demand
        shocks produce large price increases with minimal quantity response.

  \item The \textbf{user cost formula} ($\text{UC} = (r + \delta + \tau - g) \times P$)
        explains the rent-vs-buy decision. In 2021, deeply negative real
        rates and high expected capital gains drove user costs near zero,
        generating enormous ownership demand. By 2023, user costs had risen
        six-fold, making renting more economical in many markets.

  \item The \textbf{mortgage payment formula} (Equation~\ref{eq:mortgage_payment})
        shows that at 5.8\% for a \$700{,}000 home with 20\% down,
        the monthly payment is \$3{,}531, requiring gross income of at
        least \$141{,}240 to meet the 30\% affordability threshold.

  \item Vancouver (PTI 12.8) and Toronto (PTI 10.2) rank among the least
        affordable cities globally. The CMHC affordability index for
        Vancouver reached 92.1\% in 2024.

  \item House prices have risen more than four-fold nationally since 2000.
        The 2020--2022 COVID-era rate cuts triggered the most explosive
        price surge in Canadian history. The 2022--2023 rate hike cycle
        produced the first significant correction but left prices far
        above pre-pandemic levels.

  \item The supply-demand gap is structural: Canada builds 240{,}000--265{,}000
        units per year but needs 400{,}000--480{,}000 at current immigration
        levels. This shortfall drives both ownership price appreciation
        and the rental vacancy crisis.

  \item The rental vacancy rate has fallen below 2\% in most major
        Canadian cities, producing double-digit rent inflation and
        acute affordability problems for low-income renters.

  \item Most federal housing policies through 2022 were demand-side.
        Economic analysis demonstrates that demand subsidies in
        supply-constrained markets raise prices rather than improving
        affordability. Post-2022 policy has shifted toward supply-side
        instruments.

  \item International comparisons show that large non-market housing
        sectors, strong tenant protections, and permissive zoning produce
        significantly better affordability outcomes in comparable democracies.
\end{itemize}

%==========================================================
\section*{Key Terms}
%==========================================================
\addcontentsline{toc}{section}{Key Terms}

\begin{description}[style=nextline, leftmargin=3.5cm]
  \item[\term{Development charges}] Fees levied by municipalities on
        new construction to fund infrastructure costs.
  \item[\term{Financialization}] The increasing treatment of residential
        real estate as a financial asset rather than primarily as shelter.
  \item[\term{Housing starts}] New residential units on which construction
        has begun in a given period.
  \item[\term{Housing supply elasticity}] $\varepsilon_s = \%\Delta Q_s / \%\Delta P$;
        how responsive new construction is to house price changes.
  \item[\term{Normal good}] A good for which demand rises with income;
        housing has income elasticity of approximately 0.5--1.0.
  \item[\term{Mortgage stress test}] Federal requirement to qualify at
        contract rate + 2\% or the Bank of Canada benchmark rate,
        whichever is higher.
  \item[\term{Price-to-income ratio (PTI)}] Median house price divided by
        median annual household income; the primary affordability benchmark.
  \item[\term{Purpose-built rental}] Buildings constructed specifically
        for the long-term rental market.
  \item[\term{Rent control}] Government regulations limiting annual
        rent increases.
  \item[\term{Rental vacancy rate}] Percentage of purpose-built rental
        units that are unoccupied and available.
  \item[\term{Shelter cost burden}] Shelter costs as a share of household
        income; above 30\% is conventionally ``unaffordable.''
  \item[\term{Social housing}] Housing owned by government or non-profits
        and rented at below-market rates to income-qualifying households.
  \item[\term{Upzoning}] A land-use regulation change permitting higher-
        density development than was previously allowed.
  \item[\term{User cost of housing}] Annual economic cost of owning a home:
        $\text{UC} = (r + \delta + \tau - g) \times P$.
  \item[\term{Vacancy tax}] A tax on vacant residential properties,
        incentivising owners to rent units out.
  \item[\term{Zoning}] Municipal regulation specifying permitted land uses
        and development standards.
\end{description}

%==========================================================
\section*{Review Questions}
%==========================================================
\addcontentsline{toc}{section}{Review Questions}

\begin{enumerate}

\item \textbf{(Concept)} Explain why the short-run housing supply curve
is nearly vertical in major Canadian cities. What are the three most
important constraints on rapid housing supply expansion? Which is most
amenable to policy reform?

\item \textbf{(Application)} Using the user cost formula
(Equation~\ref{eq:user_cost}), calculate the annual user cost of owning
a \$900{,}000 home in each scenario. In which is buying most attractive
relative to paying \$2{,}800/month rent?
\begin{enumerate}[(a)]
  \item $r = 2\%$, $\delta = 2\%$, $\tau = 0.8\%$, $g = 5\%$
  \item $r = 5.5\%$, $\delta = 2\%$, $\tau = 0.8\%$, $g = 1\%$
  \item $r = 4.0\%$, $\delta = 2\%$, $\tau = 0.8\%$, $g = 3\%$
\end{enumerate}

\item \textbf{(Measurement)} Canada's average house price rose from
approximately \$530{,}000 in 2020 to \$796{,}000 in 2022.
\begin{enumerate}[(a)]
  \item What was the nominal percentage price increase?
  \item If median household income rose from \$84{,}000 to \$90{,}000
        over the same period, what happened to the PTI?
  \item What does the PTI change tell us that the nominal price increase alone does not?
\end{enumerate}

\item \textbf{(Policy)} A policymaker proposes a \$30{,}000 grant to
all first-time homebuyers. Using Figure~\ref{fig:housing_sd}, explain
what happens to house prices in: (a) a market with highly inelastic
supply; (b) a market with elastic supply. Which scenario reflects most
Canadian urban markets, and what does this imply about the policy's
effectiveness?

\item \textbf{(Rent control)} Using supply-and-demand analysis, explain
the short-run and long-run effects of a binding rent ceiling set below
the market rent. What is the key difference between the short-run and
long-run outcomes, and why does this matter for policy evaluation?

\item \textbf{(International comparison)} Canada's social housing stock
is approximately 3.5\% of total dwellings; the Netherlands is 30.5\%.
Explain the mechanism by which a larger social housing sector would
affect equilibrium rents in the private market. What are two obstacles
to Canada rapidly expanding its social housing stock?

\end{enumerate}

%==========================================================
\section*{Quantitative Problems}
%==========================================================
\addcontentsline{toc}{section}{Quantitative Problems}

\begin{enumerate}

\item \textbf{Mortgage Payment Calculations.}
Use Equation~\ref{eq:mortgage_payment} with Canadian semi-annual compounding
($r_m = (1 + r_{\text{annual}}/2)^{1/6} - 1$) and 25-year amortization.
\begin{enumerate}[(a)]
  \item A \$600{,}000 home, 10\% down, at 4.0\% annual rate. Monthly payment?
  \item Same home with 20\% down at 4.0\%. Monthly payment?
  \item Same home with 20\% down at 6.5\%. Monthly payment?
  \item For each scenario, what gross annual income is required for the
        payment to equal exactly 30\% of monthly gross income?
\end{enumerate}

\item \textbf{User Cost Analysis.}
A household in Kelowna, BC is choosing between renting and buying.
The home costs \$650{,}000; a comparable rental unit rents for
\$2{,}400/month (\$28{,}800/year).
\begin{enumerate}[(a)]
  \item Calculate the user cost under 2021 conditions ($r=2\%$, $\delta=2\%$,
        $\tau=0.8\%$, $g=6\%$) and under 2023 conditions ($r=5.8\%$,
        $\delta=2\%$, $\tau=0.8\%$, $g=0\%$). Is renting or buying more
        economical in each year?
  \item What annual capital gain rate $g$ (holding 2023 parameters constant)
        would make the household indifferent between renting and buying?
        Show your algebra.
  \item If the home price falls to \$560{,}000 but the rate remains at
        5.8\%, calculate the new user cost. Has affordability improved?
\end{enumerate}

\item \textbf{Housing Supply Elasticity.}
In Market A (Toronto), a 20\% price increase led to a 4\% increase in
housing starts. In Market B (Houston), the same 20\% increase led to
a 30\% increase in starts.
\begin{enumerate}[(a)]
  \item Calculate the price elasticity of housing supply in each market.
  \item Draw the supply curves, starting from common initial conditions
        of \$400{,}000 and 50{,}000 units.
  \item Suppose demand increases by 30\% in both markets (demand
        elasticity = $-0.8$). What is the new equilibrium price in each?
  \item What policy changes would move Market A toward Market B's
        supply elasticity?
\end{enumerate}

\item \textbf{Affordability Index Calculation.}
A city has the following characteristics: median benchmark price
\$820{,}000; 20\% down payment (\$164{,}000); 5.0\% mortgage rate;
25-year amortization; 0.65\% annual property tax rate; \$250/month
utilities; median annual household income \$102{,}000.
\begin{enumerate}[(a)]
  \item Calculate the monthly mortgage payment.
  \item Calculate the monthly property tax payment.
  \item Calculate the CMHC Affordability Index (total monthly housing
        costs as \% of median monthly gross income).
  \item Compare to Table~\ref{tab:affordability_index}. Which Canadian
        city does this most closely resemble?
\end{enumerate}

\item \textbf{Supply-Demand Gap.}
Canada's population grew by 1{,}053{,}000 persons in 2022.
\begin{enumerate}[(a)]
  \item At an average household size of 2.5 persons, how many new
        households were formed in 2022?
  \item Adding a demolition rate of 0.5\% of the existing stock
        ($\approx 75{,}000$ units), what total units were needed?
  \item Actual starts were 263{,}000. What was the shortfall?
  \item If this shortfall persists for 5 years, how large a housing
        deficit accumulates? What effect would you expect on rents?
\end{enumerate}

\item \textbf{Data Exercise.} Visit the CMHC website
(\url{www.cmhc-schl.gc.ca}) and find the most recent Housing Market
Outlook for your city.
\begin{enumerate}[(a)]
  \item Report the current rental vacancy rate and year-over-year
        change in average asking rents.
  \item Report housing starts over the past 12 months.
  \item Compare starts to CMHC's estimated required units. Is supply
        growing faster or slower than demand?
  \item Identify one supply-side policy your municipal or provincial
        government has announced. Using Figure~\ref{fig:housing_sd},
        evaluate whether it targets the right constraint.
\end{enumerate}

\end{enumerate}

%==========================================================
\section*{Essay Questions}
%==========================================================
\addcontentsline{toc}{section}{Essay Questions}

\begin{enumerate}

\item \textbf{(600--800 words)} ``Canada's housing crisis is fundamentally
a supply problem, and demand-side policies cannot solve it.'' Evaluate this
claim using the supply-and-demand model. Your essay should: (a) explain
why the short-run supply curve is nearly vertical; (b) use a diagram to
show why demand subsidies raise prices in supply-constrained markets;
(c) describe at least two supply-side policies and evaluate their
effectiveness; and (d) discuss whether demand-side interventions could
complement effective supply-side reforms.

\item \textbf{(500--700 words)} Compare Canada's housing outcomes to those
of Germany (Case Study 2). What are the three most important institutional
differences between the two countries? For each, assess the feasibility and
political obstacles to implementing a similar approach in Canada. Would you
recommend that Canada move toward the German model? Defend your answer.

\item \textbf{(400--600 words)} The user cost of housing in 2021 was near
zero in many Canadian markets, while the same calculation in 2023 yielded
a cost above 8\% of home value. Write an essay explaining what drove this
change, what its economic consequences were for the housing market, and
what it tells us about the relationship between monetary policy and housing
affordability. Connect your analysis explicitly to Chapter~\ref{ch:inflation}'s
discussion of the Bank of Canada's rate decisions.

\end{enumerate}

%==========================================================
\section*{Discussion Questions}
%==========================================================
\addcontentsline{toc}{section}{Discussion Questions}

\begin{enumerate}

\item The opening vignette describes three households: Vancouver engineers
who bought for \$2.1M, a nurse and teacher who cannot afford to buy, and a
Prince George mill worker paying high rent. Using concepts from this chapter,
explain why these three households face such different situations despite
the same national macroeconomic policies. Which economic force is most
responsible for each household's situation?

\item BC's 2023 upzoning legislation allows multiplexes on formerly
single-family lots. Who wins and who loses from this policy? Consider
existing homeowners, would-be first-time buyers, renters, developers,
and local municipalities. Does the aggregate economic benefit outweigh
the distributional costs?

\item Should Canada maintain the Agricultural Land Reserve in BC, given
that it constrains the supply of developable land around Metro Vancouver
and contributes to housing unaffordability? Or does food security and
environmental preservation justify the housing cost? How would an economist
structure the argument for each side?

\item The Case Study on rent control concludes that it protects existing
tenants at the expense of future tenants and new supply. Is this an
acceptable trade-off? Consider the perspective of: (a) a long-term renter
in a controlled unit in Toronto; (b) a newcomer searching for an apartment
in the same city; and (c) a developer deciding whether to build new rental
units. Is there a policy that achieves rent control's goals without its
supply-side costs?

\end{enumerate}

%==========================================================
%  CHAPTER 4: Labour Markets — Employment, Wages, and
%              Demographic Change
%  ECON 204 — Contemporary Economic Issues: A Canadian Perspective
%==========================================================

\chapter{Labour Markets: Employment, Wages, and Demographic Change}
\label{ch:labour}

%----------------------------------------------------------
%  OPENING VIGNETTE
%----------------------------------------------------------

\noindent In March 2023, two Canadians were navigating very different
labour markets --- in the same city.

Jerome, 42, had worked as a machinist at a Windsor, Ontario auto parts
plant for eighteen years. His unionized position paid \$34.20 per hour,
came with a defined-benefit pension, dental and vision coverage, and
four weeks of paid vacation. He had renegotiated his collective agreement
the previous fall and received a 6.5\% wage increase --- his largest in
a decade. His plant had just won a \$200 million retooling contract for
electric vehicle battery components. His employer was begging for more
skilled machinists and couldn't find them.

Kayla, 27, had a business degree from the University of Windsor and
worked as a ``gig coordinator'' for a logistics platform, classifying
her own work schedule and delivery routes through an app. Her pay was
\$19.50 per effective hour after accounting for vehicle costs, which
she covered herself. She had no benefits, no pension, no vacation pay,
and no job security beyond the next available delivery slot. On weeks when
demand was low, her earnings fell below minimum wage. On weeks when it was
high --- Christmas, Valentine's Day --- she worked 60 hours.

Jerome and Kayla live in the same city, are subject to the same provincial
employment standards, and are affected by the same monetary and fiscal
policies. Yet their labour market experiences are separated by a chasm:
one navigated by institutional arrangements built over decades of
collective bargaining; the other by algorithms, platform economics, and
legal classifications that exclude workers from the protections most
Canadians take for granted.

Understanding labour markets --- how wages are determined, why employment
fluctuates, how technological change and demographic aging are reshaping
the demand for and supply of work, and what policies can improve outcomes
for workers like Kayla without undermining conditions for workers like
Jerome --- is the subject of this chapter.

%----------------------------------------------------------
%  LEARNING OBJECTIVES
%----------------------------------------------------------

\begin{learningobjectives}
  \item Apply the supply-and-demand model to labour markets, distinguish
        the competitive from the monopsonistic wage-setting case, and
        use both to predict the employment effects of a minimum wage.
  \item Calculate the three key labour force statistics --- unemployment
        rate, employment rate, and labour force participation rate ---
        from raw population and employment data, and interpret each
        measure correctly.
  \item Distinguish among frictional, structural, and cyclical
        unemployment, define the natural rate of unemployment, and
        explain why the unemployment rate is an imperfect welfare measure.
  \item Interpret the Beveridge Curve and explain what an outward shift
        of the curve means for labour market matching efficiency.
  \item Analyse the wage-productivity gap in Canada, identify its causes,
        and evaluate evidence on the link between unionization decline
        and stagnant real wages.
  \item Compare standard and non-standard employment on wages, benefits,
        and security, and evaluate the policy debate over the employment
        status of platform workers.
  \item Assess the demographic pressures on Canada's labour market ---
        aging population, declining youth cohorts, and immigration ---
        and connect them to labour shortages, the dependency ratio, and
        long-run productivity.
\end{learningobjectives}

%==========================================================
\section{Labour Market Fundamentals}
\label{sec:labour_fundamentals}
%==========================================================

The \term{labour market} is the market in which workers sell their time
and skills to employers. Like any market, it has a supply side (workers
deciding how many hours to offer at different wage rates), a demand side
(employers deciding how many workers to hire at different wage rates), and
an equilibrium (the wage and employment level that clears the market).

\begin{defbox}{Labour Demand and Labour Supply}
\textbf{Labour demand} is derived from the demand for output.
A firm's profit-maximising condition for labour hiring is:
\[
  W = P \times \text{MP}_L
\]
where $W$ is the market wage, $P$ is the output price, and $\text{MP}_L$
is the marginal product of labour. This condition equates the wage to
the \term{value of the marginal product of labour (VMP$_L$)}, which is
the revenue generated by the last unit of labour hired. Labour demand
curves slope downward because, due to diminishing returns, each
additional worker adds less output (and hence less revenue) than the
previous one.

\textbf{Labour supply} reflects the trade-off between work (earning
income) and leisure (consuming non-work time). The supply curve
typically slopes upward: higher wages induce more labour supply,
as the opportunity cost of leisure rises. However, at very high wages,
the \term{income effect} (workers can afford more leisure) can dominate
the \term{substitution effect} (higher wage makes leisure more costly),
producing a backward-bending supply curve.
\end{defbox}

%----------------------------------------------------------
\subsection{The Competitive Labour Market}
%----------------------------------------------------------

In a perfectly competitive labour market, many workers and many firms
interact, no individual can influence the market wage, and the equilibrium
wage $W^*$ clears the market at employment $L^*$. This is the benchmark
model. Figure~\ref{fig:labour_markets} illustrates both the competitive
case and its main departure: monopsony.

\begin{figure}[H]
\centering
\begin{tikzpicture}[
  axisline/.style={-Stealth, thick, color=gray!70},
  lbl/.style={font=\small}
]

%---- PANEL A: COMPETITIVE MARKET ----
\begin{scope}[xshift=0cm]
\node[font=\small\bfseries, align=center] at (3.2,5.8)
  {Panel A: Competitive Market};
\draw[axisline] (0,0) -- (0,5.5) node[above, lbl] {$W$};
\draw[axisline] (0,0) -- (6.0,0) node[right, lbl] {$L$};
% Supply
\draw[casered, thick] (0.4,0.4) -- (5.4,4.9)
  node[right, font=\footnotesize, color=casered] {$S_L$};
% Demand
\draw[policyblue, thick] (0.4,4.9) -- (5.6,0.3)
  node[right, font=\footnotesize, color=policyblue] {$D_L = \text{VMP}_L$};
% Equilibrium
\fill[unbcgreen] (2.88,2.68) circle (3pt);
\draw[dashed, gray!50, thin] (0,2.68) -- (2.88,2.68)
  node[left, pos=0, font=\footnotesize] {$W^*$};
\draw[dashed, gray!50, thin] (2.88,0) -- (2.88,2.68)
  node[below, font=\footnotesize] {$L^*$};
% Minimum wage
\draw[dataorange, thick, dashed] (0,3.8) -- (5.8,3.8)
  node[right, font=\footnotesize, color=dataorange] {$\bar{W}_{\min}$};
\fill[dataorange] (1.62,3.8) circle (2.5pt);
\fill[dataorange] (4.35,3.8) circle (2.5pt);
\draw[<->, dataorange!80, thin] (1.62,4.1) -- (4.35,4.1)
  node[midway, above, font=\footnotesize, color=dataorange] {unemployment};
\node[left, font=\footnotesize, dataorange] at (0,3.8) {$\bar{W}$};
\end{scope}

%---- PANEL B: MONOPSONY ----
\begin{scope}[xshift=7.5cm]
\node[font=\small\bfseries, align=center] at (3.2,5.8)
  {Panel B: Monopsony};
\draw[axisline] (0,0) -- (0,5.5) node[above, lbl] {$W$};
\draw[axisline] (0,0) -- (6.0,0) node[right, lbl] {$L$};
% Supply (ACL)
\draw[casered, thick] (0.4,0.4) -- (5.4,4.9)
  node[right, font=\footnotesize, color=casered] {$S_L = \text{ACL}$};
% MLC (steeper than supply)
\draw[casered!60, thick, dashed] (0.4,0.8) -- (3.8,5.2)
  node[right, font=\footnotesize, color=casered!70] {MLC};
% VMP_L (demand)
\draw[policyblue, thick] (0.4,4.9) -- (5.6,0.3)
  node[right, font=\footnotesize, color=policyblue] {$D_L = \text{VMP}_L$};
% Monopsony equilibrium
\fill[casered!80] (2.15,2.15) circle (3pt);
% Monopsony wage (read off supply curve)
\draw[dashed, gray!50, thin] (0,1.55) -- (2.15,1.55)
  node[left, pos=0, font=\footnotesize] {$W_m$};
\draw[dashed, gray!50, thin] (2.15,0) -- (2.15,2.15)
  node[below, font=\footnotesize] {$L_m$};
% Competitive outcome for comparison
\fill[unbcgreen] (2.88,2.68) circle (2.5pt);
\draw[dashed, gray!20, thin] (0,2.68) -- (2.88,2.68)
  node[left, pos=0, font=\footnotesize, color=gray!60] {$W^*$};
\draw[dashed, gray!20, thin] (2.88,0) -- (2.88,2.68)
  node[below, font=\footnotesize, color=gray!60] {$L^*$};
% Minimum wage eliminating monopsony distortion
\draw[dataorange, thick, dashed] (0,2.68) -- (5.8,2.68)
  node[right, font=\footnotesize, color=dataorange] {$\bar{W}_{\min} = W^*$};
\node[font=\footnotesize, align=center, text=dataorange] at (4.2,1.3)
  {Min wage restores\\competitive outcome};
\end{scope}
\end{tikzpicture}
\caption{Labour Market Models: Competitive vs.\ Monopsony.
\textbf{Panel A} shows the competitive labour market. The minimum wage
$\bar{W}$ set above the competitive equilibrium $W^*$ creates unemployment
(excess labour supply, shown by the orange gap) --- the standard result.
\textbf{Panel B} shows a monopsony (single buyer of labour). The monopsonist
sets employment where MLC $=$ VMP$_L$, at $L_m < L^*$, and pays wage
$W_m < W^*$. Critically, a minimum wage set at $W^*$ eliminates the
monopsony distortion and \textit{increases} both employment and wages ---
the theoretical basis for why Card and Krueger (1994) found no employment
loss from minimum wage increases in some US fast-food markets. Canadian
labour markets, particularly in rural communities and in care sectors,
often have monopsonistic features.}
\label{fig:labour_markets}
\end{figure}

%----------------------------------------------------------
\subsection{Monopsony in Canadian Labour Markets}
%----------------------------------------------------------

The competitive model assumes workers can freely move between employers.
In practice, many Canadian workers face a \term{monopsonist} or
\term{oligopsonist} labour market --- a market where one or a few employers
dominate hiring. Classic examples include:

\begin{itemize}
  \item \textbf{Single-employer towns:} A sawmill in a remote Northern BC
        community, a mine in the Northwest Territories, or a fish processing
        plant in rural Newfoundland may be the only significant employer.
        Workers who do not want to relocate have little bargaining power.
  \item \textbf{Sector-specific skills:} A highly trained nurse or
        schoolteacher in a small city faces limited employer options
        even if the labour market is otherwise competitive.
  \item \textbf{Platform labour markets:} Ride-share and delivery platforms
        effectively set wages unilaterally for workers classified as
        contractors --- exhibiting monopsonistic wage-setting in practice
        even when no legal monopoly exists.
\end{itemize}

The monopsony model has an important policy implication: when labour markets
are monopsonistic, a minimum wage set up to the competitive wage level can
\textit{increase} employment and wages simultaneously. This is why the
empirical evidence on minimum wages is more nuanced than the simple
competitive model predicts. We return to this in Section~\ref{sec:min_wage}.

%==========================================================
\section{Labour Force Accounting}
\label{sec:labour_accounting}
%==========================================================

The Statistics Canada \term{Labour Force Survey (LFS)}, published monthly,
is Canada's primary source of employment and unemployment data. The LFS
divides the working-age population (WAP, persons aged 15 and over) into
three groups:

\begin{itemize}
  \item \textbf{Employed} ($E$): persons who worked at least one hour
        for pay or profit in the reference week, or who were absent from
        a job they held (due to illness, vacation, or labour dispute).
  \item \textbf{Unemployed} ($U$): persons who were without work, were
        available for work, and were actively seeking employment in the
        four weeks prior to the reference week.
  \item \textbf{Not in the labour force} ($N$): persons who are neither
        employed nor unemployed --- including retirees, students, caregivers,
        and discouraged workers who have stopped searching for jobs.
\end{itemize}

The labour force is $LF = E + U$. From these groups, three key statistics
are derived:

\begin{equation}
  \text{Unemployment Rate} = \frac{U}{LF} \times 100 = \frac{U}{E + U} \times 100
  \label{eq:unemployment_rate}
\end{equation}

\begin{equation}
  \text{Labour Force Participation Rate (LFPR)} = \frac{LF}{\text{WAP}} \times 100
  \label{eq:lfpr}
\end{equation}

\begin{equation}
  \text{Employment Rate} = \frac{E}{\text{WAP}} \times 100
  \label{eq:employment_rate}
\end{equation}

\noindent\textbf{Worked Example 4.1 --- Labour Force Accounting, Canada, Mid-2024}

\begin{center}
\renewcommand{\arraystretch}{1.35}
\begin{tabular}{lr}
\toprule
\textbf{Category} & \textbf{Value} \\
\midrule
Working-age population (WAP, 15+)             & 34{,}150{,}000 \\
Employed ($E$)                                & 20{,}615{,}000 \\
Unemployed ($U$)                              &  1{,}371{,}000 \\
Not in labour force ($N$)                     & 12{,}164{,}000 \\
\midrule
Labour force ($LF = E + U$)                   & 21{,}986{,}000 \\
\midrule
Unemployment rate ($U/LF$)                    & 6.2\% \\
Labour force participation rate ($LF$/WAP)    & 64.4\% \\
Employment rate ($E$/WAP)                     & 60.4\% \\
\bottomrule
\end{tabular}
\end{center}

\noindent\textit{Interpretation:}
\begin{itemize}
  \item A 6.2\% unemployment rate means approximately 1.37 million
        Canadians who want to work and are actively searching cannot
        find a job.
  \item A 64.4\% LFPR means that 35.6\% of working-age Canadians are
        neither employed nor looking for work.
  \item An employment rate of 60.4\% means only 60 of every 100 Canadians
        aged 15+ hold a paying job.
\end{itemize}

\begin{keytakeaway}
The unemployment rate is the most cited labour market statistic, but it
is an imperfect welfare measure. It excludes: \textbf{discouraged workers}
(who want work but stopped searching); \textbf{involuntary part-timers}
(who want full-time work but can only find part-time); and \textbf{workers
in jobs below their skill level}. Statistics Canada's R8 measure of
underemployment, which includes these groups, consistently runs 3--5
percentage points above the headline unemployment rate.
\end{keytakeaway}

%----------------------------------------------------------
\subsection{Canada's Labour Market Record, 2010--2026}
%----------------------------------------------------------

\begin{figure}[H]
\centering
\begin{tikzpicture}
\begin{axis}[
  width=14cm, height=7.5cm,
  xlabel={Year},
  ylabel={Rate (\%)},
  xlabel style={font=\small},
  ylabel style={font=\small},
  xmin=2009.5, xmax=2026.5,
  ymin=0, ymax=75,
  xtick={2010,2012,2014,2016,2018,2020,2022,2024,2026},
  xticklabel style={font=\small},
  ytick={0,10,20,30,40,50,60,70},
  yticklabel style={font=\small},
  ymajorgrids=true,
  grid style={dashed, gray!20},
  legend pos=south east,
  legend style={font=\small},
  clip=false
]
% LFPR
\addplot[policyblue, thick, mark=none] coordinates {
  (2010,67.1)(2011,66.8)(2012,66.6)(2013,66.5)(2014,65.8)
  (2015,65.9)(2016,65.7)(2017,65.8)(2018,65.5)(2019,65.6)
  (2020,62.4)(2021,64.7)(2022,64.8)(2023,65.1)(2024,64.4)(2025,64.1)(2026,64.3)
};
\addlegendentry{LFPR (\%, left)}
% Employment rate
\addplot[unbcgreen, thick, dashed, mark=none] coordinates {
  (2010,61.6)(2011,61.7)(2012,61.8)(2013,62.0)(2014,61.4)
  (2015,61.4)(2016,61.0)(2017,61.6)(2018,62.0)(2019,61.8)
  (2020,56.5)(2021,59.9)(2022,61.5)(2023,61.3)(2024,60.4)(2025,59.9)(2026,60.1)
};
\addlegendentry{Employment rate (\%)}
% Unemployment rate (scaled up ×5 for visibility on same axis)
\addplot[casered, thick, dotted, mark=none] coordinates {
  (2010,8.1)(2011,7.4)(2012,7.2)(2013,7.1)(2014,6.9)
  (2015,6.9)(2016,7.0)(2017,6.3)(2018,5.8)(2019,5.7)
  (2020,9.5)(2021,7.5)(2022,5.3)(2023,5.4)(2024,6.3)(2025,6.7)(2026,6.5)
};
\addlegendentry{Unemployment rate (\%)}
% COVID annotation
\node[font=\footnotesize, align=center, color=gray] at (axis cs:2020,50)
  {COVID-19\\lockdowns};
\draw[->, gray!60, thin] (axis cs:2020,52) -- (axis cs:2020,57);
\node[font=\footnotesize, color=casered] at (axis cs:2020,11)
  {u=9.5\%};
\node[font=\footnotesize, color=casered] at (axis cs:2022.5,4.5)
  {u=5.3\%\\(50-yr low)};
\end{axis}
\end{tikzpicture}
\caption{Canadian Labour Force Participation Rate, Employment Rate, and
Unemployment Rate, 2010--2026. The three series together tell a richer
story than any one alone. The COVID-19 collapse of April 2020 is visible
in all three: the unemployment rate spiked to 9.5\%, the employment rate
fell to 56.5\%, and the LFPR dropped to 62.4\% as many workers stopped
searching. The 2022 labour market tightness is equally clear: unemployment
at 4.9--5.3\% (its lowest since the 1970s) while LFPR had recovered to
near pre-pandemic levels. The subsequent 2024--2025 softening reflects
the economic slowdown induced by the Bank of Canada's rate-hike cycle
(Chapter~\ref{ch:inflation}).}
\label{fig:lf_statistics}
\source{Statistics Canada, Table 14-10-0287-01, Labour Force Survey,
seasonally adjusted. 2025--2026 estimated.}
\end{figure}

%==========================================================
\section{Types of Unemployment}
\label{sec:unemployment_types}
%==========================================================

Not all unemployment is the same. Economists distinguish three conceptually
distinct types, each with different causes and policy implications.

\begin{defbox}{Three Types of Unemployment}
\textbf{Frictional unemployment} arises from the normal time it takes
to match workers with jobs. Even in a healthy economy, workers quit,
are dismissed, or leave school and must search for the right job.
Firms take time to evaluate candidates. This process generates temporary
unemployment that is voluntary, transitional, and inevitable in any
dynamic market economy. Estimated at 1--2\% of the labour force.

\textbf{Structural unemployment} arises when workers' skills are
mismatched with available jobs --- either because of technological
change (automation displacing specific occupations), geographic
mismatch (jobs in growing cities, workers in declining regions), or
industry restructuring (oil sands declining while green tech expands).
Structural unemployment tends to be longer-duration and is not cured
by aggregate demand stimulus. Policy responses include retraining,
relocation assistance, and apprenticeship programs.

\textbf{Cyclical unemployment} arises from insufficient aggregate
demand during economic downturns. When GDP falls below potential,
firms reduce hiring and may lay off workers; the resulting unemployment
is correlated with the business cycle. Cyclical unemployment is the
target of countercyclical monetary and fiscal policy.

The \textbf{natural rate of unemployment} $u_n$ (also called the
NAIRU) is the unemployment rate that prevails when the economy is
at potential output, consisting of frictional and structural
unemployment only. For Canada, the Bank of Canada estimates
$u_n \approx 5.5$--$6.0\%$ in the mid-2020s.
\end{defbox}

The distinction between structural and cyclical unemployment matters
enormously for policy. When the Bank of Canada raised rates in 2022--2023,
it was targeting the \textit{cyclical} component of inflation (excess
demand). The resulting rise in unemployment from 5.3\% to 6.5--6.8\%
by 2024--25 reflects both cyclical slack (deliberate demand reduction)
and structural factors (immigration increasing labour supply faster than
job creation). A policymaker who mistakes structural unemployment for
cyclical --- and responds with rate cuts --- risks reflating the economy
before the structural adjustment is complete.

%==========================================================
\section{The Beveridge Curve: Labour Market Matching}
\label{sec:beveridge_curve}
%==========================================================

The \term{Beveridge Curve} plots the job vacancy rate (share of positions
unfilled) against the unemployment rate over time. It describes the
efficiency of the process by which workers and jobs are matched.

A well-functioning labour market traces a stable, downward-sloping curve:
when unemployment is high (recession), vacancies are low (few firms
hiring); when unemployment is low (tight market), vacancies are high
(firms competing intensely for scarce workers). The position of the curve
reflects \term{matching efficiency} --- how well the process of connecting
workers to jobs works.

\begin{figure}[H]
\centering
\begin{tikzpicture}
\begin{axis}[
  width=10cm, height=9cm,
  xlabel={Unemployment Rate (\%)},
  ylabel={Job Vacancy Rate (\%)},
  xlabel style={font=\small},
  ylabel style={font=\small},
  xmin=4.5, xmax=10.5,
  ymin=0.5, ymax=5.5,
  xtick={5,6,7,8,9,10},
  ytick={1,2,3,4,5},
  xticklabel style={font=\small},
  yticklabel style={font=\small},
  ymajorgrids=true, xmajorgrids=true,
  grid style={dashed, gray!20},
  clip=false
]
% Pre-COVID Beveridge curve (reference)
\addplot[gray!40, dashed, thick, domain=4.5:10.5]
  {1.8 + 0.8*exp(-0.45*(x-5.0))}
  node[above right, pos=0.05, font=\footnotesize, color=gray!60]
  {Pre-COVID curve};
% Post-COVID Beveridge curve (outward shift)
\addplot[unbcgreen!50, dashed, thick, domain=4.5:10.5]
  {2.5 + 1.2*exp(-0.45*(x-5.0))}
  node[above right, pos=0.05, font=\footnotesize, color=unbcgreen!70]
  {Post-COVID curve};
% Data points
\addplot[only marks, mark=*, mark size=3pt, policyblue]
  coordinates {
    (6.9,1.5)(7.0,1.6)(6.3,1.9)(5.8,2.1)(5.7,2.2)
  };
\addplot[only marks, mark=square*, mark size=3pt, casered]
  coordinates {
    (9.5,1.4)(7.5,3.2)(5.3,4.5)(5.4,3.4)(6.3,2.8)(6.7,2.5)
  };
% Year labels
\node[font=\footnotesize, policyblue, above right] at (axis cs:6.9,1.5) {15};
\node[font=\footnotesize, policyblue, above right] at (axis cs:7.0,1.6) {16};
\node[font=\footnotesize, policyblue, above right] at (axis cs:6.3,1.9) {17};
\node[font=\footnotesize, policyblue, above right] at (axis cs:5.8,2.1) {18};
\node[font=\footnotesize, policyblue, above right] at (axis cs:5.7,2.2) {19};
\node[font=\footnotesize, casered, right] at (axis cs:9.5,1.4) {20};
\node[font=\footnotesize, casered, above right] at (axis cs:7.5,3.2) {21};
\node[font=\footnotesize, casered, above] at (axis cs:5.3,4.5) {22};
\node[font=\footnotesize, casered, right] at (axis cs:5.4,3.4) {23};
\node[font=\footnotesize, casered, below right] at (axis cs:6.3,2.8) {24};
\node[font=\footnotesize, casered, below right] at (axis cs:6.7,2.5) {25};
% Arrow showing outward shift
\draw[->, unbcgreen!70, thick] (axis cs:6.5,2.3) -- (axis cs:6.3,3.1);
\node[font=\footnotesize, unbcgreen!80, align=center] at (axis cs:7.0,3.7)
  {Outward\\shift:\\matching\\inefficiency};
\end{axis}
\end{tikzpicture}
\caption{Beveridge Curve: Canada, 2015--2025. Blue circles (labelled
15--19) show the pre-pandemic period tracing a stable downward-sloping
relationship near the grey reference curve. Red squares (labelled 20--25)
show the post-COVID period. The 2021--2022 points sit far to the
upper-left of the pre-COVID curve: very high vacancies (4.5\% in 2022)
combined with still-elevated unemployment (5.3\%), suggesting that many
unfilled positions could not be matched with available workers. This
\textit{outward shift} of the Beveridge curve indicates reduced matching
efficiency --- caused by skills mismatches, geographic mismatch between
vacancies and workers, and pandemic-related labour market restructuring.
By 2024--25, the curve had shifted partially back toward its pre-COVID
position as labour markets adjusted.}
\label{fig:beveridge}
\source{Statistics Canada, Job Vacancy and Wage Survey (Table 14-10-0325-01);
Labour Force Survey (Table 14-10-0287-01). Year labels show last two digits.}
\end{figure}

An outward Beveridge shift has specific policy implications: it means
that reducing the unemployment rate by cutting interest rates (stimulating
demand) will not efficiently eliminate the vacancies. The unfilled positions
require workers with skills that the unemployed do not currently possess,
or are located where unemployed workers do not currently live. The
policy solution is \textit{structural}: retraining, apprenticeship
programs, credential recognition for immigrants, and improved geographic
labour mobility --- not monetary stimulus.

%==========================================================
\section{Wage Determination and the Wage-Productivity Gap}
\label{sec:wages}
%==========================================================

%----------------------------------------------------------
\subsection{Competitive Wage Determination}
%----------------------------------------------------------

In a competitive labour market, the equilibrium wage equals the value of
the marginal product of labour:
\begin{equation}
  W^* = P \times \text{MP}_L
  \label{eq:vmp}
\end{equation}

This means wages grow with: (1) increases in the output price $P$ (firms
can pay more when their revenues rise); and (2) increases in labour
productivity $\text{MP}_L$ (more productive workers generate more revenue
per unit of labour). The fundamental prediction is that real wages and
labour productivity should grow together over the long run.

%----------------------------------------------------------
\subsection{The Wage-Productivity Gap: Canada's Most Important Labour Statistic}
%----------------------------------------------------------

The prediction that real wages and labour productivity should grow together
has not been borne out in Canada since the 1970s. Figure~\ref{fig:wage_productivity}
shows that while labour productivity has grown substantially since 1976,
real wages have grown significantly less.

\begin{figure}[H]
\centering
\begin{tikzpicture}
\begin{axis}[
  width=14cm, height=7.5cm,
  xlabel={Year},
  ylabel={Index (1976 = 100)},
  xlabel style={font=\small},
  ylabel style={font=\small},
  xmin=1975, xmax=2026,
  ymin=90, ymax=195,
  xtick={1980,1985,1990,1995,2000,2005,2010,2015,2020,2025},
  xticklabel style={font=\small},
  ytick={100,120,140,160,180,200},
  yticklabel style={font=\small},
  ymajorgrids=true,
  grid style={dashed, gray!20},
  legend pos=north west,
  legend style={font=\small},
  clip=false
]
% Labour productivity index
\addplot[policyblue, thick, mark=none] coordinates {
  (1976,100)(1980,105)(1985,112)(1990,119)(1995,128)
  (2000,142)(2005,152)(2010,157)(2015,163)(2020,168)(2022,170)(2024,173)(2025,175)
};
\addlegendentry{Labour productivity (real output per hour worked)}
% Real wages index
\addplot[casered, thick, dashed, mark=none] coordinates {
  (1976,100)(1980,101)(1985,99)(1990,104)(1995,102)
  (2000,109)(2005,116)(2010,120)(2015,125)(2020,132)(2022,129)(2024,131)(2025,132)
};
\addlegendentry{Real average weekly earnings}
% Gap annotation
\draw[<->, gray!60, thick] (axis cs:2024,131) -- (axis cs:2024,173);
\node[right, font=\footnotesize, color=gray!60, align=left] at (axis cs:2024.5,152)
  {Gap:\\$\approx$42 pts};
\node[font=\footnotesize, policyblue] at (axis cs:2000,147)
  {Productivity};
\node[font=\footnotesize, casered] at (axis cs:2000,107)
  {Real wages};
\end{axis}
\end{tikzpicture}
\caption{Real Wages vs.\ Labour Productivity, Canada, 1976--2025 (Index 1976=100).
Labour productivity --- real output per hour worked --- has grown by approximately
75\% since 1976. Real average weekly earnings have grown by only approximately
32\% over the same period. The divergence began in the early 1980s and has
persisted, implying that workers have captured a shrinking share of the gains
from productivity growth. This gap is one of the most important distributional
facts in the Canadian economy. It is associated with declining union coverage,
rising profit margins, the growth of monopsony power, and the shift from goods
to services employment --- all discussed below.}
\label{fig:wage_productivity}
\source{Statistics Canada, Table 36-10-0222-01 (multifactor productivity);
Table 14-10-0064-01 (average weekly earnings); author index calculations.
Real wages deflated using CPI (2002=100).}
\end{figure}

The wage-productivity gap implies that the \term{labour share of income}
--- wages and salaries as a fraction of national income --- has declined
over the past four decades, while the \term{capital share} (profits,
dividends, interest) has risen. This matters for inequality (explored
in Chapter~\ref{ch:inequality}) because capital income is far more
concentrated among high-income households than labour income.

The causes of the wage-productivity gap are contested among economists,
but the leading explanations include:

\begin{enumerate}
  \item \textbf{Declining unionization:} Unions compress wage distributions
        and raise wages for medium-skill workers through collective bargaining.
        The decline in union coverage described in Section~\ref{sec:unions}
        has removed a key mechanism for translating productivity gains into
        wage gains.
  \item \textbf{Rising market power of employers (monopsony):} As large
        firms and platforms grow, their ability to set wages unilaterally
        rather than take market wages increases --- widening the gap between
        $W$ and $\text{VMP}_L$.
  \item \textbf{Trade and offshoring:} Competitive pressure from lower-wage
        countries constrains wage growth in tradable sectors, particularly
        manufacturing.
  \item \textbf{Technological change:} Automation has raised the marginal
        product of capital relative to labour in some sectors, shifting
        factor income shares toward capital.
  \item \textbf{Composition effects:} The shift from high-wage manufacturing
        to lower-wage services employment changes the average measured wage
        even if wages within each sector are growing normally.
\end{enumerate}

%==========================================================
\section{Unionization and Collective Bargaining}
\label{sec:unions}
%==========================================================

A \term{trade union} is a collective organisation of workers that bargains
with employers on behalf of its members over wages, benefits, hours, and
working conditions. \term{Collective bargaining} replaces individual
wage negotiation with a negotiation between the union (representing all
workers in a bargaining unit) and the employer. The union's market power
derives from the credible threat of collective action --- a strike.

\begin{figure}[H]
\centering
\begin{tikzpicture}
\begin{axis}[
  ybar,
  bar width=0.35cm,
  width=13cm, height=7cm,
  symbolic x coords={{All industries},{Public admin},{Education},{Healthcare},
                     {Utilities},{Construction},{Manufacturing},
                     {Finance},{Retail},{Food service}},
  xtick=data,
  x tick label style={rotate=40, anchor=east, font=\small},
  ylabel={Union Coverage Rate (\%)},
  ylabel style={font=\small},
  ymin=0, ymax=85,
  ytick={0,10,20,30,40,50,60,70,80},
  yticklabel style={font=\small},
  ymajorgrids=true,
  grid style={dashed, gray!25},
  enlarge x limits=0.05,
  legend pos=north east,
  legend style={font=\small},
  clip=false
]
% 2004 bars
\addplot[fill=policyblue!70, draw=policyblue!90] coordinates {
  ({All industries},32.2)({Public admin},74.3)({Education},72.0)({Healthcare},57.2)
  ({Utilities},59.8)({Construction},31.4)({Manufacturing},30.5)
  ({Finance},8.5)({Retail},13.8)({Food service},7.2)
};
\addlegendentry{2004}
% 2024 bars
\addplot[fill=casered!60, draw=casered!80] coordinates {
  ({All industries},27.8)({Public admin},73.1)({Education},70.5)({Healthcare},54.3)
  ({Utilities},57.2)({Construction},29.8)({Manufacturing},22.4)
  ({Finance},6.1)({Retail},10.2)({Food service},5.8)
};
\addlegendentry{2024}
\end{axis}
\end{tikzpicture}
\caption{Union Coverage Rate by Industry, Canada, 2004 vs.\ 2024.
Unions are heavily concentrated in the public sector --- government,
education, healthcare, and utilities --- where coverage exceeds 54--73\%.
Private-sector union coverage is far lower and declining: manufacturing
fell from 30.5\% to 22.4\%; retail from 13.8\% to 10.2\%. The most
dramatic decline has been in the private sector overall, where coverage
has fallen to approximately 14\%. This bifurcated structure --- heavily
unionized public sector, weakly unionized private sector --- has important
distributional consequences: public-sector workers have maintained real
wages and benefit coverage while many private-sector workers (especially
in retail and food service) have seen stagnant real wages over the same period.}
\label{fig:unionization}
\source{Statistics Canada, Labour Force Survey, Union Status supplement;
Workplace and Employee Survey. 2024 values estimated from recent LFS data.}
\end{figure}

The economic evidence on unions' effects in Canada is reasonably clear:

\begin{itemize}
  \item \textbf{Union wage premium:} Unionized workers earn approximately
        15--20\% more than comparable non-unionized workers in the same
        industry and occupation, after controlling for observable differences.
        The premium is larger for workers without post-secondary education.
  \item \textbf{Benefit premium:} Unionized workers are significantly more
        likely to have employer-sponsored pension plans, drug benefits, and
        dental coverage.
  \item \textbf{Wage compression:} Unions tend to compress the wage
        distribution, raising wages at the bottom more than at the top.
        As union coverage has declined, within-industry wage inequality
        has increased.
  \item \textbf{Productivity effects:} The evidence on unions and productivity
        is mixed. The ``collective voice'' hypothesis suggests unions improve
        productivity by reducing turnover and providing a mechanism for workers
        to communicate grievances. The ``monopoly'' hypothesis suggests unions
        create inefficiencies by raising wages above competitive levels.
        On balance, Canadian studies find small positive or neutral productivity
        effects for most unionized firms.
\end{itemize}

%==========================================================
\section{Minimum Wages in Canada}
\label{sec:min_wage}
%==========================================================

\begin{defbox}{The Minimum Wage}
A \textbf{minimum wage} is a legislated floor on the hourly wage that
employers may legally pay. In Canada, minimum wages are set provincially,
creating significant variation across jurisdictions. There is also a
federal minimum wage covering workers in federally regulated industries
(banking, telecommunications, air transport, interprovincial trucking).

\textit{Economic theory predicts opposite effects under different market
structures:}
\begin{itemize}
  \item \textbf{Competitive market:} A minimum wage above the competitive
        equilibrium causes unemployment (excess labour supply). This is
        Panel A of Figure~\ref{fig:labour_markets}.
  \item \textbf{Monopsony:} A minimum wage up to the competitive wage
        level increases both wages and employment. Panel B of
        Figure~\ref{fig:labour_markets}.
\end{itemize}

The \textit{empirical} evidence (summarised below) is closest to the
monopsony prediction: moderate minimum wage increases in Canada and the
US have generally had small or negligible employment effects, with
significant wage gains for affected workers.
\end{defbox}

\begin{figure}[H]
\centering
\begin{tikzpicture}
\begin{axis}[
  ybar,
  bar width=0.55cm,
  width=13cm, height=6.5cm,
  symbolic x coords={{NL},{NS},{NB},{PEI},{QC},{ON},{MB},{SK},{AB},{BC},{Federal}},
  xtick=data,
  x tick label style={font=\small},
  ylabel={Minimum Wage (\$ per hour), 2025},
  ylabel style={font=\small},
  ymin=13, ymax=19,
  ytick={13,14,15,16,17,18,19},
  yticklabel style={font=\small},
  ymajorgrids=true,
  grid style={dashed, gray!25},
  enlarge x limits=0.06,
  nodes near coords,
  nodes near coords style={font=\footnotesize},
  every axis plot/.append style={fill=unbcgreen!65, draw=unbcgreen!85}
]
\addplot coordinates {
  ({NL},  15.60)({NS},  15.70)({NB},  15.30)({PEI}, 16.00)({QC},  16.10)
  ({ON},  17.20)({MB},  15.80)({SK},  15.00)({AB},  15.00)({BC},  17.40)({Federal},17.30)
};
\end{axis}
\end{tikzpicture}
\caption{Provincial and Federal Minimum Wages, Canada, 2025.
British Columbia (\$17.40) and Ontario (\$17.20) have the highest minimum
wages; Alberta and Saskatchewan (\$15.00) the lowest. The federal minimum
wage (\$17.30) applies to workers in federally regulated industries and
since 2021 is indexed to the Consumer Price Index, preventing real erosion.
The range of \$15.00--\$17.40 represents a 16\% spread across provincial
jurisdictions --- significant enough to affect labour market decisions near
provincial borders and in industries where workers are geographically mobile.}
\label{fig:min_wages}
\source{Employment and Social Development Canada; provincial governments.
Rates as of mid-2025. Some provinces index automatically to CPI.}
\end{figure}

\begin{canexample}{Card and Krueger (1994) and the Canadian Evidence}
David Card and Alan Krueger's 1994 study compared fast-food employment
in New Jersey (which raised its minimum wage from \$4.25 to \$5.05) to
neighbouring Pennsylvania (no change). They found \textit{no significant
employment loss} in New Jersey --- contradicting the simple competitive
model's prediction and consistent with the monopsony model.

Card received the 2021 Nobel Prize in Economics, partly for this work.

Canadian evidence broadly supports the Card-Krueger finding for moderate
increases. A 2019 study by Dube, Lindner, and Reich using Canadian
province-level data found that a 10\% minimum wage increase reduced
employment in the lowest-wage sectors by approximately 1--3\% ---
a small effect relative to the wage gain. A 2021 PBO analysis found
that Ontario's 2018 increase to \$14/hour (a 21\% jump) had modest
negative employment effects of roughly 0.2--0.4\% --- again, small
relative to the wage increase achieved.

However, these averages conceal important heterogeneity. Small businesses
in rural, low-income markets face less monopsony power and therefore
exhibit larger negative employment effects from minimum wage increases
than large firms in urban markets. The ``right'' minimum wage is not
uniform across the vast diversity of Canadian labour markets.
\end{canexample}

%==========================================================
\section{Non-Standard Employment and the Gig Economy}
\label{sec:nonstandard}
%==========================================================

The standard employment relationship --- full-time, permanent, with a single
employer, receiving benefits and entitled to all employment standards --- has
been gradually replaced in parts of the Canadian economy by an increasingly
diverse set of \term{non-standard employment} arrangements.

\begin{figure}[H]
\centering
\begin{tikzpicture}
\begin{axis}[
  width=14cm, height=6.5cm,
  xlabel={Year},
  ylabel={Share of Total Employment (\%)},
  xlabel style={font=\small},
  ylabel style={font=\small},
  xmin=1999, xmax=2025.5,
  ymin=0, ymax=20,
  xtick={2000,2004,2008,2012,2016,2020,2024},
  xticklabel style={font=\small},
  ytick={0,5,10,15,20},
  yticklabel style={font=\small},
  ymajorgrids=true,
  grid style={dashed, gray!20},
  legend pos=north east,
  legend style={font=\small},
  clip=false
]
% Own-account self-employment
\addplot[policyblue, thick, mark=none] coordinates {
  (2000,10.5)(2004,10.8)(2008,10.2)(2012,10.5)(2016,10.9)
  (2020,11.4)(2021,10.6)(2022,10.4)(2023,10.6)(2024,10.8)
};
\addlegendentry{Own-account self-employment}
% Temporary employment
\addplot[casered, thick, dashed, mark=none] coordinates {
  (2000,11.8)(2004,12.4)(2008,12.6)(2012,12.7)(2016,13.2)
  (2020,10.4)(2021,12.5)(2022,13.1)(2023,13.5)(2024,13.7)
};
\addlegendentry{Temporary employment}
% Involuntary part-time
\addplot[dataorange, thick, dotted, mark=none] coordinates {
  (2000,4.2)(2004,4.5)(2008,4.0)(2012,5.2)(2016,5.8)
  (2020,6.5)(2021,5.2)(2022,4.2)(2023,4.5)(2024,4.8)
};
\addlegendentry{Involuntary part-time}
\end{axis}
\end{tikzpicture}
\caption{Non-Standard Employment as a Share of Total Employment, Canada,
2000--2024. Three forms of non-standard work are shown: own-account self-employment
(solo workers without employees, including many gig workers); temporary employment
(term contracts, casual, seasonal); and involuntary part-time work (workers who
want full-time hours but can only find part-time). Together, these categories
account for approximately 29\% of total Canadian employment. Temporary employment
shows a cyclical pattern, falling sharply in the COVID recession as these workers
lost jobs first and recovering quickly. Own-account self-employment has grown
modestly as platform work (delivery, ride-share, freelancing) expands.}
\label{fig:nonstandard}
\source{Statistics Canada, Labour Force Survey, Table 14-10-0190-01. Annual averages.}
\end{figure}

\begin{table}[H]
\centering
\caption{Standard vs.\ Non-Standard Employment: Wages, Benefits, and Security}
\label{tab:employment_comparison}
\renewcommand{\arraystretch}{1.30}
\begin{tabular}{lrrrr}
\toprule
\textbf{Measure} & \textbf{Permanent} & \textbf{Temporary} & \textbf{Own-Account} & \textbf{Platform} \\
                 & \textbf{Full-Time} & \textbf{Employment} & \textbf{Self-Empl.} & \textbf{(Gig)} \\
\midrule
Median hourly wage (\$)   & 31.50 & 22.40 & 24.20 & 18.50 \\
Employer pension (\%)     &  55.3 &  14.2 &   0.0 &   0.0 \\
Employer health benefits  &  64.1 &  19.8 &   0.0 &   0.0 \\
EI eligible (\%)          &  98.0 &  75.0 &   2.5 &   2.5 \\
Paid vacation (\%)        &  97.5 &  43.2 &   0.0 &   0.0 \\
Job tenure (median yrs)   &   7.2 &   1.3 &   5.8 &   1.1 \\
Union coverage (\%)       &  30.2 &   9.5 &   1.1 &   0.5 \\
\bottomrule
\end{tabular}
\source{Statistics Canada, Survey on Employment and Employee Compensation;
Gig Economy Study, Institute for Work and Health, 2024.
Platform/gig values are estimates; sector varies widely.}
\end{table}

Table~\ref{tab:employment_comparison} illustrates the profound gap between
standard and non-standard employment on every dimension: wages, benefits,
Employment Insurance eligibility, and vacation entitlements. This gap has
grown as the gig economy has expanded, and it has concentrated among
younger workers, recent immigrants, and workers without post-secondary
credentials --- precisely the groups with the least bargaining power.

\begin{thinkbox}
Are gig workers ``entrepreneurs who choose flexibility'' or ``misclassified
employees who lack protection''? The answer matters enormously for policy.
If gig workers are choosing flexibility, reclassifying them as employees
could destroy the business model that makes their income possible. If they
are misclassified employees, the current classification denies them legal
rights they are entitled to under employment standards law.

The economic evidence suggests both are partially true: some platform
workers genuinely prefer flexible arrangements, while others are platform
workers because traditional employment is not available to them. A
nuanced policy response --- as adopted in the UK (where Uber drivers
were reclassified as ``workers,'' a middle category with some but not all
employee protections) --- may serve this diversity better than a binary
classification.
\end{thinkbox}

%==========================================================
\section{Demographic Change and the Labour Market}
\label{sec:demographics}
%==========================================================

Canada's labour market is being reshaped by a demographic transition
that was decades in the making. The \term{baby boom generation} (born
1946--1964) is now retiring at the rate of approximately 450{,}000 persons
per year. The cohorts entering the labour force behind them are smaller
(fertility has been below replacement since the 1970s) and less numerous.
The result is a structural tightening of the labour supply that
complements --- and sometimes competes with --- immigration as the
primary mechanism for maintaining the labour force.

\begin{defbox}{The Dependency Ratio}
The \textbf{old-age dependency ratio} measures the number of persons aged
65 and over per 100 persons of working age (15--64):
\begin{equation}
  \text{OADR} = \frac{\text{Population}_{65+}}{\text{Population}_{15-64}} \times 100
  \label{eq:oadr}
\end{equation}

As the baby boom retires and life expectancy increases, Canada's OADR is
rising rapidly: from approximately 22.5 in 2010, to 28.2 today, and
projected to reach 40--45 by 2050 under current fertility assumptions.
Each retiree withdraws from the labour supply (reducing $LF$) while
simultaneously increasing demand for government services (healthcare,
OAS/GIS) that must be financed by the remaining working-age population.

High immigration partially offsets OADR growth because immigrants are
typically working-age adults, but the offset is incomplete: at current
immigration levels, the OADR will still reach approximately 35--38 by 2050,
compared to 40--45 under zero immigration.
\end{defbox}

The demographic transition creates several distinct labour market effects:

\begin{enumerate}
  \item \textbf{Labour shortages in care sectors:} The aging population
        dramatically increases demand for healthcare workers, personal support
        workers, and long-term care staff --- precisely the occupations facing
        the most severe shortages. Canada needs to hire approximately 1 million
        healthcare workers over the next decade simply to maintain current
        service levels for a growing senior population.

  \item \textbf{Skilled trades shortages:} Construction and infrastructure
        maintenance require skilled tradespeople whose training pipeline
        (apprenticeships, technical schools) has consistently underproduced.
        The housing expansion needed to address the affordability crisis
        (Chapter~\ref{ch:housing}) requires more construction workers than
        Canada is currently training.

  \item \textbf{Loss of experience:} When a baby boom machinist like Jerome
        retires, he takes decades of tacit knowledge with him. The
        ``experience premium'' --- the wage differential that rewards years of
        accumulated skill --- is real and measurable. Replacing experienced
        workers with younger or immigrant workers takes years of on-the-job
        learning to replicate.

  \item \textbf{Upward pressure on wages:} As labour becomes relatively
        scarcer, the bargaining power of workers increases. The 2021--2022
        labour market tightness produced the strongest wage growth in
        decades --- particularly in trades, transport, and care work.
        The aging-driven tightness is likely structural rather than cyclical:
        it will persist even as the Bank of Canada slows the economy.
\end{enumerate}

%==========================================================
\section{Case Studies}
\label{sec:ch4_casestudies}
%==========================================================

\begin{casestudy}{The Canadian Labour Market Through COVID-19 and Beyond}

\textbf{Background:} No event in recent Canadian history tested labour
market institutions more severely than the COVID-19 pandemic. The shutdown
of April 2020 --- when over 2 million jobs disappeared in a single month ---
was followed by one of the fastest recoveries in Canadian economic history.
Understanding why illuminates both the resilience and the fault lines
of Canada's labour market.

\textbf{The Collapse:} In April 2020, the unemployment rate hit 13.7\% ---
the highest since the Great Depression. But this understates the disruption:
millions of Canadians who were still employed worked zero hours (treated
as employed-but-absent by the LFS), and the participation rate fell to
62.4\% as many stopped searching. The \textit{actual} share of Canadians
with zero employment income in April 2020 was closer to 25\%.

The collapse was highly concentrated in specific sectors: accommodation
and food services (lost over 40\% of employment); arts and entertainment
(over 35\%); retail trade (over 20\%). These sectors share a key feature:
they require in-person contact and cannot be substituted with remote work.
They also share a demographic profile: disproportionately young, female,
and low-wage workers. The COVID recession was dramatically unequal.

\textbf{The Recovery:} Canada's employment recovery was faster than most
advanced economies. By September 2021, total employment had fully recovered
its pre-pandemic level. By mid-2022, the unemployment rate had fallen to
4.9\% --- a 50-year low. Three factors explain the speed of recovery:
(1) federal income supports (CERB, CEWS) prevented permanent worker
detachment from employers; (2) pent-up consumer demand was released rapidly
when restrictions lifted; and (3) the Bank of Canada maintained
ultra-accommodative monetary conditions through 2021.

\textbf{The Residual:} Full recovery in aggregate conceals persistent
scars. Participation rates for some groups --- particularly prime-age
women (25--54) who took on caregiving responsibilities during school
closures --- remained below pre-COVID levels into 2023. Long-COVID
has kept an estimated 50{,}000--100{,}000 Canadians out of the
labour force. And the Beveridge Curve's outward shift (Figure~\ref{fig:beveridge})
shows that the labour market is less efficient at matching than before
the pandemic --- a structural consequence of the COVID-era disruption
that will take years to fully unwind.
\end{casestudy}

%----------------------------------------------------------
\begin{casestudy}{The Gig Economy and Platform Work in Canada}

\textbf{Background:} Between 2015 and 2024, the number of Canadians
earning income primarily through digital platforms --- Uber, Lyft, DoorDash,
Instacart, TaskRabbit, and similar apps --- grew from negligible to
approximately 800{,}000 to 1{,}000{,}000 workers, representing roughly
4--5\% of the employed workforce. Platform work has attracted enormous
policy attention because it straddles the line between employment and
self-employment in ways that existing labour law was not designed to handle.

\textbf{The Classification Problem:} Under Canadian provincial employment
standards legislation, the key distinction is between \textit{employees}
(entitled to minimum wage, overtime, paid vacation, EI access, and workplace
protections) and \textit{independent contractors} (responsible for their
own taxes, benefits, and business expenses, but generally free to set
their own terms). Platforms classify their workers as independent contractors,
arguing that the flexibility and algorithmic matching that platforms provide
constitutes a genuine market intermediary relationship rather than employment.

\textbf{Workers' perspective:} Platform workers report an effective hourly
wage that, after deducting vehicle costs, fuel, insurance (higher for
commercial use), phone plans, and other business expenses, frequently falls
below the minimum wage. In Ontario, the minimum wage in 2025 is \$17.20/hour;
studies estimate that delivery workers' net hourly income is typically
\$12--\$16/hour after vehicle expenses.

\textbf{Policy responses:} Ontario's 2021 Bill 88 introduced a ``digital
platform worker'' category with some but not all employee protections:
minimum earnings per hour of work (not per assigned period), transparency
requirements on pay calculation, and the right to dispute deactivations.
It stopped short of full reclassification. California's Proposition 22
(2020) similarly created a third-category status for app-based workers.

\textbf{Economic assessment:} The gig economy creates real value by
improving the utilisation of underused assets (vehicles) and providing
flexible income to people who genuinely prefer episodic work. But it also
creates real externalities: the misclassification of employees as contractors
shifts costs from platforms to workers and to public programs (EI, healthcare,
CPP) funded by other workers and employers. The efficient policy
distinguishes between \textit{genuine} flexible self-employment and
misclassified employment --- a distinction that requires bright-line
legal tests that Canadian employment law has not yet fully developed.
\end{casestudy}

%----------------------------------------------------------
\begin{casestudy}{Northern BC: Labour Shortages at the Margin of the Economy}

\textbf{Background:} The University of Northern British Columbia sits in
a region that illustrates labour market dynamics largely invisible in
national statistics. Northern BC faces simultaneous labour surpluses in
some occupations (laid-off mill workers when lumber markets soften) and
acute shortages in others (healthcare workers, skilled tradespeople,
early childhood educators).

\textbf{The Healthcare Worker Shortage:} Northern Health Authority,
which serves over 300{,}000 people across 600{,}000 square kilometres,
consistently operates with nurse vacancy rates of 15--25\%. The shortage
has several causes: geographic remoteness makes recruitment from southern
Canada or internationally difficult; housing costs in resource communities
have risen dramatically with the LNG construction boom; wage competition
from Lower Mainland health authorities and Alberta offers is intense;
and the combination of demanding work conditions and geographic isolation
produces high turnover.

The economic consequence is both a welfare loss (residents receive
suboptimal healthcare) and a fiscal one (travel medicine and locum
staffing costs are far higher per patient than permanent staff costs).

\textbf{The Trades Shortage:} LNG Canada's construction at Kitimat and
related pipeline infrastructure require more skilled tradespeople than
Northern BC has historically trained. The Northern Opportunities program
--- a partnership between School District 57 (Prince George), UNBC, the
College of New Caledonia, and the Northern Development Initiative Trust
--- attempts to address the shortage by channelling Northern BC youth
into pre-apprenticeship training. The program illustrates a general
principle: effective labour market policy is geographically specific,
sector-specific, and requires coordination across educational, government,
and industry actors.

\textbf{The Policy Lesson:} The Beveridge Curve analysis (Section~\ref{sec:beveridge_curve})
showed that Canada's post-COVID outward shift reflects skills and geographic
mismatches. Northern BC is a microcosm of this: vacancies in healthcare and
trades, potential workers in resource sectors facing cyclical disruption.
The solution is not aggregate demand stimulus but targeted training,
credential recognition, geographic incentives, and housing investment
in underserved communities.
\end{casestudy}

%==========================================================
\begin{policydebate}{Should Canada Raise the Federal Minimum Wage to \$20 Per Hour?}

\textbf{The Issue:} The federal minimum wage (covering workers in federally
regulated industries) reached \$17.30/hour in 2025, indexed to CPI.
Labour advocates propose a target of \$20/hour --- a ``living wage''
benchmark derived from estimates of the income needed to cover basic costs
in major Canadian cities. Business groups oppose the increase, citing
potential job losses and cost pressures on small businesses.

\textbf{The Economic Case For \$20/Hour:}
\begin{itemize}
  \item The federal minimum wage applies to a relatively narrow group
        (banking, telecommunications, airlines, interprovincial transport
        --- fewer than 1 million workers). A \$20 federal minimum would
        have modest aggregate employment effects.
  \item The monopsony evidence (Card, Krueger, and Canadian follow-up studies)
        suggests moderate minimum wage increases have small negative employment
        effects and significant positive wage effects for affected workers.
  \item A \$17.30 minimum is below the estimated cost of living in Vancouver
        (\$22/hour), Toronto (\$21/hour), and other major cities. Real-terms
        adequacy requires periodic above-inflation increases.
  \item Higher minimum wages reduce reliance on government transfers,
        shifting income support costs from taxpayers to employers who
        benefit from low-cost labour.
\end{itemize}

\textbf{The Economic Case Against:}
\begin{itemize}
  \item A 15.6\% increase from \$17.30 to \$20.00 exceeds the range
        supported by the monopsony-compatible evidence base. The studies
        finding small employment effects generally examine 5--15\% increases.
        Larger jumps in competitive markets may produce the job losses
        the standard model predicts.
  \item The federal minimum wage creates competitive disparities between
        federally regulated firms (subject to the higher wage) and provincially
        regulated firms (subject to lower provincial minimums). This can
        distort industry structure.
  \item Small businesses, particularly in lower-wage markets, face genuine
        cash-flow constraints. A rapid increase may force closures rather than
        wage adjustments.
  \item A geographically uniform \$20 minimum applied to rural and
        northern communities with lower cost-of-living may exceed the
        competitive equilibrium wage in those markets, producing real
        employment losses there even if it is appropriate in Toronto.
\end{itemize}

\textbf{The Evidence:}
The Bank of Canada has estimated that each 10\% minimum wage increase
reduces youth employment by approximately 0.5\% and total employment by
approximately 0.2\%, with larger effects in lower-wage markets. The PBO
estimates that the proposed \$20 federal minimum would: raise wages for
approximately 200{,}000--250{,}000 workers; result in net employment effects
of approximately $-$10{,}000 to $-$25{,}000 jobs; and reduce government
expenditure on income support programs by approximately \$300--\$500 million
per year as workers earn above benefit thresholds.

\textbf{The Policy Synthesis:} The economic evidence supports moderate
minimum wage increases to a ``living wage'' benchmark over time, with
CPI indexation to prevent real erosion. A stepped increase from \$17.30
to \$20.00 over three years --- rather than a one-time jump --- would
allow labour markets to adjust more gradually and reduce the risk of
discontinuous employment effects. The geographic heterogeneity of Canadian
labour markets also argues for allowing provincial variation rather than
imposing a single federal standard on communities with very different
wage structures.
\end{policydebate}

%==========================================================
\begin{everydaylife}{Understanding Your Pay Stub: Labour Economics in Practice}

Your pay stub is a document in applied labour economics. Every line
reflects a principle from this chapter.

\textbf{Gross wages} represent your nominal wage times hours worked ---
the payment for your marginal product of labour at the market rate
(competitive) or the rate your employer sets unilaterally (monopsony).

\textbf{CPP deductions} (Canada Pension Plan): You contribute 5.95\%
of eligible earnings, matched by your employer. This is a mandatory
contribution to a defined-benefit public pension --- protecting you
against the risk that you save too little for retirement. Gig workers
classified as self-employed pay both the employee and employer share
(11.9\%).

\textbf{EI deductions} (Employment Insurance): You contribute 1.66\%
of insurable earnings. EI provides income replacement if you become
unemployed through no fault of your own --- a form of insurance against
frictional and cyclical unemployment. Platform workers classified as
contractors do not pay into EI and cannot claim it.

\textbf{Income tax withholding}: The progressive income tax system means
your employer withholds tax at your marginal rate each pay period.
The progressivity of the system (higher earners pay a higher marginal
rate) is a key tool for reducing inequality --- which is why economists
carefully track both pre-tax and after-tax income inequality.

\textbf{Union dues} (if applicable): These fund your union's collective
bargaining, legal services, and strike fund. The union wage premium ---
typically 15--20\% --- is the economic return on this investment.
Non-unionized workers receive no such institutional representation.

What your pay stub does not show: the non-wage benefits your employer
provides (pension contributions, health premiums), the employer's CPP
and EI contributions on your behalf (which equal your own contributions),
and the value of job security, training, and safe working conditions that
vary enormously between standard and non-standard employment relationships.
\end{everydaylife}

%==========================================================
\section*{Chapter Summary}
%==========================================================
\addcontentsline{toc}{section}{Chapter Summary}

\begin{itemize}
  \item In a \textbf{competitive labour market}, the equilibrium wage equals
        the value of the marginal product of labour ($W^* = P \times \text{MP}_L$).
        A minimum wage above $W^*$ creates unemployment. In a
        \textbf{monopsony}, the wage is set below the competitive level and
        employment is restricted; a minimum wage up to $W^*$ can increase
        both wages and employment.

  \item The three key labour force statistics are the
        \textbf{unemployment rate} ($U/LF$), the
        \textbf{labour force participation rate} ($LF/\text{WAP}$), and the
        \textbf{employment rate} ($E/\text{WAP}$). Each captures a different
        dimension of labour market performance; all three are needed for a
        complete picture.

  \item \textbf{Frictional unemployment} (search-related), \textbf{structural
        unemployment} (skills/geographic mismatch), and \textbf{cyclical
        unemployment} (demand deficiency) have different causes and require
        different policy responses. Canada's NAIRU is approximately
        5.5--6.0\%.

  \item The \textbf{Beveridge Curve} plots vacancy rates against unemployment.
        Canada's post-COVID outward shift indicates reduced labour market
        matching efficiency --- a structural problem requiring training and
        credential recognition reforms rather than demand stimulus.

  \item Canada's \textbf{wage-productivity gap} has widened since the 1970s:
        labour productivity has grown ~75\% since 1976 while real wages
        have grown only ~32\%. Declining unionization, rising employer
        market power, and trade competition are leading explanations.

  \item \textbf{Union coverage} in Canada's private sector has declined
        to approximately 14\%, while public-sector coverage remains above
        70\%. The union wage premium of 15--20\% and the benefit premium
        document the real costs of declining private-sector unionization.

  \item \textbf{Non-standard employment} (temporary, own-account
        self-employment, gig/platform) accounts for approximately 29\%
        of Canadian employment and offers significantly lower wages, fewer
        benefits, and less job security than permanent full-time work.

  \item \textbf{Demographic aging} --- baby boom retirements, below-replacement
        fertility, and rising dependency ratios --- is creating structural
        labour shortages in healthcare, skilled trades, and early childhood
        education that persist regardless of the business cycle.

  \item \textbf{Provincial minimum wages} range from \$15.00 (AB, SK) to
        \$17.40 (BC). The empirical evidence from Canada and the US suggests
        that moderate minimum wage increases have small negative employment
        effects and meaningful wage gains for affected workers, particularly
        in monopsonistic labour markets.
\end{itemize}

%==========================================================
\section*{Key Terms}
%==========================================================
\addcontentsline{toc}{section}{Key Terms}

\begin{description}[style=nextline, leftmargin=3.5cm]
  \item[\term{Beveridge Curve}] A plot of the job vacancy rate against
        the unemployment rate; its position reflects labour market
        matching efficiency.
  \item[\term{Collective bargaining}] Negotiation between a union
        representing workers and an employer over wages, benefits, and
        conditions.
  \item[\term{Cyclical unemployment}] Unemployment arising from insufficient
        aggregate demand during an economic downturn.
  \item[\term{Dependency ratio (old-age)}] Population 65+ per 100 persons
        aged 15--64; measures the fiscal and labour-supply burden of aging.
  \item[\term{Discouraged workers}] Persons who want work but have stopped
        actively searching and are therefore excluded from the unemployment
        count.
  \item[\term{Employment rate}] Employed persons as a share of the working-age
        population: $E/\text{WAP} \times 100$.
  \item[\term{Frictional unemployment}] Temporary unemployment arising from
        the normal search-and-matching process.
  \item[\term{Gig economy}] Economic activity mediated by digital platforms
        that match workers with short-term tasks or customers.
  \item[\term{Labour force}] Employed plus unemployed persons; $LF = E + U$.
  \item[\term{Labour Force Participation Rate (LFPR)}] Labour force as a
        share of working-age population: $LF/\text{WAP} \times 100$.
  \item[\term{Labour share of income}] Wages and salaries as a fraction
        of national income; has declined over four decades in Canada.
  \item[\term{Minimum wage}] A legislated floor on hourly wages.
  \item[\term{Monopsony}] A labour market with a single (or dominant)
        employer, who can set wages below the competitive level.
  \item[\term{NAIRU}] Non-Accelerating Inflation Rate of Unemployment;
        the unemployment rate at potential output, consisting of frictional
        and structural components; $\approx 5.5$--$6.0\%$ in Canada.
  \item[\term{Non-standard employment}] Employment arrangements that deviate
        from permanent full-time work: temporary, part-time, self-employed,
        or gig.
  \item[\term{Structural unemployment}] Unemployment arising from skills
        or geographic mismatch between available workers and available jobs.
  \item[\term{Unemployment rate}] Unemployed persons as a share of the
        labour force: $U/LF \times 100$.
  \item[\term{Union wage premium}] The percentage difference in wages between
        comparable unionized and non-unionized workers; approximately 15--20\%
        in Canada.
  \item[\term{Value of the marginal product of labour (VMP$_L$)}] The market
        value of the additional output produced by one more unit of labour:
        $\text{VMP}_L = P \times \text{MP}_L$.
  \item[\term{Wage-productivity gap}] The divergence between real wage growth
        and labour productivity growth since the 1970s.
  \item[\term{Working-age population (WAP)}] All persons aged 15 and over;
        the denominator for the employment rate and LFPR.
\end{description}

%==========================================================
\section*{Review Questions}
%==========================================================
\addcontentsline{toc}{section}{Review Questions}

\begin{enumerate}

\item \textbf{(Concept)} Explain the difference between the competitive
and monopsony models of the labour market. How does each model predict
the employment effect of a minimum wage set above the competitive equilibrium?
Why is the Card-Krueger finding (no employment loss from a moderate minimum
wage increase) consistent with the monopsony model but not the competitive model?

\item \textbf{(Calculation)} In a hypothetical province: working-age
population = 3{,}500{,}000; employed = 2{,}100{,}000; unemployed = 130{,}000;
not in labour force = 1{,}270{,}000.
\begin{enumerate}[(a)]
  \item Calculate the unemployment rate, LFPR, and employment rate.
  \item If 50{,}000 discouraged workers re-enter the labour force and
        20{,}000 find jobs, recalculate all three statistics.
  \item Why did the unemployment rate change in (b) even though 20{,}000
        people found jobs?
\end{enumerate}

\item \textbf{(Types of unemployment)} Classify each of the following
as frictional, structural, or cyclical unemployment, and explain your
reasoning:
\begin{enumerate}[(a)]
  \item A software developer in Waterloo who left her job to find one
        with a higher salary and is currently searching.
  \item A coal mine worker in Nova Scotia who lost his job when the
        mine closed due to the carbon transition.
  \item A retail sales associate in Vancouver who was laid off when
        consumer spending fell sharply during a recession.
  \item A nurse in Prince George who cannot find a job because she
        earned her credentials in the Philippines and they are not
        yet recognised in BC.
\end{enumerate}

\item \textbf{(Beveridge Curve)} Canada's job vacancy rate rose from
2.2\% in 2019 to 4.5\% in 2022 even though the unemployment rate fell
from 5.7\% to 5.3\%. Explain what this means for the Beveridge Curve
and for the interpretation of the 2022 labour market tightness. Is this
evidence of structural or cyclical problems?

\item \textbf{(Wage-productivity gap)} Explain two mechanisms through
which declining union coverage could contribute to a widening
wage-productivity gap. For each mechanism, describe a policy intervention
that could partially reverse it.

\item \textbf{(Policy)} A provincial government is considering requiring
all gig/platform workers to be reclassified as employees. Using the
concepts in this chapter, what would you predict are the effects on:
(a) platform worker wages and benefits; (b) platform company costs and
pricing; (c) the number of gig work opportunities available; and (d) the
measurement of unemployment and LFPR?

\end{enumerate}

%==========================================================
\section*{Quantitative Problems}
%==========================================================
\addcontentsline{toc}{section}{Quantitative Problems}

\begin{enumerate}

\item \textbf{Labour Force Accounting.}
Canada's Labour Force Survey for January 2024 reported: working-age
population 34{,}150{,}000; employed 20{,}580{,}000; unemployed 1{,}380{,}000.
\begin{enumerate}[(a)]
  \item Calculate the labour force, unemployment rate, LFPR, and
        employment rate.
  \item If 200{,}000 discouraged workers re-entered the labour force
        in February 2024 (without finding jobs), recalculate all four
        statistics.
  \item If those same 200{,}000 eventually all found employment by March,
        recalculate again.
  \item Explain the sequence of changes in the unemployment rate across
        these three scenarios. What does this illustrate about the
        limitations of the unemployment rate as a welfare measure?
\end{enumerate}

\item \textbf{Minimum Wage and Employment Effects.}
Suppose a city's fast-food labour market can be described as follows:
\begin{align*}
  L^d &= 15{,}000 - 600W \quad \text{(demand for workers, hours/week)}\\
  L^s &= 3{,}000 + 400W \quad \text{(supply of workers, hours/week)}
\end{align*}
\begin{enumerate}[(a)]
  \item Find the competitive equilibrium wage $W^*$ and employment $L^*$.
  \item If the provincial government sets a minimum wage of \$16.00/hour,
        how many workers are employed? How many are unemployed?
  \item Calculate the unemployment generated by the minimum wage as a
        number and as a percentage of the labour force at the minimum wage.
  \item Now suppose the market is actually monopsonistic. The monopsonist's
        marginal labour cost is $\text{MLC} = -3{,}000 + 5W$... [adapt
        from supply curve: if $W = (L^s - 3{,}000)/400$, then
        $\text{MLC} = (2L^s - 3{,}000)/400$]. Find the monopsony wage
        and employment, and show that the \$16.00 minimum wage in this
        context increases employment rather than decreasing it.
\end{enumerate}

\item \textbf{Wage-Productivity Calculation.}
Use the following Canadian data:

\begin{center}
\begin{tabular}{lrrr}
\toprule
Year & Nominal avg.\ weekly earnings (\$) & CPI (2002=100) & Labour productivity index \\
\midrule
1976 & 230.40 & 38.1 & 100.0 \\
1990 & 543.20 & 78.4 & 119.0 \\
2000 & 685.90 & 95.4 & 142.0 \\
2010 & 889.20 & 116.5 & 157.0 \\
2024 & 1{,}245.80 & 168.2 & 173.0 \\
\bottomrule
\end{tabular}
\end{center}

\begin{enumerate}[(a)]
  \item Calculate real average weekly earnings for each year
        (in 1976 dollars), indexed to 1976 = 100.
  \item Calculate the percentage growth in real wages and in labour
        productivity between 1976 and 2024.
  \item What is the ``wage-productivity gap'' in percentage points?
  \item If real wages had grown at the same rate as labour productivity
        since 1976, what would average real weekly earnings be in 2024
        (in 2024 dollars)?
  \item Using Equation~\ref{eq:vmp}, if labour productivity rose 73\%
        but real wages rose only 32\%, what must have happened to the
        ``$P$'' or to the competitive conditions in the labour market?
        Propose at least two explanations.
\end{enumerate}

\item \textbf{Dependency Ratio Projection.}
Statistics Canada projects the following population for Canada:

\begin{center}
\begin{tabular}{lrrr}
\toprule
Year & Population 65+ (M) & Population 15--64 (M) & OADR \\
\midrule
2020 &  6.6 & 24.0 & ? \\
2030 &  9.3 & 25.5 & ? \\
2040 & 11.1 & 26.0 & ? \\
2050 & 12.2 & 26.8 & ? \\
\bottomrule
\end{tabular}
\end{center}

\begin{enumerate}[(a)]
  \item Calculate the old-age dependency ratio for each year
        using Equation~\ref{eq:oadr}.
  \item If each retiree draws an average of \$25{,}000/year in
        OAS/GIS payments, calculate the total OAS/GIS cost in each year.
  \item If this cost is financed by a payroll tax split equally between
        workers and employers, calculate the implied tax rate (as a share
        of an average annual wage of \$65{,}000) in each year.
  \item What does this imply about the fiscal sustainability of Canada's
        public pension system and the potential need for immigration as
        an offset?
\end{enumerate}

\item \textbf{Non-Standard Employment Calculation.}
Using Table~\ref{tab:employment_comparison}:
\begin{enumerate}[(a)]
  \item Calculate the hourly wage gap (\$) and percentage gap between
        permanent full-time and each of the three non-standard categories.
  \item If a platform/gig worker earns \$18.50/hour but has the following
        deductible costs: vehicle depreciation \$2.40/hour, fuel
        \$1.60/hour, insurance premium \$1.20/hour, phone \$0.40/hour,
        what is the net hourly income? How does this compare to Ontario's
        minimum wage?
  \item A full-time permanent worker earning \$31.50/hour receives employer
        pension contributions of 5\% and health benefits valued at
        \$3.20/hour. What is the total effective compensation?
  \item Calculate the ``total compensation gap'' between the gig worker's
        net income from (b) and the permanent worker's total compensation
        from (c).
\end{enumerate}

\item \textbf{Data Exercise.}
Visit Statistics Canada (\url{www.statcan.gc.ca}) and find the most recent
Labour Force Survey release (Table 14-10-0287-01, seasonally adjusted).
\begin{enumerate}[(a)]
  \item Report Canada's current unemployment rate, employment rate, and LFPR.
  \item Break the unemployment rate down by: age group (15--24, 25--54,
        55+), sex (male, female), and educational attainment.
        Which group has the highest unemployment rate? The lowest?
  \item Find the job vacancy rate from the Job Vacancy and Wage Survey
        (Table 14-10-0325-01). Using Figure~\ref{fig:beveridge}, identify
        where the current data point sits relative to the pre-COVID curve.
        What does this imply about current matching efficiency?
  \item How does Canada's unemployment rate compare to the current US and
        OECD average rates? (Source: OECD.Stat, Labour Force Statistics.)
\end{enumerate}

\end{enumerate}

%==========================================================
\section*{Essay Questions}
%==========================================================
\addcontentsline{toc}{section}{Essay Questions}

\begin{enumerate}

\item \textbf{(600--800 words)} The vignette that opened this chapter
described Jerome (a unionized machinist) and Kayla (a gig logistics worker)
as inhabiting very different labour markets in the same city. Using the
economics of this chapter --- wage determination, non-standard employment,
monopsony, and the minimum wage debate --- explain the institutional and
market forces that produce such different outcomes for workers in the same
economy. What policies would you recommend to reduce the gap? What trade-offs
do your recommendations involve?

\item \textbf{(500--700 words)} The wage-productivity gap has widened in
Canada since the 1970s. Some economists argue this is primarily due to
declining unionization; others argue it reflects technological change and
globalisation. Present the best argument for each explanation, evaluate the
empirical evidence, and then state which explanation you find more convincing
and why. Your answer should connect explicitly to Figure~\ref{fig:wage_productivity}
and Figure~\ref{fig:unionization}.

\item \textbf{(400--600 words)} Canada's labour market faces two simultaneous
structural challenges: an aging population reducing labour supply in many
sectors, and an outward Beveridge Curve shift suggesting reduced matching
efficiency. Are these problems complementary (both reflecting a failure
of skills supply to meet demand) or in tension (one might offset the other)?
Evaluate the role of immigration as a potential solution to both challenges,
and identify at least one condition under which immigration could worsen
the matching efficiency problem rather than improve it.

\end{enumerate}

%==========================================================
\section*{Discussion Questions}
%==========================================================
\addcontentsline{toc}{section}{Discussion Questions}

\begin{enumerate}

\item Should Canada have a single national minimum wage rather than
10 different provincial minimums? What are the economic arguments for
standardisation? What are the arguments for provincial variation? How
does your answer relate to the monopsony model and the geographic
diversity of Canadian labour markets?

\item Jerome's unionized position pays \$34.20/hour, has a defined-benefit
pension, and includes comprehensive benefits. Kayla's gig work pays
effectively \$12--\$16/hour with no benefits or security. Is this
difference entirely the result of market forces reflecting different
marginal products of labour? Or does institutional variation (unions,
employment classifications) explain part of the gap? If the latter,
is the current distribution of labour market ``rents'' economically
efficient? Is it equitable?

\item Canada's old-age dependency ratio is rising rapidly. The three
main policy responses are: (a) increasing immigration of working-age
adults; (b) raising the retirement age; and (c) increasing the CPP
contribution rate to build up the fund. What are the economic trade-offs
of each approach? Which combination would you recommend, and why?

\item A provincial government is considering legislation to require
platforms (Uber, DoorDash, etc.) to classify all workers performing
more than 20 hours of platform work per month as ``employees.'' Predict
the effects of this policy on: (a) gig worker wages and benefits;
(b) ride-share and delivery prices; (c) total platform employment;
and (d) the composition of the gig workforce. Who wins and who loses?

\end{enumerate}

% CHAPTER STUB — full content forthcoming
\chapter{Inequality}
\label{ch:inequality}
This chapter is under development.

% CHAPTER STUB — full content forthcoming
\chapter{Immigration}
\label{ch:immigration}
This chapter is under development.

% CHAPTER STUB — full content forthcoming
\chapter{Fiscal}
\label{ch:fiscal}
This chapter is under development.

% CHAPTER STUB — full content forthcoming
\chapter{Trade}
\label{ch:trade}
This chapter is under development.

% CHAPTER STUB — full content forthcoming
\chapter{Technology}
\label{ch:technology}
This chapter is under development.

% CHAPTER STUB — full content forthcoming
\chapter{Climate}
\label{ch:climate}
This chapter is under development.


%----------------------------------------------------------
% BACK MATTER
%----------------------------------------------------------
\backmatter
\pagestyle{fancy}

%==========================================================
%  GLOSSARY OF KEY TERMS
%  ECON 204 — Contemporary Economic Issues: A Canadian Perspective
%==========================================================

\chapter*{Glossary of Key Terms}
\addcontentsline{toc}{chapter}{Glossary}
\markboth{Glossary}{Glossary}

\noindent This glossary defines key economic terms introduced in this textbook. Terms are listed alphabetically, with cross-references to the chapter in which each term first appears.

\vspace{12pt}

%----------------------------------------------------------
% A
%----------------------------------------------------------
\noindent\textbf{Adaptation (climate):} Adjustments in economic activity, infrastructure, and behaviour to reduce the harm caused by climate change that is already occurring or locked in. \textit{See} Chapter~\ref{ch:climate}.

\vspace{6pt}
\noindent\textbf{Aggregate Demand (AD):} The total quantity of goods and services demanded across the entire economy at each price level; the sum of consumption, investment, government expenditure, and net exports. \textit{See} Chapter~\ref{ch:inflation}.

\vspace{6pt}
\noindent\textbf{Aggregate Supply (AS):} The total quantity of goods and services that firms in an economy are willing and able to produce at each price level. The short-run aggregate supply (SRAS) slopes upward; the long-run aggregate supply (LRAS) is vertical at potential output. \textit{See} Chapter~\ref{ch:inflation}.

\vspace{6pt}
\noindent\textbf{Automation:} The replacement of human labour with machines, software, or robots in the performance of specific tasks. \textit{See} Chapter~\ref{ch:technology}.

%----------------------------------------------------------
% B
%----------------------------------------------------------
\vspace{6pt}
\noindent\textbf{Bank of Canada:} Canada's central bank, established in 1934 as a Crown corporation. It conducts monetary policy with the primary objective of maintaining low and stable inflation within a 1--3\% target range. \textit{See} Chapter~\ref{ch:inflation}.

\vspace{6pt}
\noindent\textbf{Built-in inflation:} Inflation arising from the expectation of future inflation; wage and price increases made in anticipation of continued price rises, creating a self-fulfilling wage-price spiral. \textit{See} Chapter~\ref{ch:inflation}.

%----------------------------------------------------------
% C
%----------------------------------------------------------
\vspace{6pt}
\noindent\textbf{Canada Carbon Rebate (CCR):} The payment made by the federal government to eligible Canadian households to return revenue from the consumer fuel charge. Paid quarterly as a flat amount per household. \textit{See} Chapter~\ref{ch:climate}.

\vspace{6pt}
\noindent\textbf{Canada Child Benefit (CCB):} A non-taxable, income-tested monthly benefit for families with children under 18. Launched in 2016 to replace previous child benefit programs; credited with significantly reducing child poverty in Canada. \textit{See} Chapter~\ref{ch:inequality}.

\vspace{6pt}
\noindent\textbf{Cap-and-trade:} An emissions reduction policy that sets a maximum (cap) on total greenhouse gas emissions and allows firms to buy and sell emissions permits. \textit{See} Chapter~\ref{ch:climate}.

\vspace{6pt}
\noindent\textbf{Capital gains:} The profit realized from selling an asset (real estate, shares, bonds) for more than its purchase price. \textit{See} Chapter~\ref{ch:inequality}.

\vspace{6pt}
\noindent\textbf{Carbon leakage:} The displacement of emissions-intensive production to jurisdictions with lower or no carbon prices as a result of domestic carbon pricing policy. \textit{See} Chapter~\ref{ch:climate}.

\vspace{6pt}
\noindent\textbf{Carbon tax:} A direct fee levied on greenhouse gas emissions, proportional to the carbon content of fossil fuels burned. A Pigouvian tax designed to internalize the externality cost of carbon emissions. \textit{See} Chapter~\ref{ch:climate}.

\vspace{6pt}
\noindent\textbf{Central bank independence:} The principle that a central bank has the authority to set monetary policy without direct political interference from the elected government. \textit{See} Chapter~\ref{ch:inflation}.

\vspace{6pt}
\noindent\textbf{Climate mitigation:} Actions that reduce greenhouse gas emissions or enhance carbon sinks, thereby slowing or reducing future climate change. \textit{See} Chapter~\ref{ch:climate}.

\vspace{6pt}
\noindent\textbf{Consumer Price Index (CPI):} A measure of the average change over time in the prices paid by urban consumers for a representative fixed basket of goods and services, produced monthly by Statistics Canada. \textit{See} Chapter~\ref{ch:inflation}.

\vspace{6pt}
\noindent\textbf{Core inflation:} Measures of inflation that exclude or reduce the influence of volatile food and energy prices, providing a cleaner signal of underlying inflationary trends. The Bank of Canada uses CPI-trim, CPI-median, and CPI-common. \textit{See} Chapter~\ref{ch:inflation}.

\vspace{6pt}
\noindent\textbf{Cost-push inflation:} Inflation caused by an increase in the cost of inputs to production (e.g., oil, wages), which reduces aggregate supply. \textit{See} Chapter~\ref{ch:inflation}.

%----------------------------------------------------------
% D
%----------------------------------------------------------
\vspace{6pt}
\noindent\textbf{Data network effects:} The process by which a platform with more users generates more data, improving its AI models and service quality, which attracts more users---a self-reinforcing competitive advantage. \textit{See} Chapter~\ref{ch:technology}.

\vspace{6pt}
\noindent\textbf{Deflation:} A sustained decrease in the general price level. Deflation can trap an economy in a downward spiral of postponed spending and rising real debt burdens. \textit{See} Chapter~\ref{ch:inflation}.

\vspace{6pt}
\noindent\textbf{Demand-pull inflation:} Inflation caused by an increase in aggregate demand, resulting in ``too much money chasing too few goods.'' \textit{See} Chapter~\ref{ch:inflation}.

\vspace{6pt}
\noindent\textbf{Development charges:} Fees levied by municipal governments on builders and developers to fund the infrastructure (roads, water, parks, transit) required to support new residential and commercial development. \textit{See} Chapter~\ref{ch:housing}.

%----------------------------------------------------------
% E
%----------------------------------------------------------
\vspace{6pt}
\noindent\textbf{Exclusionary zoning:} Land-use regulations that prohibit multi-unit residential buildings (apartments, duplexes, townhouses) in residential zones, permitting only detached single-family homes. A primary driver of housing supply constraints in Canadian cities. \textit{See} Chapter~\ref{ch:housing}.

%----------------------------------------------------------
% F
%----------------------------------------------------------
\vspace{6pt}
\noindent\textbf{Fiscal policy:} The use of government spending and taxation to influence aggregate demand and economic output. Controlled by the elected government through the Department of Finance Canada. \textit{See} Chapter~\ref{ch:introduction}.

\vspace{6pt}
\noindent\textbf{Fisher equation:} The relationship between nominal interest rates, real interest rates, and inflation: $r \approx i - \pi$, where $r$ is the real rate, $i$ is the nominal rate, and $\pi$ is inflation. \textit{See} Chapter~\ref{ch:inflation}.

\vspace{6pt}
\noindent\textbf{Free Prior and Informed Consent (FPIC):} A principle under the United Nations Declaration on the Rights of Indigenous Peoples (UNDRIP) stating that Indigenous peoples have the right to give or withhold consent to resource development on their traditional territories. \textit{See} Chapter~\ref{ch:inequality}.

\vspace{6pt}
\noindent\textbf{Frictional unemployment:} Short-term unemployment arising from the process of matching workers with jobs; inherent in any dynamic labour market. \textit{See} Chapter~\ref{ch:labour}.

%----------------------------------------------------------
% G
%----------------------------------------------------------
\vspace{6pt}
\noindent\textbf{Gender wage gap:} The difference between median earnings of men and women in the labour market, typically expressed as women's earnings as a percentage of men's. \textit{See} Chapter~\ref{ch:labour}.

\vspace{6pt}
\noindent\textbf{General-purpose technology (GPT):} An innovation that is pervasive across many sectors, capable of continuous improvement, and able to generate complementary innovations. Examples include steam power, electricity, computers, and AI. \textit{See} Chapter~\ref{ch:technology}.

\vspace{6pt}
\noindent\textbf{Gig economy:} A labour market characterized by short-term contracts and freelance work, often mediated by digital platforms. Workers are typically classified as independent contractors. \textit{See} Chapter~\ref{ch:labour}.

\vspace{6pt}
\noindent\textbf{Gini coefficient:} A summary measure of income (or wealth) inequality derived from the Lorenz curve. Ranges from 0 (perfect equality) to 1 (perfect inequality). \textit{See} Chapter~\ref{ch:inequality}.

\vspace{6pt}
\noindent\textbf{Green industrial policy:} Government interventions (subsidies, tax credits, public investment, mandates) designed to accelerate the deployment of clean energy technologies and the transition to a low-carbon economy. \textit{See} Chapter~\ref{ch:climate}.

\vspace{6pt}
\noindent\textbf{Gross Domestic Product (GDP):} The total market value of all final goods and services produced within a country's borders in a given period (typically one year). \textit{See} Chapter~\ref{ch:introduction}.

%----------------------------------------------------------
% H
%----------------------------------------------------------
\vspace{6pt}
\noindent\textbf{Housing affordability:} The relationship between housing costs (prices, rents) and household income. Commonly measured by the price-to-income ratio or the share of income spent on shelter. \textit{See} Chapter~\ref{ch:housing}.

%----------------------------------------------------------
% I
%----------------------------------------------------------
\vspace{6pt}
\noindent\textbf{Imputed rent:} The estimated rental value of owner-occupied housing---the income an owner would receive by renting their property, or equivalently, the rent they save by owning. \textit{See} Chapter~\ref{ch:housing}.

\vspace{6pt}
\noindent\textbf{Income quintile:} One-fifth of the population ranked by income. \textit{See} Chapter~\ref{ch:inequality}.

\vspace{6pt}
\noindent\textbf{Inflation:} A sustained increase in the general price level of goods and services in an economy over time. \textit{See} Chapter~\ref{ch:inflation}.

\vspace{6pt}
\noindent\textbf{Inflation targeting:} A monetary policy framework in which the central bank publicly commits to maintaining inflation at a specific numerical target (in Canada, 2\% within a 1--3\% range). \textit{See} Chapter~\ref{ch:inflation}.

%----------------------------------------------------------
% K
%----------------------------------------------------------
\vspace{6pt}
\noindent\textbf{Kuznets hypothesis:} Simon Kuznets's 1955 prediction that income inequality follows an inverted-U pattern as an economy develops: rising during early industrialization and falling as development matures. Empirically challenged by the experience of advanced economies since 1980. \textit{See} Chapter~\ref{ch:inequality}.

%----------------------------------------------------------
% L
%----------------------------------------------------------
\vspace{6pt}
\noindent\textbf{Labour force participation rate:} The percentage of the working-age population (15+) that is either employed or actively seeking work. \textit{See} Chapter~\ref{ch:labour}.

\vspace{6pt}
\noindent\textbf{Labour productivity:} Output per unit of labour input; typically measured as real GDP per hour worked. \textit{See} Chapter~\ref{ch:labour}.

\vspace{6pt}
\noindent\textbf{Lorenz curve:} A graphical representation of the income (or wealth) distribution, plotting the cumulative share of total income against the cumulative share of the population ranked from poorest to richest. \textit{See} Chapter~\ref{ch:inequality}.

\vspace{6pt}
\noindent\textbf{Low Income Cut-Off (LICO):} Statistics Canada's original poverty measure, based on the share of income devoted to food, shelter, and clothing. No longer promoted as an official measure but still published. \textit{See} Chapter~\ref{ch:inequality}.

\vspace{6pt}
\noindent\textbf{Low Income Measure (LIM):} A relative poverty measure defined as income below 50\% of median adjusted household income; the standard OECD cross-country poverty measure. \textit{See} Chapter~\ref{ch:inequality}.

%----------------------------------------------------------
% M
%----------------------------------------------------------
\vspace{6pt}
\noindent\textbf{Marginal Revenue Product of Labour ($MRP_L$):} The additional revenue generated by hiring one more unit of labour; the basis for the firm's labour demand curve. \textit{See} Chapter~\ref{ch:labour}.

\vspace{6pt}
\noindent\textbf{Market Basket Measure (MBM):} Canada's official poverty line since 2019. Measures the cost of a specific basket of goods and services needed for a modest but adequate standard of living, with regional variation. \textit{See} Chapter~\ref{ch:inequality}.

\vspace{6pt}
\noindent\textbf{Minimum wage:} A government-mandated floor on wages below which employers cannot legally pay workers. Set provincially in Canada, with a separate federal rate for federally regulated industries. \textit{See} Chapter~\ref{ch:labour}.

\vspace{6pt}
\noindent\textbf{Monetary policy:} Actions taken by a central bank to influence interest rates, the money supply, and ultimately inflation and economic activity. In Canada, conducted by the Bank of Canada through the overnight rate. \textit{See} Chapter~\ref{ch:introduction}.

\vspace{6pt}
\noindent\textbf{Monetary policy transmission:} The process by which a change in the central bank's policy rate affects economic variables such as mortgage rates, the exchange rate, credit conditions, asset prices, and ultimately inflation. \textit{See} Chapter~\ref{ch:inflation}.

\vspace{6pt}
\noindent\textbf{Monopsony:} A labour market in which there is only one buyer of labour (one employer), giving that employer market power to set wages below the competitive level. \textit{See} Chapter~\ref{ch:labour}.

\vspace{6pt}
\noindent\textbf{Mortgage stress test:} OSFI's regulatory requirement (B-20) that mortgage borrowers must qualify at the higher of their contract rate plus 2 percentage points, or 5.25\%, ensuring borrowers have a buffer against rate increases. \textit{See} Chapter~\ref{ch:inflation}.

%----------------------------------------------------------
% N
%----------------------------------------------------------
\vspace{6pt}
\noindent\textbf{Natural rate of unemployment:} The rate of unemployment consistent with full employment---the sum of frictional and structural unemployment; the rate toward which the economy gravitates in the long run. \textit{See} Chapter~\ref{ch:labour}.

\vspace{6pt}
\noindent\textbf{Negative externality:} A cost imposed on a third party (not the buyer or seller) by an economic transaction that is not reflected in the market price. \textit{See} Chapter~\ref{ch:climate}.

\vspace{6pt}
\noindent\textbf{Network effects:} The phenomenon whereby the value of a product or service to any one user increases as more users adopt it. Network effects create tendencies toward winner-take-all market outcomes. \textit{See} Chapter~\ref{ch:technology}.

\vspace{6pt}
\noindent\textbf{Normative economics:} Economic analysis that makes value judgements about what \textit{should} happen, rather than what \textit{is} or \textit{will} happen. Contrasts with positive economics. \textit{See} Chapter~\ref{ch:introduction}.

%----------------------------------------------------------
% O
%----------------------------------------------------------
\vspace{6pt}
\noindent\textbf{Opportunity cost:} The value of the next-best alternative foregone when a choice is made. The true economic cost of any decision. \textit{See} Chapter~\ref{ch:introduction}.

\vspace{6pt}
\noindent\textbf{Overnight rate:} The Bank of Canada's policy interest rate---the target for the rate at which major financial institutions lend to each other overnight. Changes in the overnight rate transmit through the financial system to affect mortgage rates, business loans, and the exchange rate. \textit{See} Chapter~\ref{ch:inflation}.

%----------------------------------------------------------
% P
%----------------------------------------------------------
\vspace{6pt}
\noindent\textbf{Pigouvian tax:} A tax set equal to the marginal external cost of a negative externality, designed to align private incentives with social costs. Named after economist Arthur C. Pigou. The carbon tax is a Pigouvian tax. \textit{See} Chapter~\ref{ch:climate}.

\vspace{6pt}
\noindent\textbf{Platform economy:} A market structure in which a digital intermediary (platform) connects two or more groups of users (e.g., buyers and sellers, drivers and riders) and generates value primarily from network effects and data. \textit{See} Chapter~\ref{ch:technology}.

\vspace{6pt}
\noindent\textbf{Positive economics:} Economic analysis that describes and explains economic phenomena without making value judgements; empirical and testable claims about how the economy works. \textit{See} Chapter~\ref{ch:introduction}.

\vspace{6pt}
\noindent\textbf{Poverty:} A condition in which a household or individual lacks sufficient income or resources to meet a defined minimum standard of living. May be measured as absolute poverty (inability to meet basic needs) or relative poverty (income below a fraction of the median). \textit{See} Chapter~\ref{ch:inequality}.

\vspace{6pt}
\noindent\textbf{Price elasticity of supply:} The percentage change in quantity supplied resulting from a 1\% increase in price. Inelastic supply ($E_s < 1$) means supply responds weakly to price. \textit{See} Chapter~\ref{ch:housing}.

\vspace{6pt}
\noindent\textbf{Productivity paradox:} The observation that investment in information and communication technologies does not immediately produce measurable aggregate productivity gains---the lag between GPT adoption and full productivity impact. \textit{See} Chapter~\ref{ch:technology}.

%----------------------------------------------------------
% R
%----------------------------------------------------------
\vspace{6pt}
\noindent\textbf{Real interest rate:} The nominal interest rate adjusted for inflation: $r \approx i - \pi$. The real interest rate measures the true cost of borrowing in terms of purchasing power. \textit{See} Chapter~\ref{ch:inflation}.

\vspace{6pt}
\noindent\textbf{Rent control:} Government regulation limiting the rent a landlord may charge or the permissible annual rent increase for a residential unit. \textit{See} Chapter~\ref{ch:housing}.

\vspace{6pt}
\noindent\textbf{Routine vs.\ non-routine tasks:} In the task-based model of automation, routine tasks are governed by explicit, codifiable rules (most automatable); non-routine tasks require judgment, creativity, or physical dexterity in variable environments (less automatable). \textit{See} Chapter~\ref{ch:technology}.

%----------------------------------------------------------
% S
%----------------------------------------------------------
\vspace{6pt}
\noindent\textbf{Skill-biased technological change (SBTC):} Technological change that increases relative demand for skilled (typically higher-educated) workers and decreases relative demand for unskilled workers, widening the wage premium for education. \textit{See} Chapter~\ref{ch:inequality}.

\vspace{6pt}
\noindent\textbf{Social cost of carbon (SCC):} The economic value of the damage caused by emitting one additional tonne of CO$_2$-equivalent into the atmosphere, including all present and future climate harms. \textit{See} Chapter~\ref{ch:climate}.

\vspace{6pt}
\noindent\textbf{Stagflation:} The simultaneous occurrence of high inflation and low economic output (or recession). A particularly challenging macroeconomic condition to address with conventional policy tools. \textit{See} Chapter~\ref{ch:inflation}.

\vspace{6pt}
\noindent\textbf{Structural unemployment:} Long-term unemployment resulting from a mismatch between workers' skills and available jobs, or from geographic mismatch. \textit{See} Chapter~\ref{ch:labour}.

%----------------------------------------------------------
% T
%----------------------------------------------------------
\vspace{6pt}
\noindent\textbf{Task-based model of automation:} An economic framework that classifies tasks (rather than whole jobs) by their susceptibility to automation, distinguishing routine from non-routine and cognitive from manual tasks. \textit{See} Chapter~\ref{ch:technology}.

\vspace{6pt}
\noindent\textbf{Time inconsistency problem:} The tendency for a policy-maker who has committed to a future policy to have incentives to deviate from that commitment once the relevant time arrives. The problem underlies the case for central bank independence. \textit{See} Chapter~\ref{ch:inflation}.

\vspace{6pt}
\noindent\textbf{Trade-off:} The situation in which choosing more of one thing means getting less of another. A fundamental concept in economics reflecting scarcity. \textit{See} Chapter~\ref{ch:introduction}.

%----------------------------------------------------------
% U
%----------------------------------------------------------
\vspace{6pt}
\noindent\textbf{Unemployment rate:} The percentage of the labour force that is without work but actively seeking employment. \textit{See} Chapter~\ref{ch:labour}.

\vspace{6pt}
\noindent\textbf{Union density:} The share of paid workers who are members of a trade union. \textit{See} Chapter~\ref{ch:labour}.

\vspace{6pt}
\noindent\textbf{UNDRIP (United Nations Declaration on the Rights of Indigenous Peoples):} An international human rights instrument adopted by the UN General Assembly in 2007. Canada passed legislation in 2021 (Bill C-15) to align Canadian laws with UNDRIP. \textit{See} Chapter~\ref{ch:inequality}.

%----------------------------------------------------------
% W
%----------------------------------------------------------
\vspace{6pt}
\noindent\textbf{Wage polarization:} The pattern of employment and wage growth observed in advanced economies: expansion at the top (high-skill, high-wage) and bottom (low-skill, low-wage) of the distribution, with contraction in the middle (middle-skill, routine occupations). Attributed to skill-biased technological change and automation. \textit{See} Chapter~\ref{ch:technology}.

\vspace{6pt}
\noindent\textbf{Wealth tax:} An annual levy on the total net worth (assets minus liabilities) of individuals above a specified threshold, as distinct from taxes on income flows. \textit{See} Chapter~\ref{ch:inequality}.

\vspace{6pt}
\noindent\textbf{Winner-take-all market:} A market in which one or a few firms capture most of the market share due to strong network effects, economies of scale, or first-mover advantages. Common in digital platform markets. \textit{See} Chapter~\ref{ch:technology}.


\bibliography{references/references}

\end{document}
